% Chapter 1, Section 1 _Linear Algebra_ Jim Hefferon
%  http://joshua.smcvt.edu/linearalgebra
%  2001-Jun-09
\chapter{Linear Systems}\leavevmode  % because no preamble?
\section{Solving Linear Systems}
Systems of linear equations are common in science and mathematics.
These two examples from high school science \cite{Onan}
give a sense of how they arise.

The first example is from 
\hypertarget{ex:Statics}{Statics}.\index{Statics problem}
Suppose that we have three objects,
we know that one has a mass of 2~kg,
and we want to find the two unknown masses.
Experimentation with a meter stick produces these two balances.
\begin{center}
  \includegraphics{gr/mp/ch1.1}
  \qquad
  \includegraphics{gr/mp/ch1.2}
\end{center}
For the masses to balance we must have that
the sum of moments on the left equals the sum of moments on
the right, where the moment of an object is its mass times its distance 
from the balance point. 
That gives a system of two linear equations.
\begin{align*}  % odd formatting to match chem problem below
      40h + 15c &= 100  \\
            25c &= 50+50h
\end{align*}

The second example
is from Chemistry.\index{Chemistry problem}
We can mix, under controlled conditions, toluene $\hbox{C}_7\hbox{H}_8$ and 
nitric acid $\hbox{H}\hbox{N}\hbox{O}_3$ to produce
trinitrotoluene $\hbox{C}_7\hbox{H}_5\hbox{O}_6\hbox{N}_3$
along with the byproduct water
(conditions have to be very well controlled\Dash trinitrotoluene 
is better known as TNT).
In what proportion should we mix them?
The number of atoms of each element present before the reaction
\begin{equation*}
    x\, {\rm C}_7{\rm H}_8\ +\ y\,{\rm H}{\rm N}{\rm O}_3
    \quad\longrightarrow\quad
    z\,{\rm C}_7{\rm H}_5{\rm O}_6{\rm N}_3\ +\ w\,{\rm H}_2{\rm O}
\tag*{}\end{equation*}
must equal the number present afterward.
Applying that in turn to the elements C, H, N, and O gives
this system.
\begin{align*}  % odd formatting to state more naturally
      7x      &= 7z  \\
      8x +1y  &= 5z+2w  \\
      1y      &= 3z  \\
      3y      &= 6z+1w
\end{align*}

Both examples come down to solving a system of equations.
In each system, the equations involve only the first power of each variable.
This chapter shows how to solve any such system of equations.















\subsection{Gauss's Method}
\begin{definition}\label{def:linearcombination}
%<*df:linearcombination>
A \definend{linear combination}\index{linear combination} of
\( x_1 \), \ldots, \( x_n \) has the form
\begin{equation*}
   a_1x_1+a_2x_2+a_3x_3+\cdots+a_nx_n
\end{equation*}
where the numbers \( a_1, \ldots ,a_n\in\Re \) are the combination's
\definend{coefficients}.\index{linear equation!coefficients}
%</df:linearcombination>
%<*df:linearequations>
A \definend{linear equation}\index{linear equation} 
in the variables $x_1$, \ldots, $x_n$ 
has the form
$a_1x_1+a_2x_2+a_3x_3+\cdots+a_nx_n=d$
where
\( d\in\Re \) is the \definend{constant}.\index{linear equation!constant}

An \( n \)-tuple \( (s_1,s_2,\ldots ,s_n)\in\Re^n \) is a 
\definend{solution}\index{linear equation!solution of} %
of, or \definend{satisfies}, that equation if substituting the numbers
$s_1$, \ldots, $s_n$ for the variables
gives a true statement,
$a_1s_1+a_2s_2+\cdots+a_ns_n=d$.
A \definend{system of linear equations}\index{linear equation!system of}%
\index{system of linear equations} 
\begin{equation*}
  \begin{linsys}{4}
    a_{1,1}x_1 &+ &a_{1,2}x_2  &+  &\cdots &+ &a_{1,n}x_n &=  &d_1  \\
    a_{2,1}x_1 &+ &a_{2,2}x_2  &+  &\cdots &+ &a_{2,n}x_n &=  &d_2  \\
             &  &           &   &       &  &          &\vdotswithin{=}  \\
    a_{m,1}x_1 &+ &a_{m,2}x_2  &+  &\cdots &+ &a_{m,n}x_n &=  &d_m
  \end{linsys}
\end{equation*}
has the solution
\( (s_1,s_2,\ldots ,s_n) \) if that $n$-tuple is a solution of all
of the equations.
%</df:linearequations>
\end{definition}

\begin{example}
The combination \( 3x_1 + 2x_2 \) of $x_1$ and $x_2$ is linear.
The combination \( 3x_1^2 + 2x_2 \) is not a linear function of
$x_1$ and $x_2$, nor is \( 3x_1 + 2\sin(x_2) \).  
\end{example}

We usually take $x_1$, \ldots, $x_n$ to be different variables because
in a sum with repeats we can rearrange to 
make the elements unique, as with $2x+3y+4x=6x+3y$. 
We sometimes include terms with a zero coefficient, as in
$x-2y+0z$, and at other times omit them, as is convenient. 

\begin{example}
The ordered pair \( (-1,5) \) is a solution of this system.
\begin{equation*}
  \begin{linsys}{2}
    3x_1 &+ &2x_2 &= &7  \\
    -x_1 &+ &x_2  &= &6
  \end{linsys}
\end{equation*}
In contrast, \( (5,-1) \) is not a solution.
\end{example}

Finding the set of all solutions is 
\definend{solving}\index{system of linear equations!solving} 
the system.
We don't need guesswork or good luck; 
there is an algorithm that always works.
This algorithm is  
\definend{Gauss's Method}\index{Gauss's Method}%
\index{system of linear equations!Gauss's Method} 
(or \definend{Gaussian elimination}\index{Gaussian elimination}%
\index{system of linear equations!Gaussian elimination}
or \definend{linear elimination}\index{linear elimination}%
\index{system of linear equations!linear elimination}%
\index{system of linear equations!elimination}%
\index{elimination, Gaussian}).
% It transforms the system, step by step, into one
% with a form that we can easily solve.
% We will first illustrate how it goes and then we will see the 
% formal statement. 

\begin{example}
To solve this system
\begin{equation*}
  \begin{linsys}{3}
                     &   &      &   &3x_3  &=  &9  \\
                 x_1 &+  &5x_2  &-  &2x_3  &=  &2  \\
      \frac{1}{3}x_1 &+  &2x_2  &   &      &=  &3  
  \end{linsys}
\end{equation*}
we transform it, step by step, until it is in a form that
we can easily solve.

The first transformation
rewrites the system by interchanging the first and third row.
\begin{align*}
  \quad
  &\grstep{ \text{swap row 1 with row 3} }
  \begin{linsys}{3}
        \frac{1}{3}x_1 &+  &2x_2  &   &      &=  &3  \\
                   x_1 &+  &5x_2  &-  &2x_3  &=  &2  \\
                       &   &      &   &3x_3  &=  &9  
    \end{linsys}                                         
\end{align*}
The second transformation rescales the first row by a factor of~$3$.
\begin{align*}
\quad
  &\grstep{ \text{multiply row 1 by 3} }
  \begin{linsys}{3}
        x_1 &+  &6x_2  &   &      &=  &9  \\
        x_1 &+  &5x_2  &-  &2x_3  &=  &2  \\
            &   &      &   &3x_3  &=  &9  
   \end{linsys}                                          
\end{align*}
The third transformation is the only nontrivial one in this example.
We mentally multiply both sides of the first row by \( -1 \),
mentally add that to the second row,
and write the result in as the new second row.
\begin{align*}
  &\grstep{ \text{add \(-1\) times row 1 to row 2} }
  \begin{linsys}{3}
        x_1 &+  &6x_2  &   &      &=  &9  \\
            &   &-x_2  &-  &2x_3  &=  &-7 \\
            &   &      &   &3x_3  &=  &9  
   \end{linsys}
\end{align*}
These steps have brought the system to a
form where we can easily find the value of each variable.
The bottom equation shows that \( x_3=3 \).
Substituting $3$ for \( x_3 \) in the middle equation shows that \( x_2=1 \).
Substituting those two into the top equation
gives that \( x_1=3 \). 
Thus the system has a unique solution; 
the solution set is \set{(3,1,3)}.
\end{example}

We will use Gauss's Method throughout the book.
It is fast and easy.
We will now show that
it is also safe: 
Gauss's Method never loses solutions nor does it ever 
pick up extraneous solutions, so that
a tuple is a solution to the system before we apply the method if and only if
it is a solution after.

\begin{theorem}[Gauss's Method]
\index{linear equation!solution of!Gauss's Method} \label{th:GaussMethod}
%<*th:GaussMethod>
If a linear system is changed to another by one of these operations
\begin{enumerate}
  \setlength{\itemsep}{0ex}
  \item
    an equation is swapped with another
  \item
    an equation has both sides multiplied by a nonzero constant
  \item
    an equation is replaced by the sum of itself and a multiple of another
\end{enumerate}
then the two systems have the same set of solutions.
%</th:GaussMethod>
\end{theorem}

Each of those three operations has a restriction.
Multiplying a row by~\( 0 \) is not allowed because 
it will lose all the information in that row and so 
potentially change the solution set.
Similarly, we disallow adding a multiple of a row to itself because
adding~\( -1 \) times the row to itself is the same as multiplying that row
by~\( 0 \).
And we disallow swapping a row with itself,
to simplify some results in the fourth chapter.
Besides, it's pointless.

\begin{proof}
We will cover the equation swap operation here. 
The other two
cases are similar and are \nearbyexercise{ex:ProveGaussMethod}.

%<*pf:GaussMethod0>
Consider a linear system.
\begin{equation*}
  \begin{linsys}{4}
    a_{1,1}x_1  &+  &a_{1,2}x_2 &+  &\cdots  &+ &a_{1,n}x_n  &=  &d_1\hfill\hbox{}  \\
                &   &           &   &        &   &     &\vdotswithin{=}  \\
    a_{i,1}x_1  &+  &a_{i,2}x_2 &+  &\cdots  &+ &a_{i,n}x_n  &=  &d_i\hfill\hbox{}  \\
                &   &           &   &        &   &     &\vdotswithin{=}  \\
    a_{j,1}x_1  &+  &a_{j,2}x_2 &+  &\cdots  &+ &a_{j,n}x_n  &=  &d_j\hfill\hbox{}  \\
                &   &           &   &        &   &     &\vdotswithin{=} \\
    a_{m,1}x_1  &+  &a_{m,2}x_2 &+  &\cdots  &+ &a_{m,n}x_n  &=  &d_m  \\
  \end{linsys}
\end{equation*} 
The tuple \( (s_1,\ldots\,,s_n) \)
satisfies it  
if and only if substituting the values for the
variables, the $s$'s for the $x$'s, gives a conjunction of true statements:
$a_{1,1}s_1+a_{1,2}s_2+\cdots+a_{1,n}s_n=d_1$
and \ldots\ 
$a_{i,1}s_1+a_{i,2}s_2+\cdots+a_{i,n}s_n=d_i$
and \ldots\  $a_{j,1}s_1+a_{j,2}s_2+\cdots+a_{j,n}s_n=d_j$
and \ldots\  $a_{m,1}s_1+a_{m,2}s_2+\cdots+a_{m,n}s_n=d_m$.
%</pf:GaussMethod0>

%<*pf:GaussMethod1>
In a list of statements joined with `and' we can  
rearrange the order of the statements. 
Thus the tuple satisfies that
if and only if
$a_{1,1}s_1+a_{1,2}s_2+\cdots+a_{1,n}s_n=d_1$
and \ldots\  $a_{j,1}s_1+a_{j,2}s_2+\cdots+a_{j,n}s_n=d_j$
and \ldots\  $a_{i,1}s_1+a_{i,2}s_2+\cdots+a_{i,n}s_n=d_i$
and \ldots\  $a_{m,1}s_1+a_{m,2}s_2+\cdots+a_{m,n}s_n=d_m$.
This exactly says that \( (s_1,\ldots\,,s_n) \) 
solves the system after the row swap.
%</pf:GaussMethod1>
\end{proof}

\begin{definition} \label{df:GaussMethod}
%<*df:GaussMethod>
The three operations from 
\nearbytheorem{th:GaussMethod}
are the
\definend{elementary reduction 
operations},\index{Gauss's Method!elementary operations}%
\index{elementary reduction operations}
or \definend{row operations}\index{elementary row operations},
or \definend{Gaussian operations}.
They are
\definend{swapping}\index{elementary reduction operations! swapping}%
\index{swapping rows},
\definend{multiplying by a scalar} (or
\definend{rescaling}\index{elementary reduction operations! rescaling}%
\index{rescaling rows}), and
\definend{row combination}\index{elementary reduction operations! row combination}%
\index{combining rows}\index{adding rows}.
%</df:GaussMethod>
\end{definition}

When writing out the calculations, we will 
abbreviate `row \(i\)' by `\( \rho_i \)'
(this is the Greek letter rho, pronounced ``row'').
For instance, we will denote a row combination operation by 
\( k\rho_i+\rho_j \), 
with the row that changes written second.
To save writing we will 
often combine addition steps when they use the same $\rho_i$, as in the
next example.

\begin{example}
Gauss's Method methodically applies the row operations to solve a system.
Here is a typical case.
\begin{equation*}
  \begin{linsys}{3}
    x  &+  &y  &   &   &=  &0  \\
   2x  &-  &y  &+  &3z &=  &3  \\
    x  &-  &2y &-  &z  &=  &3  
  \end{linsys}
\end{equation*}
We begin by using the first row to 
eliminate the $2x$ in the second row and the $x$ in the third.
To get rid of the $2x$ we mentally multiply the entire first row by $-2$, 
add that to the
second row, and write the result in as the new second row.
To eliminate the~$x$ in the third row we multiply the first row by
$-1$, add that to the third row, and write the result in as the
new third row.
\begin{align*}
  &\grstep[-\rho_1 +\rho_3]{-2\rho_1 +\rho_2}
  \begin{linsys}{3}
     x  &+  &y  &   &   &=  &0  \\
        &   &-3y&+  &3z &=  &3  \\
        &   &-3y&-  &z  &=  &3  
  \end{linsys}   
\end{align*}
% In this version of the system, the last two equations involve only two unknowns.
We finish by transforming the second system into a third, where the
bottom equation involves only one unknown. 
We do that by using 
the second row to eliminate the $y$~term from the third row.
\begin{equation*}
  \grstep{-\rho_2 +\rho_3}
  \begin{linsys}{3}
     x  &+  &y  &   &   &=  &0  \\
        &   &-3y&+  &3z &=  &3  \\
        &   &   &   &-4z&=  &0
   \end{linsys}
\end{equation*}
Now finding the system's solution is easy.
The third row gives \( z=0 \).
Substitute that back\index{Gauss's Method!back-substitution}%
\index{back-substitution}
into the second row to get \( y=-1 \).
Then substitute back into the first row to get \( x=1 \).
\end{example}

\begin{example}
For the Physics problem\index{Statics problem} from the start of this
chapter, Gauss's Method gives this.
\begin{equation*}
   \begin{linsys}{2}
     40h  &+  &15c  &=  &100      \\
     -50h &+  &25c  &=  &50         
   \end{linsys}
   \grstep{5/4\rho_1 +\rho_2}
   \begin{linsys}{2}
      40h  &+  &15c       &=  &100      \\
           &   &(175/4)c  &=  &175 
    \end{linsys}
\end{equation*}
So \( c=4 \), and back-substitution gives that \( h=1 \).
(We will solve the Chemistry problem later.)
\end{example}

\begin{example}
The reduction
\begin{align*}
   \begin{linsys}{3}
        x  &+  &y  &+  &z  &=  &9  \\
       2x  &+  &4y &-  &3z &=  &1  \\
       3x  &+  &6y &-  &5z &=  &0  
   \end{linsys}
   &\grstep[-3\rho_1 +\rho_3]{-2\rho_1 +\rho_2}
   \begin{linsys}{3}
      x  &+  &y  &+  &z  &=  &9  \\
         &   &2y &-  &5z &=  &-17\\
         &   &3y &-  &8z&=  &-27
    \end{linsys}                                    \\
   &\grstep{-(3/2)\rho_2+\rho_3}
   \begin{linsys}{3}
      x  &+  &y  &+  &z            &=  &9  \\
         &   &2y &-  &5z           &=  &-17\\
         &   &   &   &-(1/2)z      &=  &-(3/2) 
    \end{linsys}
\end{align*}
shows that \( z=3 \), \( y=-1 \), and \( x=7 \).
\end{example}

As illustrated above, the point of Gauss's Method 
is to use the elementary reduction
operations to set up back-substitution.

\begin{definition} \label{df:EchelonForm}
%<*df:EchelonForm>
In each row of a system, 
the first variable with a nonzero coefficient is the row's
\definend{leading variable}\index{echelon form!leading variable}%
\index{leading!variable}. % 
A system is in \definend{echelon form}\index{echelon form}
if each leading variable
is to the right of the leading variable in the row above it,
except for the leading variable in the first row,
and any rows with all-zero coefficients are at the bottom.
%</df:EchelonForm>
\end{definition}

\begin{example} 
The prior three examples only used the operation of row combination.
This linear system requires the swap operation
to get it into echelon form because 
after the first combination
\begin{align*}
   \begin{linsys}{4}
                x  &-  &y  &   &   &   &   &=  &0  \\
               2x  &-  &2y &+  &z  &+  &2w &=  &4  \\
                   &   &y  &   &   &+  &w  &=  &0  \\
                   &   &   &   &2z &+  &w  &=  &5  
   \end{linsys}
   &\grstep{-2\rho_1 +\rho_2}
   \begin{linsys}{4}
      x  &-  &y  &\spaceforemptycolumn   &   &   &   &=  &0  \\
         &   &   &   &z  &+  &2w &=  &4  \\
         &   &y  &   &   &+  &w  &=  &0  \\
         &   &   &   &2z &+  &w  &=  &5  
   \end{linsys}    
\end{align*}
the second equation has no leading $y$.
We exchange it for a lower-down row that has a leading $y$.
\begin{align*}
   &\grstep{\rho_2 \leftrightarrow\rho_3}
   \begin{linsys}{4}
      x  &-  &y  &\spaceforemptycolumn   &   &   &   &=  &0  \\
         &   &y  &   &   &+  &w  &=  &0  \\
         &   &   &   &z  &+  &2w &=  &4  \\
         &   &   &   &2z &+  &w  &=  &5  
    \end{linsys}    
\end{align*}
(Had there been more than one suitable row below the second
then we could have used any one.)
With that, Gauss's Method proceeds as before.
\begin{align*}
   &\grstep{-2\rho_3 +\rho_4}
   \begin{linsys}{4}
      x  &-  &y  &\spaceforemptycolumn   &   &   &   &=  &0  \\
         &   &y  &   &   &+  &w  &=  &0  \\
         &   &   &   &z  &+  &2w &=  &4  \\
         &   &   &   &   &   &-3w&=  &-3 
    \end{linsys}
\end{align*}
Back-substitution gives \( w=1 \), \( z=2 \) , \( y=-1 \), and \( x=-1 \).
\end{example}

Strictly speaking, to solve linear systems we don't need 
the row rescaling operation.  
We have introduced it here because it is natural and because we will use it 
later in this chapter as part of a variation of Gauss's Method, 
the Gauss-Jordan Method.

All of the systems so far have the same number of equations as unknowns.
All of them have a solution and for all of them there is only one solution.
We finish this subsection by seeing
other things that can happen.

\begin{example} \label{ex:MoreEqsThanUnks}
This system 
has more equations than variables.
\begin{equation*}
    \begin{linsys}{2}
      x  &+  &3y  &=  &1  \\
     2x  &+  &y   &=  &-3 \\
     2x  &+  &2y  &=  &-2 
    \end{linsys}  
\end{equation*}
Gauss's Method helps here also, since this
\begin{align*}
    &\grstep[-2\rho_1 +\rho_3]{-2\rho_1 +\rho_2}
    \begin{linsys}{2}
       x  &+  &3y  &=  &1  \\
          &   &-5y &=  &-5 \\
          &   &-4y &=  &-4 
     \end{linsys}
\end{align*}
shows that one of the equations is redundant.
Echelon form
\begin{equation*}
    \grstep{-(4/5)\rho_2 +\rho_3}
    \begin{linsys}{2}
       x  &+  &3y  &=  &1  \\
          &   &-5y &=  &-5 \\
          &   &0   &=  &0  
     \end{linsys}
\end{equation*}
gives that \( y=1 \) and \( x=-2 \).
The `\( 0=0 \)' reflects the redundancy.
\end{example}

Gauss's Method is also useful on systems with more variables than equations.
The next subsection has many examples.

Another way that linear systems can differ from the examples shown above 
is that some linear systems do not have a unique solution.
This can happen in two ways.
The first is that a system can fail to have any solution at all.

\begin{example} \label{ex:MoreEqsThanUnksInconsis}
Contrast the system in the last example with this one.
\begin{equation*}
    \begin{linsys}{2}
      x  &+  &3y  &=  &1  \\
     2x  &+  &y   &=  &-3 \\
     2x  &+  &2y  &=  &0  
    \end{linsys}
    \grstep[-2\rho_1 +\rho_3]{-2\rho_1 +\rho_2}
    \begin{linsys}{2}
       x  &+  &3y  &=  &1  \\
          &   &-5y &=  &-5 \\
          &   &-4y &=  &-2
     \end{linsys}
\end{equation*}
Here the system is inconsistent:~no pair of numbers $(s_1,s_2)$ satisfies 
all three equations simultaneously.
Echelon form makes the inconsistency obvious.
\begin{equation*}
  \grstep{-(4/5)\rho_2 +\rho_3}
  \begin{linsys}{2}
     x  &+  &3y  &=  &1  \\
        &   &-5y &=  &-5 \\
        &   &0   &=  &2 
   \end{linsys}
\end{equation*}
The solution set is empty.
\end{example}

\begin{example}
The prior system has more equations than unknowns but
that is not what causes the inconsistency\Dash 
\nearbyexample{ex:MoreEqsThanUnks}
has more equations than unknowns and yet is consistent.
Nor is having more equations than unknowns necessary for
inconsistency, as we see with this inconsistent system that has the 
same number of equations as unknowns.
\begin{equation*}
  \begin{linsys}{2}
    x  &+  &2y  &=  &8  \\
   2x  &+  &4y  &=  &8  
  \end{linsys}
  \grstep{-2\rho_1 + \rho_2}
  \begin{linsys}{2}
     x  &+  &2y  &=  &8  \\
        &   &0   &=  &-8
   \end{linsys}
\end{equation*}
Instead, 
inconsistency has to do with the interaction of the left and right sides.
In the first system above the left side's second equation is twice the
first but the right side's second constant is not twice the first.
Later we will have more to say about dependencies between a system's 
parts.
\end{example}

The other way that 
a linear system can fail to have a unique solution, besides having no solutions,
is to have many solutions.

\begin{example}
In this system
\begin{equation*}
  \begin{linsys}{2}
    x  &+  &y   &=  &4  \\
   2x  &+  &2y  &=  &8  
  \end{linsys}
\end{equation*}
any pair of numbers satisfying the first equation also
satisfies the second.
The solution set 
% \( \{ (x,y)\suchthat x+y=4 \} \) 
\( \set{ (x,y)\suchthat x+y=4} \) 
is infinite; some example
member pairs are~$(0,4)$, $(-1,5)$, and $(2.5,1.5)$.

The result of applying Gauss's Method here contrasts with the prior example
because we do not get a contradictory equation.
\begin{equation*}
  \grstep{-2\rho_1 + \rho_2}
  \begin{linsys}{2}
    x  &+  &y   &=  &4  \\
       &   &0   &=  &0  
   \end{linsys}
\end{equation*}
\end{example}

Don't be fooled by that example:~a $0=0$ equation  
is not the signal that a system has many solutions.

\begin{example}  \label{ex:NoZerosInfManySols}
The absence of a \( 0=0 \) equation does not keep a system from having
many different solutions.
This system is in echelon form,
has no $0=0$, but has infinitely many solutions,
including $(0,1,-1)$, 
$(0,1/2,-1/2)$, $(0,0,0)$, and $(0,-\pi,\pi)$.
In general, 
a triple is a solution if its 
first component is $0$ and its second component is the 
negative of the third.
\begin{equation*}
  \begin{linsys}{3}
     x  &+  &y  &+  &z  &=  &0  \\
        &   &y  &+  &z  &=  &0  
   \end{linsys}
\end{equation*}

Nor does the presence of \( 0=0 \) mean that the system must have 
many solutions.
\nearbyexample{ex:MoreEqsThanUnks} shows that.
So does this system, which does not have 
any solutions at all despite that 
in echelon form it has a $0=0$ row.
\begin{align*}
  \begin{linsys}{3}
    2x  &   &   &-   &2z  &=  &6  \\
        &   &y  &+   &z   &=  &1  \\
    2x  &+  &y  &-   &z   &=  &7  \\
        &   &3y &+   &3z  &=  &0  
  \end{linsys}
  &\grstep{-\rho_1 +\rho_3}
  \begin{linsys}{3}
     2x  &\spaceforemptycolumn   &   &-   &2z  &=  &6  \\
         &   &y  &+   &z   &=  &1  \\
         &   &y  &+   &z   &=  &1  \\
         &   &3y &+   &3z  &=  &0  
  \end{linsys}                                    \\
  &\grstep[-3\rho_2 +\rho_4]{-\rho_2 +\rho_3}
  \begin{linsys}{3}
     2x  &\spaceforemptycolumn    &   &-   &2z  &=  &6  \\
         &   &y  &+   &z   &=  &1  \\
         &   &   &    &0   &=  &0  \\
         &   &   &    &0   &=  &-3 
   \end{linsys}
\end{align*}
\end{example}

In summary,
Gauss's Method uses row operations to 
set up back substitution, by a sequence of whole system transformations
(for contrast see \nearbyexercise{exer:SolvingSystemWithoutGM}).
If any of the transformation steps results in a contradictory equation 
then we can stop and 
conclude that the system has no solutions.
If we reach echelon form without a contradictory equation,
and each variable is a leading variable in a
row, then the system has a unique solution and we find it by
back substitution.
Finally, if we reach echelon form without a contradictory equation,
and there is not a unique solution\Dash
if at least one variable is not a leading variable\Dash
then the system has many solutions.

The next subsection shows how to
describe solution sets.

\exerciseNote{In the exercises here and in the rest of the book,
you must justify your answers.
For instance, if a question asks whether a system has a solution then you
must justify a yes response by producing the solution and must justify 
a no response by showing that no solution exists.}
\begin{exercises}
  \recommended \item 
    Use Gauss's Method to find the unique solution for each system.
    \begin{exparts*}
      \partsitem 
        $\begin{linsys}[t]{2}
          2x  &+  &3y  &=  &13  \\
          x   &-  &y   &=  &-1
        \end{linsys}$
      \partsitem 
        $\begin{linsys}[t]{3}
          x   &  &  &-  &z  &=  &0  \\
          3x  &+ &y &   &   &=  &1  \\
          -x  &+ &y &+  &z  &=  &4
        \end{linsys}$
    \end{exparts*}
    \begin{answer}
      \begin{exparts}
        \partsitem Gauss's Method
          \begin{equation*}
            \grstep{-(1/2)\rho_1+\rho_2}
            \begin{linsys}{2}
               2x  &+  &3y      &=  &13  \\
                   &-  &(5/2)y  &=  &-15/2              
            \end{linsys}
          \end{equation*}
          gives that the solution is $y=3$ and $x=2$.
        \partsitem Gauss's Method here
          \begin{equation*}
            \grstep[\rho_1+\rho_3]{-3\rho_1+\rho_2}
            \begin{linsys}{3}
              x   &  &  &-  &z  &=  &0  \\
                  &  &y &+  &3z &=  &1  \\
                  &  &y &   &   &=  &4
            \end{linsys}
            \grstep{-\rho_2+\rho_3}
            \begin{linsys}{3}
              x   &  &  &-  &z    &=  &0  \\
                  &  &y &+  &3z   &=  &1  \\
                  &  &  &   &-3z  &=  &3
            \end{linsys}
          \end{equation*}
          gives $x=-1$, $y=4$, and $z=-1$.
      \end{exparts}
      % sage: attach("gauss_method_latex.sage")
      % sage: P = matrix(QQ, [[2, 3], [1, -1]])
      % sage: v = vector(QQ, [13, -1])
      % sage: P_prime = P.augment(v, subdivide=True)
      % sage: gauss_jordan(P_prime)
      % [ 2  3|13]
      % [ 1 -1|-1]
      % take -1/2 times row 1 plus row 2
      % [    2     3|   13]
      % [    0  -5/2|-15/2]
      % take 1/2 times row 1
      % take -2/5 times row 2
      % [   1  3/2|13/2]
      % [   0    1|   3]
      % take -3/2 times row 2 plus row 1
      % [1 0|2]
      % [0 1|3]
      % sage: P = matrix(QQ, [[1,0,-1], [3,1,0], [-1,1,1]])
      % sage: v = vector(QQ, [0, 1,4])
      % sage: P_prime = P.augment(v, subdivide=True)
      % sage: gauss_jordan(P_prime)
      % [ 1  0 -1| 0]
      % [ 3  1  0| 1]
      % [-1  1  1| 4]
      % take -3 times row 1 plus row 2
      % take 1 times row 1 plus row 3
      % [ 1  0 -1| 0]
      % [ 0  1  3| 1]
      % [ 0  1  0| 4]
      % take -1 times row 2 plus row 3
      % [ 1  0 -1| 0]
      % [ 0  1  3| 1]
      % [ 0  0 -3| 3]
      % take -1/3 times row 3
      % [ 1  0 -1| 0]
      % [ 0  1  3| 1]
      % [ 0  0  1|-1]
      % take 1 times row 3 plus row 1
      % take -3 times row 3 plus row 2
      % [ 1  0  0|-1]
      % [ 0  1  0| 4]
      % [ 0  0  1|-1]
    \end{answer}
  \item
    Each system is in echelon form.
    For each, say whether the system has a unique solution, 
    no solution, or infinitely many solutions.
    \begin{exparts*}
      \partsitem 
        \(\begin{linsys}[t]{2}
           -3x &+ &2y   &= &0  \\          
               &  &-2y  &= &0            
        \end{linsys}\)
      \partsitem 
        \(\begin{linsys}[t]{3}
             x &+ &y &  &   &= &4  \\          
               &  &y &- &z  &= &0  \\           
        \end{linsys}\)
      \partsitem 
        \(\begin{linsys}[t]{3}
             x &+ &y &  &   &= &4  \\          
               &  &y &  &-z  &= &0  \\ 
               &  &  &  &0  &= &0  
        \end{linsys}\)
      \partsitem 
        \(\begin{linsys}[t]{2}
             x &+ &y    &= &4  \\          
               &  &0    &= &4            
        \end{linsys}\)
      \partsitem 
        \(\begin{linsys}[t]{3}
            3x &+ &6y &+  &z  &= &-0.5  \\          
               &  &   &   &-z &= &2.5  \\           
        \end{linsys}\)
      \partsitem 
        \(\begin{linsys}[t]{2}
             x &- &3y    &= &2  \\          
               &  &0    &= &0            
        \end{linsys}\)
      \partsitem 
        \(\begin{linsys}[t]{2}
             2x &+ &2y    &= &4  \\
                &  &y     &= &1 \\          
                 &  &0    &= &4            
        \end{linsys}\)
      \partsitem 
        \(\begin{linsys}[t]{2}
              2x &+ &y   &= &0  
        \end{linsys}\)
      \partsitem 
        \(\begin{linsys}[t]{2}
              x &- &y    &= &-1  \\
                &  &0     &= &0 \\          
                 &  &0    &= &4            
        \end{linsys}\)
      \partsitem 
        \(\begin{linsys}[t]{3}
              x &+ &y  &- &3z  &= &-1  \\
                &  &y  &- &z  &= &2  \\
                &  &  &  &z    &= &0  \\
                &  &  &  &0   &= &0  
        \end{linsys}\)
    \end{exparts*}
    \begin{answer}
      If a system has a contradictory equation then it has no solution.
      Otherwise, if there are any variables that are not leading a row
      then it has infinitely many solution.
      In the final case, where there is no contradictory equation and
      every variable leads some row, it has a unique solution.
      \begin{exparts*}
        \partsitem \mbox{Unique solution}
        \partsitem \mbox{Infinitely many solutions}
        \partsitem \mbox{Infinitely many solutions}
        \partsitem \mbox{No solution}
        \partsitem \mbox{Infinitely many solutions}
        \partsitem \mbox{Infinitely many solutions}
        \partsitem \mbox{No solution}
        \partsitem \mbox{Infinitely many solutions}
        \partsitem \mbox{No solution}
        \partsitem \mbox{Unique solution}
      \end{exparts*}
    \end{answer}
  \recommended \item  
    Use Gauss's Method to solve each system
    or conclude `many solutions' or `no solutions'.
    \begin{exparts*}
      \partsitem \(
               \begin{linsys}[t]{2}
                  2x  &+  &2y  &=  &5  \\
                   x  &-  &4y  &=  &0  
               \end{linsys}
             \)
      \partsitem \(
               \begin{linsys}[t]{2}
                  -x  &+  &y   &=  &1  \\
                   x  &+  &y   &=  &2  
               \end{linsys}
             \) 
      \partsitem  \(
               \begin{linsys}[t]{3}
                   x  &-  &3y  &+  &z  &=  &1  \\
                   x  &+  &y   &+  &2z &=  &14 
                \end{linsys}
             \) 
      \partsitem  \(
               \begin{linsys}[t]{2}
                  -x  &-  &y   &=  &1  \\
                 -3x  &-  &3y  &=  &2  
               \end{linsys}
             \) 
      \partsitem  \(
               \begin{linsys}[t]{3}
                      &   &4y  &+  &z  &=  &20 \\
                  2x  &-  &2y  &+  &z  &=  &0  \\
                   x  &   &    &+  &z  &=  &5  \\
                   x  &+  &y   &-  &z  &=  &10 
                \end{linsys}
             \)
      \partsitem \( \begin{linsys}[t]{4}
                 2x  &   &   &+  &z  &+  &w  &=  &5  \\
                     &   &y  &   &   &-  &w  &=  &-1 \\
                 3x  &   &   &-  &z  &-  &w  &=  &0  \\
                 4x  &+  &y  &+  &2z &+  &w  &=  &9  
               \end{linsys}
            \)
    \end{exparts*}
    \begin{answer} 
      \begin{exparts}
       \partsitem Gaussian reduction
        \begin{equation*}
          \grstep{-(1/2)\rho_1+\rho_2}
          \begin{linsys}{2}
             2x  &+  &2y  &=  &5  \\
                 &   &-5y &=  &-5/2  
          \end{linsys}
        \end{equation*}
        shows that \( y=1/2 \) and \( x=2 \) is the unique solution.
      \partsitem Gauss's Method
        \begin{equation*}
          \grstep{\rho_1+\rho_2}
          \begin{linsys}{2}
             -x  &+  &y   &=  &1  \\
                 &   &2y  &=  &3  
           \end{linsys}
        \end{equation*}
        gives \( y=3/2 \) and \( x=1/2 \) as the only solution.
      \partsitem Row reduction
        \begin{equation*}
            \grstep{-\rho_1+\rho_2}
            \begin{linsys}{3}
                x  &-  &3y  &+  &z  &=  &1  \\
                   &   &4y  &+  &z  &=  &13 
             \end{linsys}
        \end{equation*}
        shows, because the variable $z$ is not a leading variable in any
        row, that there are many solutions.
      \partsitem Row reduction
        \begin{equation*}
          \grstep{-3\rho_1+\rho_2}
          \begin{linsys}{2}
             -x  &-  &y   &=  &1  \\
                 &   &0   &=  &-1 
           \end{linsys}
        \end{equation*}
        shows that there is no solution.
      \partsitem Gauss's Method
        \begin{align*}
            \grstep{\rho_1\leftrightarrow\rho_4}
            \begin{linsys}{3}
                x  &+  &y   &-  &z  &=  &10 \\
               2x  &-  &2y  &+  &z  &=  &0  \\
                x  &   &    &+  &z  &=  &5  \\
                   &   &4y  &+  &z  &=  &20 
             \end{linsys}
            &\grstep[-\rho_1+\rho_3]{-2\rho_1+\rho_2}
            \begin{linsys}{3}
                x  &+  &y   &-  &z  &=  &10 \\
                   &   &-4y &+  &3z &=  &-20\\
                   &   &-y  &+  &2z &=  &-5 \\
                   &   &4y  &+  &z  &=  &20 
             \end{linsys}                                      \\
            &\grstep[\rho_2+\rho_4]{-(1/4)\rho_2+\rho_3}
            \begin{linsys}{3}
                x  &+  &y   &-  &z      &=  &10 \\
                   &   &-4y &+  &3z     &=  &-20\\
                   &   &    &   &(5/4)z &=  &0  \\
                   &   &    &   &4z     &=  &0  
             \end{linsys}
        \end{align*}
        gives the unique solution \( (x,y,z)=(5,5,0) \).
      \partsitem Here Gauss's Method gives
         \begin{align*}
            &\grstep[-2\rho_1+\rho_4]{-(3/2)\rho_1+\rho_3}
            \begin{linsys}{4}
               2x  &\spaceforemptycolumn   &   &+  &z       &+  &w       &=  &5  \\
                   &   &y  &   &        &-  &w       &=  &-1 \\
                   &   &   &-  &(5/2)z  &-  &(5/2)w  &=  &-15/2  \\
                   &   &y  &   &        &-  &w       &=  &-1  
             \end{linsys}                                             \\ 
            &\grstep{-\rho_2+\rho_4}
            \begin{linsys}{4}
               2x  &\spaceforemptycolumn   &   &+  &z       &+  &w       &=  &5  \\
                   &   &y  &   &        &-  &w       &=  &-1 \\
                   &   &   &-  &(5/2)z  &-  &(5/2)w  &=  &-15/2  \\
                   &   &   &   &        &   &0       &=  &0 
             \end{linsys}
         \end{align*}
         which shows that there are many solutions.
      \end{exparts} 
    \end{answer}
  \item  
    Solve each system
    or conclude `many solutions' or `no solutions'.
    Use Gauss's Method.
    \begin{exparts*}
      \partsitem \(
               \begin{linsys}[t]{3}
                   x  &+  &y   &+  &z   &=  &5  \\
                   x  &-  &y   &   &    &=  &0  \\
                      &   &y   &+  &2z  &=  &7  
               \end{linsys}
             \)
      \partsitem \(
               \begin{linsys}[t]{3}
                   3x  &  &    &+  &z   &=  &7  \\
                    x  &-  &y   &+   &3z    &=  &4  \\
                    x   &+  &2y   &-  &5z  &=  &-1  
               \end{linsys}
             \)
      \partsitem \(
               \begin{linsys}[t]{3}
                   x   &+  &3y   &+  &z   &=  &0  \\
                   -x  &-  &y   &   &    &=  &2  \\
                   -x   &+  &y   &+  &2z  &=  &8  
               \end{linsys}
             \)
    \end{exparts*}
    \begin{answer} 
      \begin{exparts}
       \partsitem 
         % sage: P = matrix(QQ, [[1, 1, 1], [1, -1, 0], [0, 1, 2]])
         % sage: v = vector(QQ, [5, 0, 7])
         % sage: P_prime = P.augment(v, subdivide=True)
         % sage: load("gauss_method_latex.sage")
         % sage: gauss_method(P_prime)
         % [ 1  1  1| 5]
         % [ 1 -1  0| 0]
         % [ 0  1  2| 7]
         % take -1 times row 1 plus row 2
         % [ 1  1  1| 5]
         % [ 0 -2 -1|-5]
         % [ 0  1  2| 7]
         % take 1/2 times row 2 plus row 3
         % [  1   1   1|  5]
         % [  0  -2  -1| -5]
         % [  0   0 3/2|9/2]
         % sage: x,y,z=var('x,y,z')
         % sage: solve([3/2*z==9/2, -2*y-z==-5, x+y+z==5], x, y, z)
         % [[x == 1, y == 1, z == 3]]
         Gauss's Method
         \begin{align*}
          \begin{linsys}{3}
             x  &+  &y   &+  &z   &=  &5  \\
             x  &-  &y   &   &    &=  &0  \\
                &   &y   &+  &2z  &=  &7  
          \end{linsys}
          &\grstep{-\rho_1+\rho_2} 
          \begin{linsys}{3}
             x  &+  &y   &+  &z   &=  &5  \\
                &   &-2y  &-   &z    &=  &-5  \\
                &   &y   &+  &2z  &=  &7  
          \end{linsys}                                     \\
          &\grstep{(1/2)\rho_2+\rho_3} 
          \begin{linsys}{3}
             x  &+  &y   &+  &z   &=  &5  \\
                &   &-2y  &-   &1    &=  &-5  \\
                &   &    &   &(3/2)z  &=  &9/2  
          \end{linsys}
         \end{align*}
         followed by back-substitution gives $x=1$, $y=1$, and $z=3$.
       \partsitem 
         % sage: M = matrix(QQ, [[3,0,1], [1,-1,3], [1,2,-5]])
         % sage: v = vector(RDF, [7,4,-1])
         % sage: v = vector(QQ, [7,4,-1])
         % sage: M_prime = M.augment(v)
         % sage: gauss_method(M_prime)
         % [ 3  0  1  7]
         % [ 1 -1  3  4]
         % [ 1  2 -5 -1]
         %  take -1/3 times row 1 plus row 2
         %  take -1/3 times row 1 plus row 3
         % [    3     0     1     7]
         % [    0    -1   8/3   5/3]
         % [    0     2 -16/3 -10/3]
         %  take 2 times row 2 plus row 3
         % [  3   0   1   7]
         % [  0  -1 8/3 5/3]
         % [  0   0   0   0]
         Here Gauss's Method
         \begin{align*}
           \begin{linsys}{3}
               3x  &  &    &+  &z   &=  &7  \\
                x  &-  &y   &+   &3z    &=  &4  \\
                x   &+  &2y   &-  &5z  &=  &-1  
           \end{linsys}
           &\grstep[-(1/3)\rho_1+\rho_3]{-(1/3)\rho_1+\rho_2} 
           \begin{linsys}{3}
               3x  &  &    &+  &z        &=  &7  \\
                   &  &-y   &+ &(8/3)z   &=  &5/3  \\
                   &  &2y   &- &(16/3)z  &=  &-(10/3)  
           \end{linsys}                                      \\
           &\grstep{2\rho_2+\rho_3} 
           \begin{linsys}{3}
               3x  &  &    &+  &z        &=  &7  \\
                   &  &-y   &+ &(8/3)z   &=  &5/3  \\
                   &  &     &  &0        &=  &0  
           \end{linsys}
       \end{align*}
       finds that the variable~$z$ does not lead a row. 
       There are many solutions.
       \partsitem
         % sage: M = matrix(QQ, [[1,3,1], [-1,-1,0], [-1,1,2]])
         % sage: v = vector(QQ, [0,2,8])
         % sage: M_prime = M.augment(v)
         % sage: gauss_method(M_prime)
         % [ 1  3  1  0]
         % [-1 -1  0  2]
         % [-1  1  2  8]
         %  take 1 times row 1 plus row 2
         %  take 1 times row 1 plus row 3
         % [1 3 1 0]
         % [0 2 1 2]
         % [0 4 3 8]
         %  take -2 times row 2 plus row 3
         % [1 3 1 0]
         % [0 2 1 2]
         % [0 0 1 4]
         The steps
         \begin{equation*}
           \begin{linsys}{3}
              x   &+  &3y   &+  &z   &=  &0  \\
              -x  &-  &y   &   &    &=  &2  \\
              -x   &+  &y   &+  &2z  &=  &8  
           \end{linsys}
           \grstep[\rho_1+\rho_3]{\rho_1+\rho_2}
           \begin{linsys}{3}
              x   &+  &3y   &+  &z   &=  &0  \\
                  &   &2y   &+  &z   &=  &2  \\
                  &   &4y   &+  &3z  &=  &8  
           \end{linsys}
           \grstep{-2\rho_2+\rho_3}
           \begin{linsys}{3}
              x   &+  &3y   &+  &z   &=  &0  \\
                  &   &2y   &+  &z   &=  &2  \\
                  &   &     &   &z   &=  &4  
           \end{linsys}
         \end{equation*}
         give $(x,y,z)=(-1,-1,4)$.
      \end{exparts} 
    \end{answer}
  \recommended \item \label{exer:SolvingSystemWithoutGM}
    We can attack linear systems by methods other 
    than Gauss's.
    One often taught in high school is to solve one of the 
    equations for a variable, then substitute the resulting expression into
    other equations.
    Then we repeat that step until there is an equation with only one
    variable.
    From that we get the first number in the solution and then we get the
    rest with  
    back-substitution.
    This method takes longer than Gauss's Method, since it involves
    more arithmetic operations, and is also more
    likely to lead to errors.
    To illustrate how it can lead to wrong conclusions, we will use the system 
    \begin{equation*}
      \begin{linsys}{2}
            x  &+  &3y  &=  &1  \\
            2x  &+  &y   &=  &-3 \\
            2x  &+  &2y  &=  &0  
      \end{linsys}
    \end{equation*}
    from \nearbyexample{ex:MoreEqsThanUnksInconsis}.
    \begin{exparts}
      \partsitem Solve the first equation for $x$ and 
        substitute that expression into the second equation.
        Find the resulting $y$.
      \partsitem Again solve the first equation for $x$, 
        but this time substitute that expression into the third equation.
        Find this $y$.
    \end{exparts}
    Note that this approach does not process the entire system. 
    What extra step must a user of this method take to avoid 
    erroneously concluding a system has a solution?
    \begin{answer}
      \begin{exparts}
        \partsitem From $x=1-3y$ we get that $2(1-3y)+y=-3$, giving $y=1$.
        \partsitem From $x=1-3y$ we get that $2(1-3y)+2y=0$, leading to 
           the conclusion that $y=1/2$.
      \end{exparts}
      Users of this method must check any potential solutions by
      substituting back into all the equations.
    \end{answer}
  \recommended \item 
    For which values of \( k \) are
    there no solutions, many solutions, or a unique solution?
    \begin{equation*}
       \begin{linsys}{2}
          x  &-  &y  &=  &1  \\
         3x  &-  &3y &=  &k  
       \end{linsys}
    \end{equation*}
    \begin{answer}
      Do the reduction
      \begin{equation*}
        \grstep{-3\rho_1+\rho_2}
        \begin{linsys}{2}
          x  &-  &y  &=  &1\hfill  \\
             &   &0  &=  &-3+k\hfill  
        \end{linsys}
      \end{equation*}
      to conclude this system has no solutions if \( k\neq 3 \) and if
      \( k=3 \) then it has infinitely many solutions.
      It never has a unique solution.  
    \end{answer}
  \item 
    This system is not linear in that it says $\tan\alpha$ instead of $\alpha$
    \begin{equation*}
      \begin{linsys}{3}
         2\tan\alpha  &-  &\cos\beta  &+  &3\sin\gamma  &=  &3  \\
         4\tan\alpha  &+  &2\cos\beta &-  &2\sin\gamma  &=  &10  \\
         6\tan\alpha  &-  &3\cos\beta &+  &\sin\gamma   &=  &9  
      \end{linsys}
    \end{equation*}
    and yet we can apply Gauss's Method.
    Do so.
    Does the system have a solution?
    \begin{answer}
      Let \( x=\tan\alpha \), \( y=\cos\beta \), and \( z=\sin\gamma \).
      \begin{equation*}
        \begin{linsys}{3}
           2x  &-  &y  &+  &3z  &=  &3  \\
           4x  &+  &2y &-  &2z  &=  &10  \\
           6x  &-  &3y &+  &z   &=  &9  
        \end{linsys}
         \grstep[-3\rho_1+\rho_3]{-2\rho_1+\rho_2}
         \begin{linsys}{3}
           2x  &-  &y  &+  &3z  &=  &3  \\
               &   &4y &-  &8z  &=  &4   \\
               &   &   &   &-8z &=  &0  
         \end{linsys}
      \end{equation*}
      gives \( z=0 \), \( y=1 \), and \( x=2 \).
      Equations such as $\tan\alpha=2$ have infinitely many answers, 
      but the principle values give $\alpha\approx 1.11$~rad, $\beta= 0$~rad,
      and $\gamma= 0$~rad.
      \end{answer}
  \recommended \item 
    What conditions must the constants, the $b$'s,
    satisfy so that each of these systems has a solution?
    \textit{Hint.} 
    Apply Gauss's Method and see what happens to the right side.
    \begin{exparts*}
      \partsitem  \(
        \begin{linsys}[t]{2}
           x  &-  &3y  &=  &b_1 \\
          3x  &+  &y   &=  &b_2 \\
           x  &+  &7y  &=  &b_3 \\
          2x  &+  &4y  &=  &b_4 
        \end{linsys}   \)
      \partsitem \(
        \begin{linsys}[t]{3}
           x_1  &+  &2x_2  &+  &3x_3  &=  &b_1  \\
          2x_1  &+  &5x_2  &+  &3x_3  &=  &b_2  \\
           x_1  &   &      &+  &8x_3  &=  &b_3  
        \end{linsys}   \)
    \end{exparts*}
    \begin{answer} 
      \begin{exparts}
       \partsitem Gauss's Method
         \begin{equation*}
           \grstep[-\rho_1+\rho_3 \\ -2\rho_1+\rho_4]{-3\rho_1+\rho_2}
           \begin{linsys}{2}
              x  &-  &3y  &=  &b_1\hfill \\
                 &   &10y &=  &-3b_1+b_2\hfill \\
                 &   &10y &=  &-b_1+b_3\hfill \\
                 &   &10y &=  &-2b_1+b_4\hfill 
            \end{linsys}         
           \grstep[-\rho_2+\rho_4]{-\rho_2+\rho_3}
           \begin{linsys}{2}
              x  &-  &3y  &=  &b_1\hfill \\
                 &   &10y &=  &-3b_1+b_2\hfill \\
                 &   &0   &=  &2b_1-b_2+b_3\hfill \\
                 &   &0   &=  &b_1-b_2+b_4\hfill 
            \end{linsys}
         \end{equation*}
         shows that this system is consistent if and only if both
         \( b_3=-2b_1+b_2 \) and \( b_4=-b_1+b_2 \).
       \partsitem Reduction
         \begin{align*}
            &\grstep[-\rho_1+\rho_3]{-2\rho_1+\rho_2}
            \begin{linsys}{3}
              x_1  &+  &2x_2  &+  &3x_3  &=  &b_1\hfill  \\
                   &   &x_2   &-  &3x_3  &=  &-2b_1+b_2\hfill  \\
                   &   &-2x_2 &+  &5x_3  &=  &-b_1+b_3\hfill  
             \end{linsys}                                          \\
            &\grstep{2\rho_2+\rho_3}
            \begin{linsys}{3}
              x_1  &+  &2x_2  &+  &3x_3  &=  &b_1\hfill  \\
                   &   &x_2   &-  &3x_3  &=  &-2b_1+b_2\hfill  \\
                   &   &      &   &-x_3  &=  &-5b_1+2b_2+b_3\hfill  
             \end{linsys}
         \end{align*}
         shows that each of \( b_1 \), \( b_2 \), and \( b_3 \) can be any
         real number\Dash this system always has a solution.
      \end{exparts}   
      \end{answer}
  \item 
    True or false: a system with more unknowns than equations
    has at least one solution.
    \begin{answer}
      False.
      This system with more unknowns than equations
      \begin{equation*}
        \begin{linsys}{3}
          x  &+  &y  &+  &z  &=  &0  \\
          x  &+  &y  &+  &z  &=  &1  
        \end{linsys}
      \end{equation*}
      has no solution.   
      \end{answer}
  \item 
    Argue that any Chemistry\index{Chemistry problem} problem like
    the one that starts this subsection\Dash
    a balance the reaction problem\Dash must have many solutions.
    \begin{answer}
      The fact that we can have the same reaction 
      in two different flasks shows that twice any solution is another,
      different, solution (if a physical reaction occurs then there must be
      at least one nonzero solution).
      Likewise for three, etc.
    \end{answer}
  \recommended \item 
    Find the coefficients
    \( a \), \( b \), and \( c \) so that the graph of \( f(x)=ax^2+bx+c \) 
    passes through the points \( (1,2) \), \( (-1,6) \), and \( (2,3) \).
    \begin{answer}
      Because \( f(1)=2 \), \( f(-1)=6 \), and \( f(2)=3 \) we get
      a linear system.
      \begin{equation*}
        \begin{linsys}{3}
          1a  &+  &1b  &+  &c  &=  &2  \\
          1a  &-  &1b  &+  &c  &=  &6  \\
          4a  &+  &2b  &+  &c  &=  &3  
         \end{linsys}
      \end{equation*}
      Gauss's Method
      \begin{equation*}
         \grstep[-4\rho_1+\rho_3]{-\rho_1+\rho_2}
         \begin{linsys}{3}
            a  &+  &b  &+  &c  &=  &2  \\
               &   &-2b&   &   &=  &4  \\
               &   &-2b&-  &3c &=  &-5 
          \end{linsys}               
         \grstep{-\rho_2+\rho_3}
         \begin{linsys}{3}
            a  &+  &b  &+  &c  &=  &2  \\
               &   &-2b&   &   &=  &4  \\
               &   &   &   &-3c&=  &-9 
           \end{linsys}
      \end{equation*}
      shows that the solution is $a = 1$, $b = -2$,  and $c = 3$, 
      so the quadratic is \( f(x)=1x^2-2x+3 \).  
      \end{answer}
  \item After \nearbytheorem{th:GaussMethod} we note that multiplying a 
   row by~$0$ is not allowed because that could change a solution set. 
   Give an example of a system where multiplying a row by zero changes
   the solution set and also give an example where it doesn't.
   \begin{answer}
     Here the solution set of the starting system is~$\set{(1,1)}$
     \begin{equation*}
         \begin{linsys}{2}
            x  &+  &y  &=  &2  \\
            x  &-  &y  &=  &0
          \end{linsys}               
         \grstep{0\rho_2}
         \begin{linsys}{2}
            x  &+  &y  &=  &2  \\
               &   &0  &=  &0
          \end{linsys}               
     \end{equation*}
     while the solution set of the system on the right 
     is a proper superset because it
     contains at least two points, $(1,1)$ and~$(2,0)$.

     In this example the solution set does not change.  
     \begin{equation*}
         \begin{linsys}{2}
            x  &+  &y  &=  &2  \\
           2x  &+  &2y &=  &4
          \end{linsys}               
         \grstep{0\rho_2}
         \begin{linsys}{2}
            x  &+  &y  &=  &2  \\
               &   &0  &=  &0
          \end{linsys}               
     \end{equation*}
   \end{answer}
  \item 
    Gauss's Method works by combining the equations in a system to make new
    equations.
    \begin{exparts}
      \partsitem Can we derive the equation \( 3x-2y=5 \) by a sequence of
        Gaussian reduction steps from the equations in this system?
        \begin{equation*}
          \begin{linsys}{2}
             x  &+  &y  &=  &1  \\
            4x  &-  &y  &=  &6
          \end{linsys}
        \end{equation*}
      \partsitem Can we derive \( 5x-3y=2 \)
        from the equations in this system?
        \begin{equation*}
          \begin{linsys}{2}
            2x  &+  &2y &=  &6  \\
            3x  &+  &3y &=  &9
          \end{linsys}
        \end{equation*}
      \partsitem Can we derive \( 6x-9y+5z=-2 \)  
        from these?
        \begin{equation*}
          \begin{linsys}{3}
            2x  &+  &y  &-  &z  &=  &4  \\
            6x  &-  &3y &+  &z  &=  &5
          \end{linsys}
        \end{equation*}
    \end{exparts}
    \begin{answer} 
       \begin{exparts} 
        \partsitem
          We can represent the equation $x+y=1$ with the row vector
          $\rowvec{1 &1 &1}$, 
          represent $4x-y=6$ with $\rowvec{4 &-1 &6}$,
          and represent $3x-2y=5$ with $\rowvec{3 &-2 &5}$.
          This relationship
          \begin{equation*}
            a\rowvec{1 &1 &1}
            +b\rowvec{4 &-1 &6}
            =\rowvec{3 &-2 &5}
          \end{equation*}
          gives a linear system.
          \begin{equation*}
            \begin{linsys}{2}
              a  &+  &4b  &=  &3  \\
              a  &-  &b   &=  &-2 \\
              a  &+  &6b  &=  &5
            \end{linsys}
          \end{equation*}
          Gauss's Method gives the solution, $a = -1$ and $b = 1$.
          \begin{align*}
            \begin{amat}{2}
              1  &4  &3  \\ 
              1  &-1  &-2  \\ 
              1  &6  &5  \\ 
            \end{amat}
            &\grstep[-\rho_{1}+\rho_{3}]{-\rho_{1}+\rho_{2}}
            \begin{amat}{2}
              1  &4  &3  \\ 
              0  &-5  &-5  \\ 
              0  &2  &2  \\ 
            \end{amat}                           \\
            &\grstep{(2/5)\rho_{2}+\rho_{3}}
            \begin{amat}{2}
              1  &4  &3  \\ 
              0  &-5  &-5  \\ 
              0  &0  &0  \\ 
            \end{amat}
          \end{align*}
          In short, the equation $3x-2y=5$ is the result of 
          $-\rho_1+\rho_2$.

          \textit{Comment.} 
          This is a wrong approach, but a common one:~add $3x-2y=5$
          as a third equation in the system and see whether the result 
          has a solution or not.
          True, if there is no solution, if
          the three-equation system is inconsistent, then 
          $3x-2y=5$ is not a combination of the other two.
          But the converse does not hold;~if there is a solution
          (or many solutions)
          then the added equation must be consistent with the others,
          but that does not imply that it is a combination.
          An example is this system
          \begin{equation*}
            \begin{linsys}{3}
              x   &+  &0y  &+  &0z  &=  &1  \\
              0x  &+  &y   &+  &0z  &=  &2  \\
            \end{linsys}
          \end{equation*}
          along with the extra equation $z=3$.
          Adding the extra one to the system and applying Gauss's Method 
          gives a solution.
          But this extra equation is not a combination of the ones in the
          system, because neither of those has a~$z$ term.
        \partsitem No.
          This
          \begin{equation*}
            a\cdot\colvec{2 &2 &6} +b\cdot\colvec{3 &3 &9} = \colvec{5 &-3 &2}
          \end{equation*}
          gives a system
          \begin{equation*}
            \begin{linsys}{2}
              2a  &+  &3b  &=  &5     \\
              2a  &+  &3b  &=  &-3   \\
              6a  &+  &9b  &=  &2
            \end{linsys}
          \end{equation*}
          with no solution, since the first and second equations conflict.
        \partsitem Yes.
          Set up.
          \begin{equation*}
            c_1\cdot\colvec{2 &1 &-1 &4}
             +c_2\cdot\colvec{6 &-3 &1 &5}
             =\colvec{6 &-9 &5 &-2}
          \end{equation*}
          To the resulting linear system
          \begin{equation*}
            \begin{linsys}{2}
               2c_1  &+  &6c_2  &=  &6  \\
                c_1  &-  &3c_2  &=  &-9 \\
               -c_1  &+  &c_2   &=  &5  \\
               4c_1  &+  &5c_2  &=  &-2 
            \end{linsys}
          \end{equation*}
          apply Gauss's Method.
          \begin{equation*}
            \begin{amat}{2}
              2  &6  &6  \\ 
              1  &-3  &-9  \\ 
              -1  &1  &5  \\ 
              4  &5  &-2  \\ 
            \end{amat}
            \grstep[(1/2)\rho_{1}+\rho_{3} \\ -2\rho_{1}+\rho_{4}]{(-1/2)\rho_{1}+\rho_{2}}
            \begin{amat}{2}
              2  &6  &6  \\ 
              0  &-6  &-12  \\ 
              0  &4  &8  \\ 
              0  &-7  &-14  \\ 
            \end{amat}                
            \grstep[(-7/6)\rho_{2}+\rho_{4}]{(2/3)\rho_{2}+\rho_{3}}
            \begin{amat}{2}
              2  &6  &6  \\ 
              0  &-6  &-12  \\ 
              0  &0  &0  \\ 
              0  &0  &0  \\ 
            \end{amat}
          \end{equation*}
          Back substitution gives $c_1=-3$ and $c_2=2$, and so the given 
          equation is $-3$~times the first equation plus $2$~times the second.
      \end{exparts}  
     \end{answer}
  \item 
    Prove that, where \( a,b,c,d,e \) are real numbers
    with \( a\neq 0 \), if this linear equation
    \begin{equation*}
       ax+by=c
    \end{equation*}
    has the same solution set as this one 
    \begin{equation*}
       ax+dy=e
    \end{equation*}
    then they are the same equation.
    What if \( a=0 \)?
    \begin{answer}
      If \( a\neq 0 \) then the solution set of the first equation is
      this.
      \begin{equation*}
        \set{(x,y)\in\Re^2\suchthat 
             \text{$x=(c-by)/a=(c/a)-(b/a)\cdot y$}}
        \tag{$*$} 
      \end{equation*}
      Thus, given~$y$ we can compute the associated~$x$.
      Taking $y=0$ gives the solution $(c/a,0)$, and since the second
      equation $ax+dy=e$ is supposed to have the same solution set, 
      substituting into
      it gives that $a(c/a)+d\cdot 0=e$, so $c=e$.
      Taking $y=1$ in ($*$) gives $a((c-b)/a)+d\cdot 1=e$,
      and so $b=d$.
      Hence they are the same equation.

      When \( a=0 \) the equations can be different and still have the 
      same solution set:~e.g.,
      \( 0x+3y=6 \) and \( 0x+6y=12 \).   
     \end{answer}
  \item 
    Show that if \( ad-bc\neq 0 \) then
    \begin{equation*}
      \begin{linsys}{2}
        ax  &+  &by  &=  &j  \\
        cx  &+  &dy  &=  &k  
      \end{linsys}
    \end{equation*}
    has a unique solution.
    \begin{answer}
      We will solve the general two-equation, two-unknown system.
      Along the way we will find  
      three cases:~that $a\neq 0$, that $a=0$ and 
      $c\neq 0$, and that both $a=0$ and $c=0$.

      For the first, Gauss's Method has us starting with the reduction step
      $-(c/a)\rho_1+\rho_2$.
      For that step we must assume that \( a\neq 0 \).
      Then the reduction
      \begin{equation*}
        \grstep{-(c/a)\rho_1+\rho_2}
        \begin{linsys}{2}
          ax  &+  &by                  &=  &j \hfill\hbox{} \\
              &   &(-(cb/a)+d)y  &=  &-(cj/a)+k \hfill  
         \end{linsys}
      \end{equation*}
      shows that this system has a unique solution if and only if
      \( -(cb/a)+d\neq 0   \); remember that \( a\neq 0 \) so 
      that back substitution yields a unique \( x \)
      (observe, by the way, that \( j \) and \( k \) play no role in the
      conclusion that there is a unique solution, although if there is a 
      unique solution then they contribute to its value).
      But \( -(cb/a)+d = (ad-bc)/a \) and a fraction is not equal to \( 0 \) 
      if and only if its numerator is not equal to \( 0 \).
      Thus in this first case, there is a unique solution if and only if
      $ad-bc\neq 0$.

      Otherwise $a=0$.
      If \( a=0 \) but \( c\neq 0 \), then swap
      \begin{equation*}
        \begin{linsys}{2}
          cx  &+  &dy  &=  &k  \\
              &   &by  &=  &j  
        \end{linsys}
      \end{equation*}
      and conclude that the system has a unique solution if and only if 
      \( b\neq 0 \)
      (for this conclusion we need the case assumption that \( c\neq 0 \) 
      to get a unique
      \( x \) in back substitution).
      But in a context where \( a=0 \) and \( c\neq 0 \),
      the condition ``\( b\neq 0 \)''
      is equivalent to the condition ``\( ad-bc\neq 0 \)''.
      That finishes the second case.

      Finally,
      if both \( a \) and \( c \) are \( 0 \) then the system
      \begin{equation*}
        \begin{linsys}{2}
          0x  &+  &by  &=  &j  \\
          0x  &+  &dy  &=  &k  
        \end{linsys}
      \end{equation*}
      might have no solutions (if the second equation is not a multiple of the
      first), or it might have infinitely many solutions (if the second
      equation is a multiple of the first then for each \( y \) satisfying
      both equations, any pair \( (x,y) \) will do). 
      But it never has a unique solution.
      Note that \( a=c=0 \) gives that \( ad-bc=0 \).  
    \end{answer}
  \recommended \item 
    In the system
    \begin{equation*}
      \begin{linsys}{2}
         ax  &+  &by  &=  &c  \\
         dx  &+  &ey  &=  &f  
      \end{linsys}
    \end{equation*}
    each of the equations describes a line in the \( xy \)-plane.
    By geometrical reasoning, show that there are three possibilities:
    there is a unique solution, there is no solution, 
    and there are infinitely many solutions.
    \begin{answer}
      Recall that if a pair of lines share two distinct points then
      they are the same line. 
      That's because two points determine a line, so these
      two points determine each of the two lines, 
      and so they are the same line.

      Thus the lines can share one point (giving a unique solution), 
      share no points (giving no solutions), or
      share at least two points (which makes them the same line).  
    \end{answer}
  \item \label{ex:ProveGaussMethod}
    Finish the proof of \nearbytheorem{th:GaussMethod}.
    \begin{answer}
     We start with the reduction operation of multiplying $\rho_i$ by a nonzero
     real number $k$.
     We have that \( (s_1,\ldots,s_n) \) satisfies
     \begin{equation*}
       \begin{linsys}{4}
         a_{1,1}x_1  &+  &a_{1,2}x_2 &+  &\cdots  &+  &a_{1,n}x_n  &=  &d_1  \\
                   &   &          &  &        &   &           &\vdotswithin{=}   \\
        ka_{i,1}x_1  &+  &ka_{i,2}x_2 &+  &\cdots  &+  &ka_{i,n}x_n &=  &kd_i  \\
                   &   &           &   &        &   &            &\vdotswithin{=}   \\
         a_{m,1}x_1  &+  &a_{m,2}x_2 &+  &\cdots  &+  &a_{m,n}x_n  &=  &d_m  
       \end{linsys}
     \end{equation*}
     if and only if
     % \begin{align*}
     %    a_{1,1}s_1+a_{1,2}s_2+\cdots+a_{1,n}s_n
     %    &=d_1                                              \\
     %    &\alignedvdots                                     \\
     %    \text{and\ } ka_{i,1}s_1+ka_{i,2}s_2+\cdots+ka_{i,n}s_n
     %    &=kd_i                                              \\
     %    &\alignedvdots                                      \\
     %    \text{and\ } a_{m,1}s_1+a_{m,2}s_2+\cdots+a_{m,n}s_n
     %    &=d_m
     % \end{align*}
     $a_{1,1}s_1+a_{1,2}s_2+\cdots+a_{1,n}s_n=d_1$ 
     and \ldots{} $ka_{i,1}s_1+ka_{i,2}s_2+\cdots+ka_{i,n}s_n=kd_i$
     and \ldots{} 
     $a_{m,1}s_1+a_{m,2}s_2+\cdots+a_{m,n}s_n=d_m$
     by the definition of `satisfies'.
     Because \( k\neq 0 \), that's true if and only if
     % \begin{align*}
     %    a_{1,1}s_1+a_{1,2}s_2+\cdots+a_{1,n}s_n
     %    &=d_1                                              \\
     %    &\alignedvdots                                     \\
     %    \text{and\ } a_{i,1}s_1+a_{i,2}s_2+\cdots+a_{i,n}s_n
     %    &=d_i                                              \\
     %    &\alignedvdots                                      \\
     %    \text{and\ } a_{m,1}s_1+a_{m,2}s_2+\cdots+a_{m,n}s_n
     %    &=d_m
     % \end{align*}
     $a_{1,1}s_1+a_{1,2}s_2+\cdots+a_{1,n}s_n=d_1$
     and \ldots{}
     $a_{i,1}s_1+a_{i,2}s_2+\cdots+a_{i,n}s_n=d_i$
     and \ldots{}
     $a_{m,1}s_1+a_{m,2}s_2+\cdots+a_{m,n}s_n=d_m$
     (this is simple canceling of a nonzero factor 
     on both sides of the $i$-th equation).
     In turn this says that \( (s_1,\ldots,s_n) \) solves this system.
     \begin{equation*}
       \begin{linsys}{4}
         a_{1,1}x_1  &+  &a_{1,2}x_2 &+  &\cdots  &+  &a_{1,n}x_n  &=  &d_1  \\
                     &   &           &   &        &   &            &\vdotswithin{=}   \\
         a_{i,1}x_1  &+  &a_{i,2}x_2 &+  &\cdots  &+  &a_{i,n}x_n  &=  &d_i  \\
                     &   &           &   &        &   &            &\vdotswithin{=}   \\
         a_{m,1}x_1  &+  &a_{m,2}x_2 &+  &\cdots  &+  &a_{m,n}x_n  &=
              &d_m  
         \end{linsys}
     \end{equation*}

     We finish by addressing the combination operation $k\rho_i+\rho_j$
     (recall the restriction on this operation that $i\neq j$). 
     The tuple
     \( (s_1,\ldots,s_n) \) satisfies this system
     % (the equation with the complicated coefficients is the $j$-th equation)
     \begin{equation*}
       \begin{linsys}{4}
         a_{1,1}x_1             &+  &\cdots  &+  &a_{1,n}x_n  &=  &d_1\hfill\hbox{} \\
                                &   &        &   &            &\vdotswithin{=}   \\
         a_{i,1}x_1             &+  &\cdots  &+  &a_{i,n}x_n  &=  &d_i\hfill\hbox{} \\
                                &   &        &   &            &\vdotswithin{=}   \\
         (ka_{i,1}+a_{j,1})x_1  &+  &\cdots  &+  &(ka_{i,n}+a_{j,n})x_n
               &=  &kd_i+d_j \hfill \\
                                &   &        &   &            &\vdotswithin{=}   \\
         a_{m,1}x_1             &+   &\cdots  &+  &a_{m,n}x_n  &=  
          &d_m\hfill\hbox{} 
        \end{linsys}
     \end{equation*}
     if and only if
     $a_{1,1}s_1+\cdots+a_{1,n}s_n=d_1$
     and \ldots{}
     $a_{i,1}s_1+\cdots+a_{i,n}s_n=d_i$
     and \ldots{} 
     $(ka_{i,1}+a_{j,1})s_1+\cdots+(ka_{i,n}+a_{j,n})s_n=kd_i+d_j$
     and \ldots{}
     $a_{m,1}s_1+a_{m,2}s_2+\cdots+a_{m,n}s_n=d_m$,
     again by the definition of `satisfies'.
     Multiply both sides of equation~\( i \) by $-k$ and add it to 
     equation~\( j \).
     The previous compound statement holds if and only if
     $a_{1,1}s_1+\cdots+a_{1,n}s_n=d_1$
     and \ldots{}
     $a_{i,1}s_1+\cdots+a_{i,n}s_n=d_i$
     and\ldots{} $
     (-ka_{i,1}s_1+ka_{i,1}+a_{j,1})s_1+\cdots+(-ka_{i,n}s_n+ka_{i,n}+a_{j,n})s_n
         =-kd_ikd_i+d_j$
      and\ldots{} $a_{m,1}s_1+\cdots+a_{m,n}s_n=d_m$.
     After cancellation (see the note below), 
     this says that \( (s_1,\ldots,s_n) \) solves
     \begin{equation*}
       \begin{linsys}{4}
         a_{1,1}x_1  &+   &\cdots  &+  &a_{1,n}x_n  &=  &d_1  \\
                     &    &        &   &            &\vdotswithin{=}   \\
         a_{i,1}x_1  &+   &\cdots  &+  &a_{i,n}x_n  &=  &d_i  \\
                     &    &        &   &            &\vdotswithin{=}   \\
         a_{j,1}x_1  &+  &\cdots  &+  &a_{j,n}x_n  &=  &d_j  \\
                     &   &        &   &            &\vdotswithin{=}   \\
         a_{m,1}x_1  &+  &\cdots  &+  &a_{m,n}x_n  &=
              &d_m\hfill\hbox{}
       \end{linsys}
     \end{equation*}  
     as required.

     \textit{Note}  As mentioned, 
     the row combination operation is subject to the restriction 
     that \( i\neq j \).
     Suppose that 
     we are considering the case of a row combination where $i=j$.
     Then in  row~$j$ (which is also row~$i$) the term on the left
     in column~$c$ is $(ka_{i,c}+a_{i,c})\cdot s_c=(k+1)a_{i,c}s_c$.
     That line will not be equivalent to the one where the term is $a_{i,c}s_c$
     if $k= -1$.)
   \end{answer}
  \item 
    Is there a two-unknowns
    linear system whose solution set is all of \( \Re^2 \)?
    \begin{answer}
      Yes, this one-equation system:
      \begin{equation*}
         0x+0y=0
      \end{equation*}
      is satisfied by every \( (x,y)\in\Re^2 \)\!.  
    \end{answer}
  \recommended \item 
    Are any of the operations used in Gauss's Method
    redundant?
    That is, can we make any of the operations from a combination
    of the others?
    \begin{answer}
      Yes.
      This sequence of operations swaps rows \( i \) and \( j \).
      \begin{equation*}
         \grstep{\rho_i+\rho_j}
         \repeatedgrstep{-\rho_j+\rho_i}
         \repeatedgrstep{\rho_i+\rho_j}
         \repeatedgrstep{-1\rho_i}
      \end{equation*}  
      This illustrates the effect of the sequence.
      \begin{multline*}
        \begin{linsys}{2}
          ax  &+  &by  &=  &e  \\
          cx  &+  &dy  &=  &f
        \end{linsys}
        \grstep{\rho_1+\rho_2}
        \begin{linsys}{2}
          ax      &+  &by      &=  &e  \\
          (a+c)x  &+  &(b+d)y  &=  &e+f
        \end{linsys}                      
        \grstep{-\rho_2+\rho_1}               \\
        \begin{linsys}{2}
          -cx      &-  &dy      &=  &-f  \\
          (a+c)x   &+  &(b+d)y  &=  &e+f
        \end{linsys}                        
        \grstep{\rho_1+\rho_2}
        \begin{linsys}{2}
          -cx      &-  &dy      &=  &-f  \\
          ax       &+  &by      &=  &e
        \end{linsys}
        \grstep{-1\rho_i}
        \begin{linsys}{2}
          cx      &+  &dy      &=  &f  \\
          ax      &+  &by      &=  &e
        \end{linsys}
      \end{multline*}
      The row-swap operation is redundant in the presence of the other two.
     \end{answer}
  \item 
    Prove that each operation of Gauss's Method is reversible.
    That is, show that if two systems are related by a row operation
    $S_1\rightarrow S_2$ then there is a row operation to go back
    $S_2\rightarrow S_1$.
    \begin{answer}
      Reverse a row swap $\rho_i\leftrightarrow\rho_j$ by swapping 
      back $\rho_j\leftrightarrow\rho_i$.
      Reverse the $k\rho_i$ step of multiplying  \( k\neq 0  \) 
      on both sides of a row
      by  dividing through,~$(1/k)\rho_i$.

      The row combination case is the nontrivial one.
      The operation $k\rho_i+\rho_j$
      results in this $j$-th row.
      \begin{equation*}
        (k\cdot a_{i,1}+a_{j,1})x_1+
         \cdots+(k\cdot a_{i,n}+a_{j,n})x_n=k\cdot d_i+d_j
      \end{equation*}
      The $i$-th row unchanged 
      (because of the $i\neq j$ restriction on the combination operation).
      Because of this, the operation $-k\rho_i+\rho_j$ 
      returns the $j$-th row to its original state.

      \textit{Comment}. 
       That \( i=j \) condition on the $k\rho_i+\rho_j$ 
       is needed, or this would happen
       \begin{equation*}
         \begin{linsys}{2}
           3x  &+  &2y  &=  &7  
         \end{linsys}
         \grstep{2\rho_1+\rho_1}
         \begin{linsys}{2}
           9x  &+  &6y  &=  &21 
          \end{linsys}                 
         \grstep{-2\rho_1+\rho_1}
         \begin{linsys}{2}
          -9x  &-  &6y  &=  &-21 
         \end{linsys}
       \end{equation*}
       and so the result wouldn't hold.
    \end{answer}
  \puzzle \item  
    \cite{Anton}
    A box holding pennies, nickels and dimes contains
    thirteen coins with a total value of \( 83 \) cents.
    How many coins of each type are in the box?
    (These are US coins;~a
    penny is $1$~cent, a nickel is $5$~cents, and
    a dime is $10$~cents.)
    \begin{answer}
      Let \( p \), \( n \), and \( d \) be the number of
      pennies, nickels, and dimes.
      For variables that are real numbers, this system
      \begin{equation*}
         \begin{linsys}{3}
              p  &+ &n   &+  &d   &=  &13   \\
              p  &+ &5n  &+  &10d &=  &83    
         \end{linsys}
         \grstep{-\rho_1+\rho_2}
         \begin{linsys}{3}
              p  &+ &n   &+  &d   &=  &13   \\
                 &  &4n  &+  &9d  &=  &70    
          \end{linsys}
      \end{equation*}
      has more than one solution; in fact, it has infinitely many of them.
      However, it has a limited number of solutions in which \( p \), \( n \),
      and \( d \) are non-negative integers.
      Running through \( d=0 \), \ldots, \( d=8 \), whether by hand or on the
      computer 
      \begin{lstlisting}
sage: def solve(d):
....:     n = (70-(9*d))/4
....:     p = 13 - d -n
....:     if n.is_integer() and p.is_integer():
....:         print("d=", d, "gives integers: n=", n, " p=", p)
....:     else:
....:         print("d=", d, "gives nonintegers")
....: 
sage: for d in range(9):
....:     solve(d)
....: 
d= 0 gives nonintegers
d= 1 gives nonintegers
d= 2 gives integers: n= 13  p= -2
d= 3 gives nonintegers
d= 4 gives nonintegers
d= 5 gives nonintegers
d= 6 gives integers: n= 4  p= 3
d= 7 gives nonintegers
d= 8 gives nonintegers
      \end{lstlisting}
      shows that 
      \( (p,n,d)=(3,4,6) \)
      is the only solution using natural numbers.  
    \end{answer}
  \puzzle \item 
    \cite{ContestProb1955no38}
    Four positive integers are given.
    Select any three of the integers, find their arithmetic average,
    and add this result to the fourth integer.
    Thus the numbers 29, 23, 21, and 17 are obtained.
    One of the original integers is:
    \begin{exparts*}
      \partsitem 19
      \partsitem 21
      \partsitem 23
      \partsitem 29
      \partsitem 17
    \end{exparts*}
    \begin{answer}
      The wording of the question does not tell us which of the values 
      $29$, $23$, $21$, or~$17$ is obtained,
      say, from $(1/3)(a+b+c)+d$.
      However, 
      if we assume that the question is correct\Dash that 
      there is a fixed choice 
      among these five numbers $19$, $21$, $23$, $29$, and~$17$
      so that for every way to assign values to equations then that
      choice will be one of the four variables\Dash then we can try assigning
      the values in any way and apply Gauss's Method.
      Solving the system
      \begin{equation*}
        \begin{linsys}{2}
        (1/3)(a+b+c)  &+  &d  &=  &29  \\
        (1/3)(b+c+d)  &+  &a  &=  &23  \\
        (1/3)(c+d+a)  &+  &b  &=  &21  \\
        (1/3)(d+a+b)  &+  &c  &=  &17
        \end{linsys}
      \end{equation*}
      via
      \begin{equation*}
        \grstep[-\rho_{1}+\rho_{3} \\ -\rho_{1}+\rho_{4}]{-3\rho_{1}+\rho_{2}}
        % \begin{amat}{4}
        %   1/3  &1/3  &1/3  &1  &29  \\ 
        %   0  &-2/3  &-2/3  &-8/3  &-64  \\ 
        %   0  &2/3  &0  &-2/3  &-8  \\ 
        %   0  &0  &2/3  &-2/3  &-12  \\ 
        % \end{amat}                            \\                      
        \grstep{\rho_{2}+\rho_{3}}
          % \begin{amat}{4}
          %   1/3  &1/3  &1/3  &1  &29  \\ 
          %   0  &-2/3  &-2/3  &-8/3  &-64  \\ 
          %   0  &0  &-2/3  &-10/3  &-72  \\ 
          %   0  &0  &2/3  &-2/3  &-12  \\ 
          % \end{amat}                           \\
        \grstep{\rho_{3}+\rho_{4}}
        \begin{amat}{4}
          1/3  &1/3  &1/3  &1  &29  \\ 
          0  &-2/3  &-2/3  &-8/3  &-64  \\ 
          0  &0  &-2/3  &-10/3  &-72  \\ 
          0  &0  &0  &-4  &-84  \\ 
        \end{amat}                               
      \end{equation*}
      gives $a=12$, $b=9$, $c=3$, $d=21$.
      Since it is the only one of these numbers that appears, 
      the second item, $21$, is the correct answer.

      \textit{Comment}.
      That was perhaps cheating
      (and we got lucky in that if more than one of the available choices had
      been an option then we would have had to do more work).
      To properly solve we can add all four equations, giving
      $2a+2b+2c+2d=90$, and multiply by $1/6$ to get
      $(1/3)a+(1/3)b+(1/3)c+(1/3)d=15$.
      Compare this with the expression $a+(1/3)b+(1/3)c+(1/3)d$ to conclude that
      $(2/3)a$ is either $29-15=14$, $23-15=8$, $21-15=6$, or~$17-15=2$.
      Thus $a$ is either $21$, $12$, $9$, or~$3$.
      The same reasoning applies to the other variables.
     \end{answer}
  \puzzle \item  
      \cite{Monthly35p47}
      Laugh at this:  \( \mbox{AHAHA}+\mbox{TEHE}=\mbox{TEHAW} \).
      It resulted from substituting a code letter for each digit of a simple
      example in addition, and it is required to identify the letters
      and prove the solution unique.
      \begin{answer}
        \answerasgiven
        A comparison of the units and hundreds columns of this
        addition shows that there must be a carry from the tens column.
        The tens column then tells us that \( A<H \), so there
        can be no carry from the units or hundreds columns.
        The five columns then give the following five equations.
        \begin{align*}
          A+E  &=  W  \\
          2H   &=  A+10  \\
          H    &=  W+1  \\
          H+T  &=  E+10  \\
          A+1  &=  T
        \end{align*}
        The five linear equations in five unknowns, if solved simultaneously,
        produce the unique solution: \( A=4 \), \( T=5 \), \( H=7 \),
        \( W=6 \) and \( E=2 \), so that the original example in addition
        was \( 47474+5272=52746 \).  
      \end{answer}
  \puzzle \item 
     \cite{Wohascum2}
     The Wohascum County Board of Commissioners, which has 20 members, 
     recently had to elect a President.
     There were three candidates ($A$, $B$, and $C$); on each ballot
     the three
     candidates were to be listed in order of preference, with no abstentions.
     It was found that 11 members, a majority, preferred $A$ over $B$
     (thus the other 9 preferred $B$ over $A$).
     Similarly, it was found that 12 members preferred $C$ over $A$.
     Given these results, it was suggested that $B$ should withdraw, to enable
     a runoff election between $A$ and $C$.
     However, $B$ protested, and it was then found that 14 members preferred
     $B$ over $C$!
     The Board has not yet recovered from the resulting confusion.
     Given that every possible order of $A$, $B$, $C$ appeared on at least 
     one ballot, how many members voted for $B$ as their first choice?
     \begin{answer}
       \answerasgiven{} 
       \textit{Some additional material was added from \cite{joriki}.}
       Eight commissioners voted for $B$.
       To see this, we will use the given information to study how many voters
       chose each order of $A$, $B$, $C$.

       The six orders of preference are $ABC$, $ACB$, $BAC$, $BCA$, $CAB$,
       $CBA$; assume they receive $a$, $b$, $c$, $d$, $e$, $f$ votes 
       respectively.
       We know that
       \begin{equation*}
         \begin{linsys}{3}
           a  &+  &b  &+  &e  &=  &11  \\
           d  &+  &e  &+  &f  &=  &12  \\
           a  &+  &c  &+  &d  &=  &14
         \end{linsys}
       \end{equation*}
       from the number preferring $A$ over $B$, the number preferring
       $C$ over $A$, and the number preferring $B$ over $C$.
       Because 20 votes were cast, 
       $a+b+\cdots+f=20$.
       Subtracting the sum of the three above equations from twice the 
       prior equation gives $b+c+f=3$.
       We've specified that each preference order got at 
       least one vote, so that means $b=c=f=1$. 

       From the above three equations the complete solution is then 
       $a=6$, $b=1$, $c=1$, $d=7$, $e=4$, and~$f=1$,
       as we can find with Gauss's Method.
       The number of commissioners voting for $B$ as their first choice 
       is therefore $c+d=1+7=8$.

       \par\noindent {\em Comment.}
       % The answer to this question would have been the same had we known only
       % that \emph{at least} $14$ commissioners preferred $B$ over $C$.
       %
       The seemingly paradoxical nature of the commissioner's preferences
       ($A$ is preferred to $B$, and $B$ is preferred to $C$, and $C$ is 
       preferred to $A$), an example of ``non-transitive dominance,'' is
       common when individual choices are pooled.
     \end{answer}
  \puzzle \item   
     \cite{Monthly63p93}
    ``This system
     of \( n \) linear equations with
     \( n \) unknowns,'' said the Great Mathematician, ``has a curious
     property.''

     ``Good heavens!'' said the Poor Nut,  ``What is it?''

     ``Note,'' said the Great Mathematician, ``that the constants are in
     arithmetic progression.''

     ``It's all so clear when you explain it!'' said the Poor Nut.
     ``Do you mean like \( 6x+9y=12 \) and \( 15x+18y=21 \)?''

     ``Quite so,'' said the Great Mathematician, pulling out his bassoon.
     ``Indeed, the system has a unique solution.
     Can you find it?''

     ``Good heavens!'' cried the Poor Nut, ``I am baffled.''

     Are you?
     \begin{answer}
       \answerasgiven
       \textit{We have not used ``dependent'' yet; 
       it means here that Gauss's
       Method shows that there is not a unique solution.}
       If \( n\geq 3 \) the system is dependent and the solution is not
       unique.
       Hence \( n<3 \).
       But the term ``system'' implies \( n>1 \).
       Hence \( n=2 \).
       If the equations are
       \begin{equation*}
         \begin{linsys}{2}
              ax  &+ &(a+d)y  &=  &a+2d  \\
         (a+3d)x  &+ &(a+4d)y &=  &a+5d  
         \end{linsys}
       \end{equation*}
       then \( x=-1 \), \( y=2 \).  
    \end{answer}
\end{exercises}







\subsection{Describing the Solution Set}
Linear systems come in three types:~ones with no solution, 
ones with a single solution, 
and ones with many solutions.
For the first two types the solution sets are easy to describe.
Solution sets are a challenge to describe only when they contain many elements.

\begin{example}
\nearbytheorem{th:GaussMethod} shows that $(x,y,z)$  
satisfies the first system if and only if it satisfies the
other two.
\begin{align*}
  \begin{linsys}{3}
    2x  &   &   &+  &z  &=  &3 \\
     x  &-  &y  &-  &z  &=  &1 \\
    3x  &-  &y  &   &   &=  &4 
  \end{linsys}
  &\grstep[-(3/2)\rho_1 +\rho_3]{-(1/2)\rho_1+\rho_2}
  \begin{linsys}{3}
     2x  &   &   &+  &z      &=  &3    \\
         &   &-y &-  &(3/2)z &=  &-1/2 \\
         &   &-y &-  &(3/2)z &=  &-1/2 
   \end{linsys}                                   \\
  &\grstep{-\rho_2+\rho_3}
  \begin{linsys}{3}
     2x  &   &   &+  &z      &=  &3    \\
         &   &-y &-  &(3/2)z &=  &-1/2 \\
         &   &   &   &0      &=  &0    
   \end{linsys}
\end{align*}
So we can describe the solution set of the starting system 
\begin{equation*}
\set{(x,y,z)\suchthat\text{$2x+z=3$ and $x-y-z=1$ and $3x-y=4$}}
\end{equation*}
in this way.
\begin{equation*}
  \set{(x,y,z)\suchthat\text{$2x+z=3$ and $-y-3z/2=-1/2$}}
  \tag{$*$}
\end{equation*}
This is simpler because 
it has two equations instead of three 
but it still has hard to understand interactions among the variables.
For instance, is there a solution with $x=4$ and~$y=3$?
It turns out that there is not, because with $x=4$ we have $z=-5$, which
is incompatible with~$y=3$.

For a better description, start with the 
variable that does not lead any equation,~$z$. 
Use it to express
the variables that do lead, $x$ and~$y$.
The bottom equation gives
$y=(1/2)-(3/2)z$ 
and the top equation gives
$x=(3/2)-(1/2)z$.
We now have expressed the solution set as this set of triples.   
\begin{equation*}
  \set{\bigl( (3/2)-(1/2)z,\,(1/2)-(3/2)z,\,z\bigr)\suchthat z\in\Re}
  \tag{$**$}
\end{equation*}
\end{example}

An advantage of ($**$) over  ($*$) 
is the `$z\in\R$', that is, $z$ is unrestricted in that it can be any real number.
This makes deciding what is and is not in the solution set much easier.
For instance, setting $z=2$ shows that $(1/2,-5/2,2)$ is a solution.
And, $(3,2,1)$ is not a solution because if $z=1$ then the first component
in ($**$) is not~$3$.

Another advantage of ($**$) is that it shows how many solutions
there are.
There is a different solution for each~$z\in\R$, because each solution
has a different third coordinate. 
Therefore the number of solutions is infinite.

\begin{definition} \label{df:FreeVars}
%<*df:FreeVars>
In an echelon form linear system the variables that do not lead
any row are 
\definend{free}.\index{echelon form!free variable}\index{free variable}
%</df:FreeVars>
\end{definition}

\begin{example}   \label{ex:Parametrize2}
Reduction of a linear system can end with more than one free variable.
Gauss's Method on this system
leaves  \( x \) and \( y \) leading, and both \( z \) and~\( w \) free.
\begin{align*}
   \begin{linsys}{4}
               x  &+  &y   &+  &z   &-  &w   &=  &1  \\
                  &   &y   &-  &z   &+  &w   &=  &-1 \\
              3x  &   &    &+  &6z  &-  &6w  &=  &6  \\
                  &   &-y  &+  &z   &-  &w   &=  &1  
   \end{linsys}
  &\grstep{-3\rho_1 +\rho_3}
  \begin{linsys}{4}
     x  &+  &y   &+  &z   &-  &w   &=  &1  \\
        &   &y   &-  &z   &+  &w   &=  &-1 \\
        &   &-3y &+  &3z  &-  &3w  &=  &3  \\
        &   &-y  &+  &z   &-  &w   &=  &1  
  \end{linsys}                                      \\
  &\grstep[\rho_2 +\rho_4]{3\rho_2 +\rho_3}
  \begin{linsys}{4}
     x  &+  &y   &+  &z   &-  &w   &=  &1  \\
        &   &y   &-  &z   &+  &w   &=  &-1 \\
        &   &    &   &    &   &0   &=  &0  \\
        &   &    &   &    &   &0   &=  &0  
   \end{linsys}
\end{align*}
To get the description that we prefer, start from the bottom.
First express the leading variable~$y$ in terms of
$z$ and~$w$, as $y=-1+z-w$.
Moving up to the top equation,
substituting for $y$ gives
$x+(-1+z-w)+z-w=1$ and solving for~$x$ leaves $x=2-2z+2w$.
The solution set 
\begin{equation*}
   \set{(2-2z+2w,-1+z-w,z,w)\suchthat z,w\in\Re}
  % \tag{$**$}
\end{equation*}
has the leading variables expressed in terms of the variables that are free.
\end{example}

\begin{example}    \label{ex:Parametrize1}
The list of leading variables may skip some columns.
After this reduction
\begin{align*}
  \begin{linsys}{4}
              2x  &-  &2y  &   &    &   &    &=  &0  \\
                  &   &    &   &z   &+  &3w  &=  &2  \\
              3x  &-  &3y  &   &    &   &    &=  &0  \\
               x  &-  &y   &+  &2z  &+  &6w  &=  &4  
  \end{linsys}
  &\grstep[-(1/2)\rho_1+ \rho_4]{-(3/2)\rho_1 +\rho_3}
  \begin{linsys}{4}
     2x  &-  &2y  &\spaceforemptycolumn   &    &   &    &=  &0  \\
         &   &    &   &z   &+  &3w  &=  &2  \\
         &   &    &   &    &   &0   &=  &0  \\
         &   &    &   &2z  &+  &6w  &=  &4  
   \end{linsys}                                    \\
  &\grstep{-2\rho_2 +\rho_4}
  \begin{linsys}{4}
     2x  &-  &2y  &\spaceforemptycolumn   &    &   &    &=  &0  \\
         &   &    &   &z   &+  &3w  &=  &2  \\
         &   &    &   &    &   &0   &=  &0  \\
         &   &    &   &    &   &0   &=  &0  
   \end{linsys}
\end{align*}
the leading variables are $x$ and~$z$, not $x$ and~$y$.
The free variables are $y$ and~$w$. 
The solution set is
$\set{ (y,y,2-3w,w)\suchthat y,w\in\Re }$.
For instance, \( (1,1,2,0) \) satisfies the system\Dash take 
$y=1$ and $w=0$.
And, the four-tuple \( (1,0,5,4) \) is not a solution
since its first coordinate does not equal its second.
\end{example}

%<*df:Parameter>
A variable used to describe a family of solutions
is a \definend{parameter}.\index{parameter}  
%</df:Parameter>
We say that the solution set in the prior example 
is \definend{parametrized\/}\index{parametrized} 
with $y$ and~$w$.

The terms `parameter' and `free variable' do not mean the same thing.
In the prior example
$y$ and~$w$ are free because in the echelon form system they
do not lead. 
They are parameters because 
we used them to describe the set of solutions.
Had we instead 
rewritten the second equation as $w=2/3-(1/3)z$ then
the free variables would still be $y$ and~$w$ but the parameters 
would be $y$ and~$z$.

In the rest of this book 
we will solve linear systems by bringing them to
echelon form and then parametrizing with the free variables.

\begin{example}
This is another system with infinitely many solutions.
\begin{align*}
   \begin{linsys}{4}
               x  &+  &2y  &   &   &   &   &=  &1  \\
              2x  &   &    &+  &z  &   &   &=  &2  \\
              3x  &+  &2y  &+  &z  &-  &w  &=  &4  
   \end{linsys}
  &\grstep[-3\rho_1 +\rho_3]{-2\rho_1+\rho_2}
  \begin{linsys}{4}
     x  &+  &2y  &   &   &   &   &=  &1  \\
        &   &-4y &+  &z  &   &   &=  &0  \\
        &   &-4y &+  &z  &-  &w  &=  &1  
   \end{linsys}                                    \\
  &\grstep{-\rho_2+\rho_3}
  \begin{linsys}{4}
     x  &+  &2y  &   &   &   &   &=  &1  \\
        &   &-4y &+  &z  &   &   &=  &0  \\
        &   &    &   &   &   &-w &=  &1  
   \end{linsys}
\end{align*}
The leading variables are \( x \), \( y \), and \( w \).
The variable \( z \) is free.
Notice that although the system has infinitely many 
solutions, $w$ is fixed at~$-1$.
To parametrize, write \( w \) in terms of \( z \) with \( w=-1+0z \).
Then \( y=(1/4)z \).
In the first 
equation, substitute for \( y \) to get \( x=1-(1/2)z \).
The solution set is $\set{(1-(1/2)z,(1/4)z,z,-1)\suchthat z\in\Re}$.
\end{example}

We finish this subsection by developing the final
notation for linear systems 
and their solution sets that we will use for the rest of this book.
It is both efficient and clear.

\begin{definition}\label{df:matrix}
%<*df:matrix>
A \definend{matrix}\index{matrix}
is a rectangular array of numbers.
Where it has
\( m \)~\definend{rows}\index{matrix!row}\index{row} 
and \( n \)~\definend{columns}\index{matrix!column}\index{column}
we say that it is \( \nbym{m}{n} \).
Each number in the matrix is an 
\definend{entry}\index{matrix!entry}\index{entry, matrix}.
%</df:matrix>
\end{definition}

We usually denote a matrix with an uppercase % One word is the most common: https://books.google.com/ngrams/graph?content=lowercase%2Clower%20case%2Clower-case&year_start=1800&year_end=2022&corpus=en-US&smoothing=3&case_insensitive=false
roman letter.
For instance,
\begin{equation*}
  A=
  \begin{mat}[r]
    1  &2.2  &5  \\
    3  &4    &-7
  \end{mat}
\end{equation*}
has $2$~rows and $3$~columns and so
is \( \nbym{2}{3} \)
(read aloud as ``$2$ by~$3$'';
always state the number of rows first).
The parentheses are there 
so that when two matrices are adjacent,
we can tell where one ends and the other begins.
We name matrix entries with the corresponding lowercase letter,
so that the entry in the second row and third column 
of the above array is \( a_{2,3}=-7 \).
Note that the order of the subscripts matters;~for 
instance, \( a_{1,2}=2.2 \), while \( a_{2,1}=3 \). 
We write
the set of all \( \nbym{m}{n} \)
matrices as \( \matspace_{\nbym{m}{n}} \).\index{matrices|set of}

Do Gauss's Method with matrices in the same
way as with systems of equations:~a matrix row's
\definend{leading entry}\index{echelon form!leading entry}%
\index{leading!entry} % 
is its first nonzero entry (if it has one) and
we perform row operations to arrive at 
\definend{matrix echelon form},\index{echelon form!matrix}%
\index{matrix!echelon form} %
where any leading entry in lower rows is to the right of those in 
the rows above.

\begin{example}
We can abbreviate this linear system
with this matrix.
\begin{equation*}
  \begin{linsys}{3}
    x  &+  &2y  &   &    &=  &4   \\
       &   &y   &-  &z   &=  &0   \\
    x  &   &    &+  &2z  &=  &4   
  \end{linsys}
  \qquad\qquad
  \begin{amat}[r]{3}
  1  &2  &0  &4  \\
  0  &1  &-1 &0  \\
  1  &0  &2  &4
  \end{amat}
\end{equation*}
% \begin{equation*}
%     \begin{amat}[r]{3}
%       1  &2  &0  &4  \\
%       0  &1  &-1 &0  \\
%       1  &0  &2  &4
%     \end{amat}
% \end{equation*}
The vertical bar reminds a reader of the difference between the 
coefficients on the system's left hand side and the constants on the right.
With a bar, this is an
\definend{augmented\/}\index{matrix!augmented}\index{augmented matrix} matrix.
\begin{equation*}
    \begin{amat}[r]{3}
      1  &2  &0  &4  \\
      0  &1  &-1 &0  \\
      1  &0  &2  &4
    \end{amat}
  \grstep{-\rho_1 +\rho_3}
  \begin{amat}[r]{3}
       1  &2  &0  &4  \\
       0  &1  &-1 &0  \\
       0  &-2 &2  &0
     \end{amat}                        
  \grstep{2\rho_2 +\rho_3}
  \begin{amat}[r]{3}
       1  &2  &0  &4  \\
       0  &1  &-1 &0  \\
       0  &0  &0  &0
     \end{amat}
\end{equation*}
On the right, the second row stands for $y-z=0$ and the first stands for
$x+2y=4$, so the solution set is
\( \set{(4-2z,z,z)\suchthat z\in\Re} \).
\end{example}

Matrix notation lightens
the clerical load such as the copying of variables and the
writing of $+$'s and $=$'s. 
We will also use it to
clarify the descriptions of solution sets.
\nearbyexample{ex:Parametrize2}'s
$\set{(2-2z+2w,-1+z-w,z,w)\suchthat z,w\in\Re}$ 
is hard to read.
We will rewrite it to group all of the
constants together, all of the
coefficients of~\( z \) together, and all of the coefficients of~\( w \)
together.
We write them vertically, in single-column matrices.
\begin{equation*}
  \set{\colvec{x \\ y \\ z \\ w}
       =\colvec{2 \\ -1 \\ 0 \\ 0}
       +\colvec{-2 \\ 1 \\ 1 \\ 0}\cdot z
       +\colvec{2 \\ -1 \\ 0 \\ 1}\cdot w
       \suchthat z,w\in\Re}
\end{equation*}
For instance, the top line says that \( x=2-2z+2w \)
and the second line says that \( y= -1+z-w \).
(Our next section gives a geometric interpretation to help 
picture the solution sets.)

\begin{definition}\label{df:vector}
%<*df:vector>
A 
\definend{column vector},\index{column!vector}\index{vector!column}
often just called a \definend{vector},\index{vector} 
is a matrix with a single column.
A matrix with a single row is a
\definend{row vector}.\index{row!vector}\index{vector!row}
In a vector the entries are sometimes called
\definend{components}.\index{component of a vector}\index{vector!component}
A column or row vector whose components are all zeros is a 
\definend{zero vector}.\index{zero vector}\index{vector!zero}
%</df:vector>
\end{definition}

Vectors are an exception to the convention of representing matrices with 
uppercase roman letters.
We use lowercase roman or greek letters overlined
with an arrow:
\( \vec{a} \), \( \vec{b} \), \ldots{} or
\( \vec{\alpha} \), \( \vec{\beta} \), \ldots{}
(also common is boldface:
{\boldmath \( a \)} or {\boldmath \( \alpha \)}).
For instance, this is a column vector
with a third component of \( 7 \).
\begin{equation*}
  \vec{v}=
  \colvec{ 1  \\  3  \\ 7}
\end{equation*}

A zero vector is denoted \( \zero \).
There are many different zero vectors\Dash the
one-tall zero vector, the two-tall zero vector, etc.\Dash but
nonetheless we will often say ``the'' zero vector, expecting 
that the size will be clear from the context.

\begin{definition}
The linear equation
\( a_1x_1+a_2x_2+\,\cdots\,+a_nx_n=d \)
with unknowns \( x_1,\ldots\,,x_n \)
is \definend{satisfied}\index{vector!satisfies an equation}%
\index{linear equation!satisfied by a vector} by the vector
\begin{equation*}
  \vec{s}=\colvec{s_1 \\ \vdotswithin{s_1} \\ s_n}
\end{equation*}
if \( a_1s_1+a_2s_2+\,\cdots\,+a_ns_n=d \).
A vector satisfies a system if it satisfies each equation in 
that system.
\end{definition}

The vector description style that we will use for solution sets
involves adding vectors and 
multiplying them by real numbers.
Before we do examples, 
we should formally define these operations.

\begin{definition} \label{df:VectorSum} \label{df:VectorScalarMultiplication}
%<*df:VectorSum>
The 
\definend{sum}\index{vector!sum}\index{sum!vector}\index{addition of vectors} 
of two same-sized vectors
is the vector of the sums of components.
\begin{equation*}
  \vec{u}+\vec{v}=
  \colvec{u_1 \\ \vdotswithin{u_1} \\ u_n}
   +
  \colvec{v_1 \\ \vdotswithin{v_1} \\ v_n}
   =
  \colvec{u_1+v_1 \\ \vdotswithin{u_1+v_1} \\ u_n+v_n}
\end{equation*}
%</df:VectorSum>
%<*df:VectorScalarMultiplication>
The \definend{scalar multiplication\/}\index{vector!scalar multiple}%
\index{scalar multiple!vector} of a real number
and a vector is the vector of the multiples of components.
\begin{equation*}
  r\cdot\vec{v}=
  r\cdot\colvec{v_1 \\ \vdotswithin{v_1} \\ v_n}
  =
  \colvec{rv_1 \\ \vdotswithin{rv_1} \\ rv_n}
\end{equation*}
%</df:VectorScalarMultiplication>
Both of these operations work for matrices in general, 
not just vectors, 
provided for addition that the matrices have the same number of rows and 
columns.\index{matrix!sum}\index{sum!matrix}
\end{definition}

We write scalar multiplication either as \( r\cdot\vec{v} \) or
\( \vec{v}\cdot r \), and sometimes omit the `$\cdot$' symbol:~$r\vec{v}$.
(Do not refer to scalar multiplication 
as `scalar product' because that names a different operation.)

\begin{example}
\begin{equation*}
  \colvec{2 \\ 3 \\ 1}
   +
  \colvec{3 \\ -1 \\ 4}
   =
  \colvec{5 \\ 2 \\ 5}
  \qquad
  7\cdot\colvec{1 \\ 4 \\ -1 \\ -3}
  =
  \colvec{7 \\ 28 \\ -7 \\ -21}
\end{equation*}
\end{example}

Observe that the definitions of addition and scalar multiplication agree
where they overlap, as with \( \vec{v} +\vec{v} = 2\vec{v} \).

\begin{example} \label{ex:ManyParamsInfManySolsSystem}
This system
\begin{equation*}
   \begin{linsys}{5}
      2x  &+  &y  &  &  &-  &w  &   &   &=  &4  \\
          &   &y  &  &  &+  &w  &+  &u  &=  &4  \\
       x  &   &   &- &z &+  &2w &   &   &=  &0  
   \end{linsys}
\end{equation*}
reduces 
\begin{align*}
  \begin{amat}[r]{5}
    2  &1  &0  &-1  &0  &4  \\
    0  &1  &0  &1   &1  &4  \\
    1  &0  &-1 &2   &0  &0
  \end{amat}
  &\grstep{-(1/2)\rho_1+\rho_3}
  \begin{amat}[r]{5}
    2  &1     &0  &-1    &0  &4  \\
    0  &1     &0  &1     &1  &4  \\
    0  &-1/2  &-1 &5/2   &0  &-2
  \end{amat}                                 \\
  &\grstep{(1/2)\rho_2+\rho_3}
  \begin{amat}[r]{5}
    2  &1     &0  &-1    &0    &4  \\
    0  &1     &0  &1     &1    &4  \\
    0  &0     &-1 &3     &1/2  &0
  \end{amat}
\end{align*}
to give the solution set
\( \set{(w+(1/2)u,4-w-u,3w+(1/2)u,w,u)\suchthat w,u\in\Re} \).
Here it is in vector form.
\begin{equation*}
  \set{\colvec{x \\ y \\ z \\ w \\ u}=
       \colvec{0 \\ 4 \\ 0 \\ 0 \\ 0}+
       \colvec{1 \\ -1 \\ 3 \\ 1 \\ 0}w+
       \colvec{1/2 \\ -1 \\ 1/2 \\ 0 \\ 1}u
       \suchthat w,u\in\Re}
\end{equation*}
Note how well this notation sets off 
the coefficients of each parameter.
For instance, the third row shows plainly that if \( u \) is
fixed then \( z \) increases three times as fast as \( w \).
Another thing shown plainly is that setting both \( w \) and \( u \) to zero
gives that
\begin{equation*}
  \colvec{x \\ y \\ z \\ w \\ u}
  =\colvec{0 \\ 4 \\ 0 \\ 0 \\ 0}
\end{equation*}
is a particular solution of the linear system.
\end{example}

\begin{example}
In the same way, the system
\begin{equation*}
   \begin{linsys}{3}
     x  &-  &y  &+  &z  &=  &1  \\
    3x  &   &   &+  &z  &=  &3  \\
    5x  &-  &2y &+  &3z &=  &5  
  \end{linsys}
\end{equation*}
reduces
\begin{align*}
  \begin{amat}[r]{3}
    1  &-1  &1  &1  \\
    3  &0   &1  &3  \\
    5  &-2  &3  &5
  \end{amat}
  &\grstep[-5\rho_1+\rho_3]{-3\rho_1+\rho_2}
  \begin{amat}[r]{3}
    1  &-1  &1  &1  \\
    0  &3   &-2 &0  \\
    0  &3   &-2 &0
  \end{amat}                                    \\
  &\grstep{-\rho_2+\rho_3}
  \begin{amat}[r]{3}
    1  &-1  &1  &1  \\
    0  &3   &-2 &0  \\
    0  &0   &0  &0
  \end{amat}
\end{align*}
to give a one-parameter solution set.
\begin{equation*}
  \set{\colvec{1 \\ 0 \\ 0}
       +\colvec{-1/3 \\ 2/3 \\ 1}z
       \suchthat z\in\Re}
\end{equation*}
As in the prior example, the vector not associated with the parameter
\begin{equation*}
   \colvec{1 \\ 0 \\ 0}
\end{equation*}
is a particular solution of the system.
\end{example}

Before the exercises, a preview.
We have focused so far on the mechanics of Gauss's Method.
Except for proving \nearbytheorem{th:GaussMethod}\Dash which 
justifies the doing the method at all by showing that it gives 
the right answers\Dash we have not stopped to consider any of the questions
that arise,
 
One question is:~can we 
always describe solution sets as above, with
a particular solution vector added to an unrestricted linear combination of 
some other vectors?
We've noted that the solution sets described in this way 
have infinitely many members,
so answering this question
would tell us about the size of solution sets.
The following subsection shows that the answer is ``yes.''
This chapter's second section then
uses that answer to describe the geometry of solution sets, what they look like.

Other questions arise from the observation that we can do Gauss's Method in 
more than one way (for instance, when swapping we may have a choice of 
rows).
\nearbytheorem{th:GaussMethod} says that we must get the same solution set
no matter how we proceed but
if we do Gauss's Method in two ways
must we get the same number of free variables in each echelon form system?
Must those be the same variables, that is, is it impossible to
solve a problem
one way to get $y$ and~$w$ free and solve it another way to get $y$ and~$z$ 
free?
The third section of this chapter answers ``yes'' to both questions.

In this way, by the end of the chapter we will not only have a 
solid grounding in the practice of Gauss's Method but 
we will also be solid on the theory.
We will know what can and cannot happen in a reduction.

\exerciseNote{%
Sets can have more than one description.
When we solve systems by hand there is a natural way to go from the
system to a solution set description, so usually
we don't find that the answer book description 
is not the same as the one that we got.
But we can run across different ones, for instance 
when using a computer algebra system to check our work.
Subsection~\ref{subsec:ComparingSetDescriptions}
shows how to verify that two differently-described sets are equal.}

\begin{exercises}
  \recommended \item  
    Find the indicated entry of the matrix,
    if it is defined.
    \begin{equation*}
      A=\begin{mat}[r]
        1  &3  &1  \\
        2  &-1 &4
      \end{mat}
    \end{equation*}
    \begin{exparts*}
      \partsitem \( a_{2,1} \)
      \partsitem \( a_{1,2} \)
      \partsitem \( a_{2,2} \)
      \partsitem \( a_{3,1} \)
    \end{exparts*}
    \begin{answer}
      \begin{exparts*}
        \partsitem \( 2 \)
        \partsitem \( 3 \)
        \partsitem \(-1 \)
        \partsitem Not defined.
      \end{exparts*}  
    \end{answer}
  \recommended \item 
    State the size of each matrix.
    \begin{exparts*}
      \partsitem \(
        \begin{mat}[r]
          1  &0  &4  \\
          2  &1  &5
        \end{mat}  \)
      \partsitem \(
        \begin{mat}[r]
          1  &1  \\
         -1  &1  \\
          3  &-1
        \end{mat}  \)
      \partsitem \(
        \begin{mat}[r]
          5  &10 \\
         10  &5
        \end{mat}  \)
    \end{exparts*}
    \begin{answer}
      \begin{exparts*}
        \partsitem \( \nbym{2}{3} \)
        \partsitem \( \nbym{3}{2} \)
        \partsitem \( \nbym{2}{2} \)
      \end{exparts*}  
    \end{answer}
  \recommended \item 
    Give the \( \nbyn{4} \) matrix whose
    \( i,j \)-th entry is
    \begin{exparts*}
      \partsitem \( i+j \);
      \partsitem \( -1 \) to the \( i+j \) power.
    \end{exparts*}
    \begin{answer}
      \begin{exparts*}
        \partsitem \(
           \begin{mat}[r]
             2  &3  &4  &5  \\
             3  &4  &5  &6  \\
             4  &5  &6  &7  \\
             5  &6  &7  &8
           \end{mat} \)
        \partsitem \(
           \begin{mat}[r]
             1  &-1  &1   &-1  \\
            -1  &1   &-1  &1  \\
             1  &-1  &1   &-1  \\
            -1  &1   &-1  &1
           \end{mat} \)
      \end{exparts*}  
    \end{answer}
  \recommended \item 
    Do the indicated vector operation, if it is defined.
    \begin{exparts*}
      \partsitem \( \colvec{2 \\ 1 \\ 1}
               +\colvec{3 \\ 0 \\ 4} \)
      \partsitem \( 5\colvec{4 \\ -1} \)
      \partsitem \( \colvec{1 \\ 5 \\ 1}
               -\colvec{3 \\ 1 \\ 1} \)
      \partsitem \( 7\colvec{2 \\ 1}
               +9\colvec{3 \\ 5} \)
      \partsitem \( \colvec{1 \\ 2}
               +\colvec{1 \\ 2 \\ 3} \)
      \partsitem \( 6\colvec{3 \\ 1 \\ 1}
               -4\colvec{2 \\ 0 \\ 3}
               +2\colvec{1 \\ 1 \\ 5} \)
    \end{exparts*}
    \begin{answer}
      \begin{exparts*}
        \partsitem \( \colvec{5 \\ 1 \\ 5} \)
        \partsitem \( \colvec{20 \\ -5} \)
        \partsitem \( \colvec{-2 \\ 4 \\ 0} \)
        \partsitem \( \colvec{41 \\ 52} \)
        \partsitem Not defined.
        \partsitem \( \colvec{12 \\ 8 \\ 4} \)
      \end{exparts*}  
     \end{answer}
  \recommended \item  \label{exer:SolveInMatrixNotation}
    Solve each system using matrix notation.
    Express the solution using vectors.
    \begin{exparts*}
      \partsitem \( \begin{linsys}[t]{2}
                  3x  &+  &6y  &=  &18  \\
                   x  &+  &2y  &=  &6   
                   \end{linsys}  \)
      \partsitem \( \begin{linsys}[t]{2}
                   x  &+  &y   &=  &1  \\
                   x  &-  &y   &=  &-1   
                    \end{linsys}  \)
      \partsitem \( \begin{linsys}[t]{3}
                   x_1  &   &     &+  &x_3   &=  &4  \\
                   x_1  &-  &x_2  &+  &2x_3  &=  &5  \\
                  4x_1  &-  &x_2  &+  &5x_3  &=  &17  
                   \end{linsys}  \)
      \partsitem \( \begin{linsys}[t]{3}
                   2a   &+  &b    &-  &c     &=  &2  \\
                   2a   &   &     &+  &c     &=  &3  \\
                    a   &-  &b    &   &      &=  &0   
                    \end{linsys}  \)
      \partsitem \( \begin{linsys}[t]{4}
                     x  &+  &2y   &-   &z   &    &    &=  &3  \\
                    2x  &+  &y    &    &    &+   &w   &=  &4  \\
                     x  &-  &y    &+   &z   &+   &w   &=  &1  
                    \end{linsys}  \)
      \partsitem \( \begin{linsys}[t]{4}
                     x  &   &     &+   &z   &+   &w   &=  &4  \\
                    2x  &+  &y    &    &    &-   &w   &=  &2  \\
                    3x  &+  &y    &+   &z   &    &    &=  &7  
                     \end{linsys}  \)
    \end{exparts*}
    \begin{answer}
      \begin{exparts}
        \partsitem This reduction
          \begin{equation*}
            \begin{amat}[r]{2}
              3  &6  &18 \\
              1  &2  &6
            \end{amat}
            \grstep{(-1/3)\rho_1+\rho_2}
            \begin{amat}[r]{2}
              3  &6  &18 \\
              0  &0  &0
            \end{amat}
          \end{equation*}
          leaves \( x \) leading and \( y \) free.
          Making \( y \) the parameter, gives \( x=6-2y \) and this solution
          set.
          \begin{equation*}
            \set{\colvec{6 \\ 0}+\colvec{-2 \\ 1}y
              \suchthat y\in\Re}
          \end{equation*}
        \partsitem A reduction
          \begin{equation*}
            \begin{amat}[r]{2}
              1  &1  &1  \\
              1  &-1 &-1
            \end{amat}
            \grstep{-\rho_1+\rho_2}
            \begin{amat}[r]{2}
              1  &1  &1  \\
              0  &-2 &-2
            \end{amat}
          \end{equation*}
          gives the unique solution \( y=1 \), \( x=0 \).
          The solution set is this.
          \begin{equation*}
            \set{\colvec{0 \\ 1} }
          \end{equation*}
        \partsitem Gauss's Method
          \begin{equation*}
            \begin{amat}[r]{3}
              1  &0  &1  &4  \\
              1  &-1 &2  &5  \\
              4  &-1 &5  &17
            \end{amat}
            \grstep[-4\rho_1+\rho_3]{-\rho_1+\rho_2}
            \begin{amat}[r]{3}
              1  &0  &1  &4  \\
              0  &-1 &1  &1  \\
              0  &-1 &1  &1
            \end{amat}     
            \grstep{-\rho_2+\rho_3}
            \begin{amat}[r]{3}
              1  &0  &1  &4  \\
              0  &-1 &1  &1  \\
              0  &0  &0  &0
            \end{amat}
          \end{equation*}
          leaves \( x_1 \) and \( x_2 \) leading with \( x_3 \) free.
          The solution set is this.
          \begin{equation*}
            \set{\colvec{4 \\ -1 \\ 0}+\colvec{-1 \\ 1 \\ 1}x_3
              \suchthat x_3\in\Re}
          \end{equation*}
        \partsitem This reduction
          \begin{align*}
            \begin{amat}[r]{3}
              2  &1  &-1 &2  \\
              2  &0  &1  &3  \\
              1  &-1 &0  &0
            \end{amat}
            &\grstep[-(1/2)\rho_1+\rho_3]{-\rho_1+\rho_2}
            \begin{amat}[r]{3}
              2  &1    &-1   &2  \\
              0  &-1   &2    &1  \\
              0  &-3/2 &1/2  &-1
            \end{amat}                                          \\
            &\grstep{(-3/2)\rho_2+\rho_3}
            \begin{amat}[r]{3}
              2  &1  &-1   &2  \\
              0  &-1 &2    &1  \\
              0  &0  &-5/2 &-5/2
            \end{amat}
          \end{align*}
          shows that the solution set is a singleton set.
          \begin{equation*}
            \set{\colvec{1 \\ 1 \\ 1}}
          \end{equation*}
        \partsitem This reduction is straightforward
          \begin{align*}
            \begin{amat}[r]{4}
              1  &2  &-1 &0  &3 \\
              2  &1  &0  &1  &4 \\
              1  &-1 &1  &1  &1
            \end{amat}
            &\grstep[-\rho_1+\rho_3]{-2\rho_1+\rho_2}
            \begin{amat}[r]{4}
              1  &2  &-1 &0  &3  \\
              0  &-3 &2  &1  &-2 \\
              0  &-3 &2  &1  &-2
            \end{amat}                                       \\
            &\grstep{-\rho_2+\rho_3}
            \begin{amat}[r]{4}
              1  &2  &-1 &0  &3  \\
              0  &-3 &2  &1  &-2 \\
              0  &0  &0  &0  &0
            \end{amat}
          \end{align*}
          and ends with \( x \) and $y$ leading while \( z \) and \( w \) are
          free.
          Solving for \( y \) gives \( y=(2+2z+w)/3 \) and substitution shows
          that \( x+2(2+2z+w)/3-z=3 \) so \( x=(5/3)-(1/3)z-(2/3)w \),
          making this the solution set.
          \begin{equation*}
            \set{\colvec{5/3 \\ 2/3 \\ 0 \\ 0}
                 +\colvec{-1/3 \\ 2/3 \\ 1 \\ 0}z
                 +\colvec{-2/3 \\ 1/3 \\ 0 \\ 1}w
                 \suchthat z,w\in\Re}
          \end{equation*}
        \partsitem The reduction
          \begin{align*}
            \begin{amat}[r]{4}
              1  &0  &1  &1  &4 \\
              2  &1  &0  &-1 &2 \\
              3  &1  &1  &0  &7
            \end{amat}
            &\grstep[-3\rho_1+\rho_3]{-2\rho_1+\rho_2}
            \begin{amat}[r]{4}
              1  &0  &1  &1  &4 \\
              0  &1  &-2 &-3 &-6\\
              0  &1  &-2 &-3 &-5
            \end{amat}                                       \\
            &\grstep{-\rho_2+\rho_3}
            \begin{amat}[r]{4}
              1  &0  &1  &1  &4 \\
              0  &1  &-2 &-3 &-6\\
              0  &0  &0  &0  &1
            \end{amat}
          \end{align*}
          shows that there is no solution\Dash the solution set is empty.
      \end{exparts} 
    \end{answer}
      % sage: P = matrix(QQ, [[1, 2, -1, 0], [2, 1, 0, 1], [1, -1, 1, 1]])
      % sage: v = vector(QQ, [3, 4, 1])
      % sage: P_prime = P.augment(v, subdivide=True)
      % sage: gauss_method(P_prime)
      % [ 1  2 -1  0| 3]
      % [ 2  1  0  1| 4]
      % [ 1 -1  1  1| 1]
      % take -2 times row 1 plus row 2
      % take -1 times row 1 plus row 3
      % [ 1  2 -1  0| 3]
      % [ 0 -3  2  1|-2]
      % [ 0 -3  2  1|-2]
      % take -1 times row 2 plus row 3
      % [ 1  2 -1  0| 3]
      % [ 0 -3  2  1|-2]
      % [ 0  0  0  0| 0]
      % sage: x,y,z,w=var('x,y,z,w')
      % sage: solve([x+2*y-z==3, -3*y+2*z+w==-2], x, y)
      % [[x == -2/3*w - 1/3*z + 5/3, y == 1/3*w + 2/3*z + 2/3]]
  \item \label{exer:SlvMatNot}
    Solve each system using matrix notation.
    Give each solution set in vector notation.
    \begin{exparts*}
      \partsitem \( \begin{linsys}[t]{3}
                  2x  &+  &y  &-  &z  &=  &1  \\
                  4x  &-  &y  &   &   &=  &3  
                \end{linsys}  \)
      \partsitem \( \begin{linsys}[t]{4}
                   x  &   &   &-  &z  &   &   &=  &1  \\
                      &   &y  &+  &2z &-  &w  &=  &3  \\
                   x  &+  &2y &+  &3z &-  &w  &=  &7  
               \end{linsys}  \)
      \partsitem \( \begin{linsys}[t]{4}
                   x  &-  &y  &+  &z  &   &   &=  &0  \\
                      &   &y  &   &   &+  &w  &=  &0  \\
                  3x  &-  &2y &+  &3z &+  &w  &=  &0  \\
                      &   &-y &   &   &-  &w  &=  &0  
               \end{linsys}  \)
      \partsitem \( \begin{linsys}[t]{5}
                   a  &+  &2b &+  &3c &+  &d  &-  &e  &=  &1  \\
                  3a  &-  &b  &+  &c  &+  &d  &+  &e  &=  &3  
               \end{linsys}  \)
    \end{exparts*}
    \begin{answer}
      \begin{exparts}
      \partsitem The reduction
        \begin{equation*}
          \begin{amat}[r]{3}
            2  &1  &-1  &1  \\
            4  &-1 &0   &3
          \end{amat}
          \grstep{-2\rho_1+\rho_2}
          \begin{amat}[r]{3}
            2  &1  &-1  &1  \\
            0  &-3 &2   &1
          \end{amat}
        \end{equation*}
        ends with \( x \) and \( y \) leading, and with \( z \) free.
        Solving for \( y \) gives \( y=(1-2z)/(-3) \), and then substitution
        \( 2x+(1-2z)/(-3)-z=1 \) shows that \( x=((4/3)+(1/3)z)/2 \).
        Hence the solution set is this.
        \begin{equation*}
          \set{\colvec{2/3 \\ -1/3 \\ 0}
               +\colvec{1/6 \\ 2/3 \\ 1}z
              \suchthat z\in\Re}
        \end{equation*}
      \partsitem This application of Gauss's Method
        \begin{align*}
          \begin{amat}[r]{4}
            1  &0  &-1  &0  &1 \\
            0  &1  &2   &-1 &3 \\
            1  &2  &3   &-1 &7
          \end{amat}
          &\grstep{-\rho_1+\rho_3}
          \begin{amat}[r]{4}
            1  &0  &-1  &0  &1 \\
            0  &1  &2   &-1 &3 \\
            0  &2  &4   &-1 &6
          \end{amat}                                   \\
          &\grstep{-2\rho_2+\rho_3}
          \begin{amat}[r]{4}
            1  &0  &-1  &0  &1 \\
            0  &1  &2   &-1 &3 \\
            0  &0  &0   &1  &0
          \end{amat}
        \end{align*}
        leaves  \( x \), \( y \), and \( w \)  leading.
        The solution set is here.
        \begin{equation*}
          \set{\colvec{1 \\ 3 \\ 0 \\ 0}
               +\colvec{1 \\ -2 \\ 1 \\ 0}z
              \suchthat z\in\Re}
        \end{equation*}
      \partsitem This row reduction
        \begin{align*}
          \begin{amat}[r]{4}
            1  &-1 &1   &0  &0 \\
            0  &1  &0   &1  &0 \\
            3  &-2 &3   &1  &0 \\
            0  &-1 &0   &-1 &0
          \end{amat}
          &\grstep{-3\rho_1+\rho_3}
          \begin{amat}[r]{4}
            1  &-1 &1   &0  &0 \\
            0  &1  &0   &1  &0 \\
            0  &1  &0   &1  &0 \\
            0  &-1 &0   &-1 &0
          \end{amat}                                       \\
          &\grstep[\rho_2+\rho_4]{-\rho_2+\rho_3}
          \begin{amat}[r]{4}
            1  &-1 &1   &0  &0 \\
            0  &1  &0   &1  &0 \\
            0  &0  &0   &0  &0 \\
            0  &0  &0   &0  &0
          \end{amat}
        \end{align*}
        ends with \( z \) and \( w \) free.
        We have this solution set.
        \begin{equation*}
          \set{\colvec{0 \\ 0 \\ 0 \\ 0}
               +\colvec{-1 \\ 0 \\ 1 \\ 0}z
               +\colvec{-1 \\ -1 \\ 0 \\ 1}w
              \suchthat z,w\in\Re}
        \end{equation*}
      \partsitem Gauss's Method
        \begin{equation*}
          \begin{amat}[r]{5}
            1  &2  &3   &1  &-1 &1  \\
            3  &-1 &1   &1  &1  &3
          \end{amat}
          \grstep{-3\rho_1+\rho_2}
          \begin{amat}[r]{5}
            1  &2  &3   &1  &-1 &1  \\
            0  &-7 &-8  &-2 &4  &0
          \end{amat}
        \end{equation*}
        ends with \( c \), \( d \), and \( e \) free.
        Solving for \( b \) shows that \( b=(8c+2d-4e)/(-7) \) and then 
        substitution
        \( a+2(8c+2d-4e)/(-7)+3c+1d-1e=1 \) shows that 
        \( a=1-(5/7)c-(3/7)d-(1/7)e \) and we have the solution set.
        \begin{equation*}
          \set{\colvec{1 \\ 0 \\ 0 \\ 0 \\ 0}
               +\colvec{-5/7 \\ -8/7 \\ 1 \\ 0 \\ 0}c
               +\colvec{-3/7 \\ -2/7 \\ 0 \\ 1 \\ 0}d
               +\colvec{-1/7 \\ 4/7 \\ 0 \\ 0 \\ 1}e
              \suchthat c,d,e\in\Re}
        \end{equation*}
    \end{exparts}  
   \end{answer}
  \item 
    Solve each system using matrix notation.
    Express the solution set using vectors.
    \begin{exparts*}
      \partsitem 
        $\begin{linsys}{3}
          3x &+ &2y &+ &z &= &1 \\
          x  &- &y  &+ &z &= &2 \\
          5x &+ &5y &+ &z &= &0
        \end{linsys}$
      \partsitem
        $\begin{linsys}{4}
          x  &+ &y  &- &2z &= &0 \\
          x  &- &y  &  &   &= &-3 \\
          3x &- &y  &- &2z &= &-6  \\
             &  &2y &- &2z &= &3  
        \end{linsys}$
      \partsitem
        $\begin{linsys}{5}
          2x  &- &y  &- &z &+ &w &= &4 \\
           x  &+ &y  &+ &z &  &  &= &-1 
        \end{linsys}$
      \partsitem
         $\begin{linsys}{3}
          x  &+ &y  &- &2z &= &0 \\
          x  &- &y  &  &   &= &-3 \\
          3x &- &y  &- &2z &= &0    
        \end{linsys}$ 
    \end{exparts*}
    \begin{answer}
      \begin{exparts}
        \partsitem This reduction
          \begin{align*}
            \begin{amat}{3}
              3 &2  &1  &1 \\
              1 &-1 &1  &2 \\
              5 &5  &1  &0
            \end{amat}
            &\grstep[-(5/3)\rho_1+\rho_3]{-(1/3)\rho_1+\rho_2}
            \begin{amat}{3}
              3 &2    &1     &1 \\
              0 &-5/3 &2/3   &5/3 \\
              0 &5/3  &-2/3  &-5/3
            \end{amat}                                     \\
            &\grstep{\rho_2+\rho_3}
            \begin{amat}{3}
              3 &2    &1     &1 \\
              0 &-5/3 &2/3   &5/3 \\
              0 &0    &0     &0
            \end{amat}
          \end{align*}
          gives this solution set.
          \begin{equation*}
            \set{\colvec{x \\ y \\ z}=\colvec{1 \\ -1 \\ 0}
                                       +\colvec{-3/5 \\ 2/5 \\ 1}z
                             \suchthat z\in\Re}
          \end{equation*}
        \partsitem
          This is the reduction.
          \begin{align*}
            \begin{amat}{3}
              1 &1   &-2 &0 \\
              1 &-1  &0  &3  \\
              3 &-1  &-2 &-6 \\
              0 &2   &-2 &3
            \end{amat}
            &\grstep[-3\rho_1+\rho_3]{-\rho_1+\rho_2}
            \begin{amat}{3}
              1 &1   &-2 &0 \\
              0 &-2  &2  &-3  \\
              0 &-4  &4  &-6 \\
              0 &2   &-2 &3
            \end{amat}                                  \\
            &\grstep[\rho_2+\rho_4]{-2\rho_2+\rho_3}
            \begin{amat}{3}
              1 &1   &-2 &0 \\
              0 &-2  &2  &-3  \\
              0 &0   &0  &0 \\
              0 &0   &0  &0
            \end{amat}
          \end{align*}
          The solution set is this.
          \begin{equation*}
            \set{\colvec{-3/2 \\ 3/2 \\ 0}
                 +\colvec{1 \\ 1 \\ 1}z
                 \suchthat z\in\Re}
          \end{equation*}
      \partsitem
         Gauss's Method
         \begin{equation*}
           \begin{amat}{4}
             2 &-1 &-1 &1 &4 \\
             1 &1  &1  &0 &-1
           \end{amat}
           \grstep{-(1/2)\rho_1+\rho_2}
           \begin{amat}{4}
             2 &-1   &-1   &1    &4 \\
             0 &3/2  &3/2  &-1/2 &-3
           \end{amat}
         \end{equation*}
         gives the solution set.
         \begin{equation*}
           \set{\colvec{1 \\ -2 \\ 0 \\ 0}
                  +\colvec{0 \\ -1 \\ 1 \\ 0}z
                   \colvec{-1/3 \\ 1/3 \\ 0 \\ 1}w
                 \suchthat z,w\in\Re}
         \end{equation*}
       \partsitem
          Here is the reduction.
          \begin{equation*}
            \begin{amat}{3}
              1 &1  &-2 &0  \\
              1 &-1 &0  &-3 \\
              3 &-1 &-2 &0
           \end{amat}
           \grstep[-3\rho_1+\rho_3]{-\rho_1+\rho_2}
           \begin{amat}{3}
             1 &1  &-2 &0  \\
             0 &-2 &2  &-3 \\
             0 &-4  &4  &0
          \end{amat}
          \grstep{-2\rho_2+\rho_3}
          \begin{amat}{3}
            1 &1  &-2 &0  \\
            0 &-2 &2  &-3 \\
            0 &0  &0  &6
          \end{amat}
        \end{equation*}
        The solution set is empty,~$\set{}$.
      \end{exparts}
    \end{answer}
  \recommended \item 
    The vector is in the set.
    What value of the parameters produces that vector?
    \begin{exparts}
      \partsitem $\colvec{5 \\ -5}$,
        $\set{\colvec{1 \\ -1}k\suchthat k\in\Re}$
      \partsitem $\colvec{-1 \\ 2 \\ 1}$,
        $\set{\colvec{-2 \\ 1 \\ 0}i
           +\colvec{3 \\ 0 \\ 1}j\suchthat i,j\in\Re}$
      \partsitem $\colvec{0 \\ -4 \\ 2}$,
        $\set{\colvec{1 \\ 1 \\ 0}m
               +\colvec{2 \\ 0 \\ 1}n\suchthat m,n\in\Re}$
    \end{exparts}
    \begin{answer}
      For each problem we get a system of linear equations by looking at the 
      equations of components.
      \begin{exparts}
       \partsitem By inspection $k=5$.
       \partsitem The second components show that $i=2$ and the third
       components show that $j=1$.
       \partsitem The second components give $m=-4$, the 
       third $n=2$.
      \end{exparts} 
    \end{answer}
  \item 
    Decide if the vector is in the set.
    \begin{exparts}
      \partsitem $\colvec{3 \\ -1}$,
        $\set{\colvec{-6 \\ 2}k\suchthat k\in\Re}$
      \partsitem $\colvec{5 \\ 4}$,
        $\set{\colvec{5 \\ -4}j\suchthat j\in\Re}$
      \partsitem $\colvec{2 \\ 1 \\ -1}$,
        $\set{\colvec{0 \\ 3 \\ -7}
             +\colvec{1 \\ -1 \\ 3}r\suchthat r\in\Re}$
      \partsitem $\colvec{1 \\ 0 \\ 1}$,
        $\set{\colvec{2 \\ 0 \\ 1}j
            +\colvec{-3 \\ -1 \\ 1}k\suchthat j,k\in\Re}$
    \end{exparts}
    \begin{answer}
      \begin{exparts}
        \partsitem Yes, the vector is in the set.
            By inspection, take $k=-1/2$.
        \partsitem No, the vector is not in the set.
            The two equations  $5=5j$ and
            $4=-4j$ have no common solution.
        \partsitem Yes, the vector is in the set.
            This system
            \begin{equation*}
              \begin{linsys}{2}
                0  &+   &r  &=  &2  \\
                3  &-   &r  &=  &1  \\
               -7  &+   &3r &=  &-1  \\
              \end{linsys}
            \end{equation*}
            has the solution $r=2$.
        \partsitem No.
            This system has no solution.
            \begin{equation*}
              \begin{linsys}{2}
                2j  &-  &3k  &=  1  \\
                0j  &-  &k   &=  0  \\
                 j  &+  &k   &=  1  \\
              \end{linsys}
            \end{equation*}
            The second components give $k=0$
            and then the third components give $j=1$.
            But the first components are not consistent with those. 
      \end{exparts}
     \end{answer}
  \item \cite{Cleary}
    A farmer with 1200 acres is considering planting three different crops, 
    corn, soybeans, and oats.   
    The farmer wants to use all~$1200$ acres.  
    Seed corn costs \$$20$ per acre, while soybean and oat seed cost 
    \$$50$ and~\$$12$ per acre respectively.  
    The farmer has \$$40\,000$ available to buy seed and intends to 
    spend it all.
    \begin{exparts}  
      \item Use the information above to formulate two linear equations 
        with three unknowns and solve it.
     \item Solutions to the system are choices that the farmer can make.  
        Write down two reasonable solutions.
     \item Suppose that in the fall when the crops mature, the farmer 
        can bring in revenue of \$$100$ per acre for corn, 
        \$$300$ per acre for soybeans and \$$80$ per acre for oats.  
        Which of your two solutions in the prior part would have resulted 
        in a larger revenue? 
    \end{exparts}
    \begin{answer}
      \begin{exparts}
        \item Let $c$ be the number of acres of corn, $s$ be the number of 
          acres of soy, and $a$ be the number of acres of oats.
          \begin{equation*}
            \begin{linsys}{3}
              c   &+   &s   &+   &a   &=   &1200 \\ 
            20c   &+   &50s &+   &12a &=   &40\,000  
            \end{linsys}
            \grstep{-20\rho_1+\rho_2}
            \begin{linsys}{3}
              c   &+   &s   &+   &a   &=   &1200 \\ 
                  &    &30s &-   &8a  &=   &16\,000  
            \end{linsys}
          \end{equation*}
          To describe the solution set we can parametrize using $a$.
          \begin{equation*}
            \set{\colvec{c \\ s \\ a}
                 =\colvec{2\,000/3 \\ 1\,600/3 \\ 0}
                  +\colvec{-19/15 \\ 4/15 \\ 1}a
                 \suchthat a\in\Re}
          \end{equation*}
        \item There are many possible answers. 
          We can take $a=0$, which leads to $c=2000/3\approx 666.67$ and
          $s=1600/3\approx 533.33$.
          We can also set $c=0$ and solve, giving  
          $a=10\,000/19\approx 526.32$
          and $s=12\,800/19\approx 673.68$.
          % sage: P = matrix(QQ, [[1, 1, 1], [20, 50, 12]])
          % sage: v = vector(QQ, [1200, 40000])
          % sage: P_prime = P.augment(v, subdivide=True)
          % sage: gauss_method(P_prime)
          % [    1     1     1| 1200]
          % [   20    50    12|40000]
          % take -20 times row 1 plus row 2
          % [    1     1     1| 1200]
          % [    0    30    -8|16000]
          % sage: c,s,a=var('c,s,a')
          % sage: solve([30*s-8*a==16000, c+s+a==1200], c,s)
          % [[c == -19/15*a + 2000/3, s == 4/15*a + 1600/3]]
          % sage: round(2000/3, ndigits=2)
          % 666.67
          % sage: round(1600/3, ndigits=2)
          % 533.33
          % sage: (-2000/3)*(-15/19)
          % 10000/19
          % sage: round(10000/19, ndigits=2)
          % 526.32
          % sage: (4/15)*(10000/19)+(1600/3)
          % 12800/19
          % sage: round(12800/19, ndigits=2)
          % 673.68
          % sage: c=0
          % sage: a=10000/19
          % sage: s=12800/19
          % sage: c+s+a
          % 1200
          % sage: 20*c+50*s+12*a
          % 40000
        \item Plugging in the prior part's first example answers into 
          $100c+300s+80a$ gives $680\,000/3$, approximately~$226\,666.67$.
          The second example answer gives 
          $100\cdot 0+300\cdot 12800/19s+80\cdot 10000/19=4\,640\,000/19$,
          about~$244\,210.53$.
          % sage: a=0
          % sage: c=2000/3
          % sage: s=1600/3
          % sage: 100*c+300*s+80*a
          % 680000/3
          % sage: round(680000/3, ndigits=2)
          % 226666.67
          % sage: c=0
          % sage: a=10000/19
          % sage: s=12800/19
          % sage: 100*c+300*s+80*a
          % 4640000/19
          % sage: round(4640000/19, ndigits=2)
          % 244210.53
      \end{exparts}
    \end{answer}
  \item 
    Parametrize the solution set of this one-equation system.
    \begin{equation*}
      x_1+x_2+\cdots+x_n=0
    \end{equation*}
    \begin{answer}
      This system has one equation.
      The leading variable is \( x_1 \), the other variables are free.
      \begin{equation*}
        \set{\colvec{-1 \\ 1 \\ \vdotswithin{-1} \\ 0}x_2
             +\cdots+
             \colvec{-1 \\ 0 \\ \vdotswithin{-1} \\ 1}x_n
             \suchthat x_2,\ldots,x_n\in\Re}
      \end{equation*}  
     \end{answer}
  \recommended \item 
    \begin{exparts}
    \partsitem Apply Gauss's Method to the left-hand side to solve
      \begin{equation*}
        \begin{linsys}{4}
          x  &+  &2y  &    &    &-   &w   &=   &a   \\
         2x  &   &    &+   &z   &    &    &=   &b   \\
          x  &+  &y   &    &    &+   &2w  &=   &c   
        \end{linsys}
      \end{equation*}
      for \( x \), \( y \), \( z \), and \(  w \), in terms of the 
      constants $a$, $b$, and $c$.
    \partsitem Use your answer from the prior part to solve this.
      \begin{equation*}
        \begin{linsys}{4}
          x  &+  &2y  &    &    &-   &w   &=   &3   \\
         2x  &   &    &+   &z   &    &    &=   &1   \\
          x  &+  &y   &    &    &+   &2w  &=   &-2
        \end{linsys}
      \end{equation*}
    \end{exparts}
    \begin{answer}
      \begin{exparts}
        \partsitem Gauss's Method here gives
          \begin{multline*}
            \begin{amat}{4}
              1  &2  &0  &-1  &a  \\
              2  &0  &1  &0   &b  \\
              1  &1  &0  &2   &c
            \end{amat}
            \grstep[-\rho_1+\rho_3]{-2\rho_1+\rho_2}
            \begin{amat}{4}
              1  &2  &0  &-1  &\multicolumn{1}{l}{a}  \\
              0  &-4 &1  &2   &\multicolumn{1}{l}{-2a+b}  \\
              0  &-1 &0  &3   &\multicolumn{1}{l}{-a+c}
            \end{amat}                                  \\
            \grstep{-(1/4)\rho_2+\rho_3}
            \begin{amat}{4}
              1  &2  &0    &-1  &\multicolumn{1}{l}{a}  \\
              0  &-4 &1    &2   &\multicolumn{1}{l}{-2a+b}  \\
              0  &0  &-1/4 &5/2 &\multicolumn{1}{l}{-(1/2)a-(1/4)b+c}
            \end{amat}
          \end{multline*}
          leaving \( w \) free.
          Solve: \(  z=2a+b-4c+10w \),
          and \( -4y=-2a+b-(2a+b-4c+10w)-2w \) so
          \( y=a-c+3w \), and
          \( x=a-2(a-c+3w)+w=-a+2c-5w. \)
          Therefore the solution set is this.
          \begin{equation*}
             \set{\colvec{-a+2c \\ a-c \\ 2a+b-4c \\ 0}
                  +\colvec{-5 \\ 3 \\ 10 \\ 1}w
                  \suchthat w\in\Re}
          \end{equation*}
        \partsitem Plug in with \( a=3 \), \( b=1 \), and \( c=-2 \).
          \begin{equation*}
             \set{\colvec{-7 \\ 5 \\ 15 \\ 0}
                  +\colvec{-5 \\ 3 \\ 10 \\ 1}w
                  \suchthat w\in\Re}
          \end{equation*}
      \end{exparts}  
     \end{answer}
  \item 
    Why is the comma needed in the notation `\( a_{i,j} \)'
    for matrix entries?
    \begin{answer}
       Leaving the comma out, say by writing \( a_{123} \),
       is ambiguous.
       It could mean $a_{1,23}$ or $a_{12,3}$.  
    \end{answer}
  \item  
    For any matrix \( A \), the 
    \definend{transpose}\index{transpose}
    \index{matrix!transpose}
    of \( A \), written
    \( \trans{A} \), is the matrix whose columns are the rows of \( A \).
    Find the transpose of each of these.
    \begin{exparts*}
      \partsitem \( \begin{mat}[r]
                  1  &2  &3  \\
                  4  &5  &6
               \end{mat}  \)
      \partsitem \( \begin{mat}[r]
                  2  &-3 \\
                  1  &1
               \end{mat}  \)
      \partsitem \( \begin{mat}[r]
                  5  &10 \\
                 10  &5
               \end{mat}  \)
      \partsitem \( \colvec{1 \\ 1 \\ 0} \)
    \end{exparts*}
    \begin{answer}
      \begin{exparts*}
        \partsitem \( \begin{mat}[r]
                   1  &4  \\
                   2  &5  \\
                   3  &6
                 \end{mat}  \)
        \partsitem \( \begin{mat}[r]
                   2  &1  \\
                  -3  &1
                 \end{mat}  \)
        \partsitem \( \begin{mat}[r]
                   5  &10 \\
                  10  &5
                 \end{mat}  \)
        \partsitem \( \rowvec{1 &1 &0}  \)
      \end{exparts*}  
     \end{answer}
  \recommended \item 
    \begin{exparts}
      \partsitem Describe all functions \( f(x)=ax^2+bx+c \) 
        such that \( f(1)=2 \) and \( f(-1)=6 \).
      \partsitem Describe all functions \( f(x)=ax^2+bx+c \) 
        such that \( f(1)=2 \).
    \end{exparts}
    \begin{answer}
      \begin{exparts}
        \partsitem Plugging in \( x=1 \) and \( x=-1 \) gives
          \begin{equation*}
            \begin{linsys}{3}
              a  &+  &b   &+  &c  &=  &2  \\
              a  &-  &b   &+  &c  &=  &6  
            \end{linsys}
            \grstep{-\rho_1+\rho_2}
            \begin{linsys}{3}
              a  &+  &b   &+  &c  &=  &2  \\
                 &   &-2b &   &   &=  &4  
              \end{linsys}
          \end{equation*}
          so the set of functions is
          \( \set{f(x)=(4-c)x^2-2x+c\suchthat c\in\Re} \).
        \partsitem Putting in \( x=1 \) gives
          \begin{equation*}
            \begin{linsys}{3}
              a  &+  &b   &+  &c  &=  &2  
            \end{linsys}
          \end{equation*}
          so the set of functions is
          \( \set{f(x)=(2-b-c)x^2+bx+c\suchthat b,c\in\Re} \).
      \end{exparts}  
    \end{answer}
  \item Show that any set of five points from the plane \( \Re^2 \) lie on a
    common conic section, that is, they all satisfy some equation of the
    form \( ax^2+by^2+cxy+dx+ey+f=0 \) where some of \( a,\,\ldots\,,f \)
    are nonzero.
    \begin{answer}
      Plugging in the five pairs of the form $(x,y)$ gives a system with
      five equations in the six unknowns $a$, \ldots~$f$.
      There are more unknowns than equations so unless there is some 
      inconsistency between the equations, at least one variable will be free.
      Thus if there is no inconsistency then the system has not just one 
      solution, it has infinitely many.

      No inconsistency can exist\Dash these five equations 
      must have at least one solution in 
      common\Dash because each of these five
      is solved by $a=0$, \ldots~$f=0$. 
    \end{answer}
  \item 
    Make up a four equations/four unknowns system having
    \begin{exparts}
      \partsitem a one-parameter solution set;
      \partsitem a two-parameter solution set;
      \partsitem a three-parameter solution set.
    \end{exparts}
    \begin{answer}
      \begin{exparts}
      \partsitem Here is one\Dash the fourth equation is redundant 
        but still OK.
        \begin{equation*}
          \begin{linsys}{4}
             x  &+  &y  &-  &z  &+  &w  &=  &0  \\
                &   &y  &-  &z  &   &   &=  &0  \\
                &   &   &   &2z &+  &2w &=  &0  \\
                &   &   &   &z  &+  &w  &=  &0
          \end{linsys}
        \end{equation*}
       We could have made the fourth equation redundant in a more subtle way,
       say by making it the sum of the first two equations, but in either case
       we have one redundant equation.
      \partsitem Here is one.
        \begin{equation*}
          \begin{linsys}{4}
             x  &+  &y  &-  &z  &+  &w  &=  &0  \\
                &   &   &   &   &   &w  &=  &0  \\
                &   &   &   &   &   &w  &=  &0  \\
                &   &   &   &   &   &w  &=  &0
          \end{linsys}
        \end{equation*}
      \partsitem This is one.
        \begin{equation*}
          \begin{linsys}{4}
             x  &+  &y  &-  &z  &+  &w  &=  &0  \\
             x  &+  &y  &-  &z  &+  &w  &=  &0  \\
             x  &+  &y  &-  &z  &+  &w  &=  &0  \\
             x  &+  &y  &-  &z  &+  &w  &=  &0 
          \end{linsys}
        \end{equation*}
    \end{exparts}  
   \end{answer}
  \puzzle \item 
    \cite{Shepelev}
    This  puzzle  is  from  a  Russian web site
    % \texttt{http://www.arbuz.uz/}  
    and  there are many solutions
    to it, but mine uses  linear  algebra  and  is  very
    naive.   There's   a   planet  inhabited  by  arbuzoids
   (watermeloners, to  translate  from  Russian). 
   Those creatures are found in three colors: red, green and blue.  
   There  are  $13$~red
   arbuzoids,  $15$~blue  ones, and $17$~green. 
   When
   two differently colored arbuzoids meet,  they
   both change to the third color.

   The question is, can it ever happen that all
   of them assume the same color?
    \begin{answer}
       \answerasgiven
       My solution was to define the numbers  of  arbuzoids
       as $3$-dimensional vectors, and express all possible
       elementary transitions as such vectors, too.
       \begin{center}
         \begin{tabular}{rr}
           R: &$13$  \\
           G: &$15$  \\
           B: &$17$
         \end{tabular}
         \qquad
         Operations:
         $\colvec{-1 \\ -1 \\ 2}$, 
         $\colvec{-1 \\ 2 \\ -1}$, 
         and 
         $\colvec{2 \\ -1 \\ -1}$
       \end{center}
       Now, it is enough to check whether the  solution  to
       one  of  the  following  systems of linear equations
       exists:
       \begin{equation*}
         \colvec{13 \\ 15 \\ 17}
         +x\colvec{-1 \\ -1  \\ 2}
         +y\colvec{-1 \\ 2 \\ -1}
         +z\colvec{2 \\ -1 \\ -1}
         =\colvec{0 \\ 0 \\ 45}
         \qquad
         \text{(or $\colvec{0 \\ 45 \\ 0}$ or $\colvec{45 \\ 0 \\ 0}$)}
       \end{equation*}
       Solving
       \begin{equation*}
         \begin{amat}[r]{3}
          -1  &-1 &2  &-13  \\
          -1  &2  &-1 &-15  \\
           2  &-1 &-1 &28
         \end{amat}
         \grstep[2\rho_1+\rho_3]{-\rho_1+\rho_2}
         \repeatedgrstep{\rho_2+\rho_3}
         \begin{amat}[r]{3}
          -1  &-1 &2  &-13  \\
           0  &3  &-3 &-2  \\
           0  &0  &0  &0
         \end{amat}
       \end{equation*}
       gives $y+2/3=z$ so if the number of transformations $z$ is an integer
       then $y$ is not.
       The other two systems give similar conclusions so there is no
       solution.
    \end{answer}
  % \puzzle \item 
  %   \cite{USSROlympiad174}
  %   \begin{exparts}
  %     \partsitem Solve the system of equations.
  %       \begin{equation*}
  %         \begin{linsys}{2}
  %           ax  &+  &y  &=  &a^2  \\
  %            x  &+  &ay &=  &1
  %        \end{linsys}
  %       \end{equation*}
  %       For what values of $a$ does the system fail to have solutions, and
  %       for what values of $a$ are there infinitely many solutions?
  %     \partsitem Answer the above question for the system.
  %       \begin{equation*}
  %         \begin{linsys}{2}
  %           ax  &+  &y  &=  &a^3  \\    
  %            x  &+  &ay &=  &1
  %         \end{linsys}
  %       \end{equation*}
  %   \end{exparts}
  %   \begin{answer}
  %      \answerasgiven
  %      \begin{exparts}
  %       \partsitem Formal solution of the system yields
  %         \begin{equation*}
  %           x=\frac{a^3-1}{a^2-1}  
  %           \qquad
  %           y=\frac{-a^2+a}{a^2-1}.
  %         \end{equation*}
  %         If $a+1\neq 0$ and $a-1\neq 0$, then the system has the single
  %         solution
  %         \begin{equation*}
  %           x=\frac{a^2+a+1}{a+1}
  %           \qquad
  %           y=\frac{-a}{a+1}.
  %         \end{equation*}
  %         If $a=-1$, or if $a=+1$, then the formulas are meaningless; in the
  %         first instance we arrive at the system
  %         \begin{equation*}
  %           \left\{ 
  %           \begin{linsys}{2}
  %             -x &+  &y  &=  &1 \\
  %              x &-  &y  &=  &1
  %           \end{linsys}\right.
  %         \end{equation*}
  %         which is a contradictory system.
  %         In the second instance we have
  %         \begin{equation*}
  %           \left\{
  %           \begin{linsys}{2}
  %              x &+  &y  &=  &1 \\
  %              x &+  &y  &=  &1
  %           \end{linsys}\right.
  %         \end{equation*}
  %         which has an infinite number of solutions (for example, for 
  %         $x$ arbitrary, $y=1-x$).
  %       \partsitem Solution of the system yields
  %         \begin{equation*}
  %           x=\frac{a^4-1}{a^2-1}
  %           \qquad
  %           y=\frac{-a^3+a}{a^2-1}.
  %         \end{equation*}
  %         Here, is $a^2-1\neq 0$, the system has the single solution
  %         $x=a^2+1$, $y=-a$.
  %         For $a=-1$ and $a=1$, we obtain the systems
  %         \begin{equation*}
  %           \left\{
  %           \begin{linsys}{2}
  %             -x &+  &y  &=  &-1 \\
  %              x &-  &y  &=  &1
  %           \end{linsys}\right.
  %           \qquad
  %           \left\{
  %           \begin{linsys}{2}
  %              x &+  &y  &=  &1 \\
  %              x &+  &y  &=  &1
  %           \end{linsys}\right.
  %        \end{equation*}
  %        both of which have an infinite number of solutions.
  %     \end{exparts}
  %   \end{answer}
  \puzzle \item 
    \cite{MathMag52p48}
    In air a gold-sur\-faced sphere weighs \( 7588 \)
    grams.
    It is known that it may contain one or more of the metals aluminum,
    copper, silver, or lead.
    When weighed successively under standard conditions in water, benzene,
    alcohol, and glycerin its respective weights are \( 6588 \), \( 6688 \),
    \( 6778 \), and \( 6328 \) grams.
    How much, if any, of the forenamed metals does it contain if the
    specific gravities of the designated substances are 
    % taken to be as follows?
    % \begin{center}
    %    \begin{tabular}{lrclr}
    %      Aluminum  &\( 2.7 \)
    %         &\makebox[3em]{\mbox{}\hfill\mbox{}} &Alcohol &$0.81$ \\
    %      Copper    &\( 8.9 \)  &     &Benzene   &\( 0.90 \) \\
    %      Gold      &\( 19.3 \) &     &Glycerin &\( 1.26 \) \\
    %      Lead      &\( 11.3 \) &     &Water     &\( 1.00 \) \\
    %      Silver    &\( 10.8 \)
    %    \end{tabular}
    % \end{center}
    aluminum~$2.7$,  copper~$8.9$, gold~$19.3$, lead~$11.3$,
    silver~$10.8$, alcohol~$0.81$, benzene~$0.90$, glycerin~$1.26$, 
    and water~$1.00$.
    \begin{answer}
      \answerasgiven
      Let \( u \), \( v \), \( x \), \( y \), \( z \) be the volumes in
      \( {\rm cm}^3 \) of Al, Cu, Pb, Ag, and Au, respectively, contained in
      the sphere, which we assume to be not hollow.
      Since the loss of weight in water (specific gravity \( 1.00 \)) is
      \( 1000 \) grams, the volume of the sphere is \( 1000\mbox{ cm}^3 \)\!.
      Then the data, some of which is superfluous, though consistent, leads to
      only \( 2 \) independent equations, one relating volumes and the
      other, weights.
      \begin{equation*}
        \begin{linsys}{5}
           u  &+  &v    &+  &x     &+  &y     &+  &z     &=  &1000  \\
        2.7u  &+  &8.9v &+  &11.3x &+  &10.5y &+  &19.3z &=  &7558
        \end{linsys}
      \end{equation*}
      Clearly the sphere must contain some aluminum to bring its mean specific
      gravity below the specific gravities of all the other metals.
      There is no unique result to this part of the problem, for the amounts
      of three metals may be chosen arbitrarily, provided that the choices
      will not result in negative amounts of any metal.

      If the ball contains only aluminum and gold, there are
      \( 294.5\mbox{ cm}^3 \) of gold and \( 705.5\mbox{ cm}^3 \) of aluminum.
      Another possibility is \( 124.7\mbox{ cm}^3 \) each of Cu, Au, Pb, and
      Ag and \( 501.2\mbox{ cm}^3  \) of Al.

      (\authorsNote{I get slightly different numbers.
       If the metals besides Al and Au are zero then I get that 
       Al is $u= 58710/83 \approx 707.3 $ and Au is $z= 24290/83 \approx 292.7 $.
       If Al is $u= 501.2 $ and and Ag is $y= 124.7 $ and Au is $z= 124.7 $
       then I get that Cu is $v= 137.3 $ and Pb is $x= 112.1 $.})   
    \end{answer}
\end{exercises}

























\subsection{\texorpdfstring{$\text{General}=\text{Particular}+\text{Homogeneous}$}{General=Particular+Homogeneous}}
In the prior subsection's examples the descriptions of solution sets
all
have a vector that is a particular solution 
of the system added to an unrestricted combination of some other vectors.
The solution set from 
\nearbyexample{ex:ManyParamsInfManySolsSystem} illustrates.
\begin{equation*}
  \set{
   \underbracket[.7pt]{
     \colvec{0 \\ 4 \\ 0 \\ 0 \\ 0}}_{\text{\shortstack{\rule{0pt}{2ex}particular \\
                                                    solution}}}+
   \underbracket[.7pt]{w\colvec{1 \\ -1 \\ 3 \\ 1 \\ 0}+
       u\colvec{1/2 \\ -1 \\ 1/2 \\ 0 \\ 1}}_{\text{\shortstack{\rule{0pt}{2ex}unrestricted\\
                                                                combination}}}
       \suchthat w,u\in\Re}
\end{equation*}
The combination is unrestricted in that 
$w$ and $u$ can be any real numbers;~there
is no condition like ``such that $2w-u=0\mspace{1mu}$'' to restrict 
valid pairs $w,u$.

% That example shows a solution set that is infinite.
% The other two kinds also fit the pattern.
% A one-element solution set 
% has an unrestricted combination part that is trivial. 
% Instead of being a combination of two vectors or
% of one vector, it is a combination of no vectors
% (by convention, the sum of an empty set of vectors
% is the zero vector).
% A solution set that is empty also fits the pattern because there is no 
% particular solution and thus there are no sums of matching the form
% `particular plus unrestricted combination'.

\begin{theorem} \label{th:GenEqPartPlusHomo}
%<*th:GenEqPartPlusHomo>
Any linear system's 
solution set has the form 
\begin{equation*}
   \set{\vec{p}+c_1\vec{\beta}_1+\,\cdots\,+c_k\vec{\beta}_k
     \suchthat c_1,\,\ldots\,,c_k\in\Re}
\end{equation*}
where \( \vec{p} \) is any particular solution  
and the $c_i$ are the 
free variables in an echelon form version of that system.
%</th:GenEqPartPlusHomo>
\end{theorem}

This result shows that we can always parametrize solution sets
with the free variables.
We shall prove it in two setps, \nearbylemma{le:HomoSltnSpanVecs} and
\nearbylemma{th:GenEqPartHomo}.
The first talks about unrestricted linear combinations and
the second adds the particular solutions.

Before those lemmas, first consider systems that have the zero vector
as a particular solution.
% so that we can shorten $\vec{p}+c_1\vec{\beta}_1+\dots+c_k\vec{\beta}_k$
% to $c_1\vec{\beta}_1+\dots+c_k\vec{\beta}_k$.

\begin{definition} \label{df:HomogeneousEquation}
%<*df:HomogeneousEquation>
A linear equation is \definend{homogeneous}\index{homogeneous equation}%
\index{linear equation!homogeneous}
if it can be put into a form where all terms are linear.
Thus, $a_1x_1+\,\cdots\,+a_{n-1}x_{n-1}=a_nx_n$ and
$a_1x_1+\,\cdots\,+a_{n-2}x_{n-2}=a_{n-1}x_{n-1}+a_nx_n$, etc.,\ are homogeneous. 
Equivalently, the equation   
can be put in the form $a_1x_1+a_2x_2+\,\cdots\,+a_nx_n=0$.
% A \definend{homogeneous system} is composed of homogeneous equations.
%</df:HomogeneousEquation>
\end{definition}

% The equation $x+2y=3z$ is homogeneous because all three terms are of 
% first degree in their variables.
% We like to rearrange so that terms with variables are on the left, 
% to $x+2y-3z=0$.
% Thus another way to define a homogeneous linear equation is as one  
% that can be put in the form $a_1x_1+a_2x_2+\,\cdots\,+a_nx_n=0$.
\begin{example}  \label{ex:FirstExHomoSys}
With any linear system such as the one below on the left,
we associate a system of homogeneous equations by setting the constants to
zeros, as on the right.
\begin{equation*}
  \begin{linsys}{2}
    3x  &+  &4y  &=  3  \\
    2x  &-  &y   &=  1  
  \end{linsys}
  \qquad
  \begin{linsys}{2}
    3x  &+  &4y  &=  0  \\
    2x  &-  &y   &=  0  
  \end{linsys}
\end{equation*}
Compare the reduction of the original system
\begin{equation*}
  \begin{linsys}{2}
    3x  &+  &4y  &=  3  \\
    2x  &-  &y   &=  1  
  \end{linsys}
  \grstep{-(2/3)\rho_1+\rho_2}
  \begin{linsys}{2}
    3x  &+  &4y        &=  3  \\
        &   &-(11/3)y   &=  -1  
   \end{linsys}
\end{equation*}
with the reduction of its associated system.
\begin{equation*}
  \begin{linsys}{2}
    3x  &+  &4y  &=  0  \\
    2x  &-  &y   &=  0  
  \end{linsys}
  \grstep{-(2/3)\rho_1+\rho_2}
  \begin{linsys}{2}
    3x  &+  &4y        &=  0  \\
        &   &-(11/3)y   &=  0
   \end{linsys}
\end{equation*}
The two reductions go the same way.
We can study reducing a linear system by instead studying
reducing the associated system of homogeneous equations.
\end{example}

Studying the associated system has an advantage over
studying the original system.
Systems of homogeneous equations must be consistent 
since there is always at least one solution, the zero vector.
% Nonhomogeneous systems can be inconsistent, that is, can fail to have
% any solution.

\begin{example} \label{ex:HomoZeroOnlySol}
Some homogeneous systems have the zero vector as their only solution.
\begin{equation*}
   \begin{linsys}{3}
     3x  &+  &2y  &+  &z  &=  &0  \\
     6x  &+  &4y  &   &   &=  &0  \\
         &   &y   &+  &z  &=  &0  
   \end{linsys}
   \grstep{-2\rho_1 +\rho_2}
   \begin{linsys}{3}
      3x  &+  &2y  &+  &z  &=  &0  \\
          &   &    &   &-2z&=  &0  \\
          &   &y   &+  &z  &=  &0  
    \end{linsys}
   \grstep{\rho_2 \leftrightarrow\rho_3}
   \begin{linsys}{3}
      3x  &+  &2y  &+  &z  &=  &0  \\
          &   &y   &+  &z  &=  &0  \\
          &   &    &   &-2z&=  &0
    \end{linsys}
\end{equation*}
\end{example}

\begin{example} \label{ex:SolnChemProb}
Some homogeneous systems have many solutions.
One is the Chemistry problem\index{Chemistry problem} 
from the first page of the first subsection.
\begin{multline*}
  \begin{linsys}{4}
              7x  &   &   &-  &7z  &   &   &=  &0  \\
              8x  &+  &y  &-  &5z  &-  &2w &=  &0  \\
                  &   &y  &-  &3z  &   &   &=  &0  \\
                  &   &3y &-  &6z  &-  &w  &=  &0  
  \end{linsys}
  \grstep{-(8/7)\rho_1+\rho_2}
  \begin{linsys}{4}
                7x &   &   &-  &7z  &   &   &=  &0  \\
                   &   &y  &+  &3z  &-  &2w &=  &0  \\
                   &   &y  &-  &3z  &   &   &=  &0  \\
                   &   &3y &-  &6z  &-  &w  &=  &0  
   \end{linsys}                                        \\
  \grstep[-3\rho_2+\rho_4]{-\rho_2+\rho_3}
  \begin{linsys}{4}
                7x &   &   &-  &7z  &   &   &=  &0  \\
                   &   &y  &+  &3z  &-  &2w &=  &0  \\
                   &   &   &   &-6z &+  &2w &=  &0  \\
                   &   &   &   &-15z&+  &5w &=  &0  
   \end{linsys}                                        
  \grstep{-(5/2)\rho_3+\rho_4}
  \begin{linsys}{4}
                7x &   &   &-  &7z  &   &   &=  &0  \\
                   &   &y  &+  &3z  &-  &2w &=  &0  \\
                   &   &   &   &-6z &+  &2w &=  &0  \\
                   &   &   &   &    &   &0  &=  &0  
   \end{linsys}
\end{multline*}
The solution set
\begin{equation*}
  \set{\colvec{1/3 \\ 1 \\ 1/3 \\ 1}w \suchthat w\in\Re}
\end{equation*}
contains infinitely many vectors, not just the zero vector.
% (if we take \( w \) to be a number of molecules then solutions
% make sense only when \( w \) is a nonnegative multiple of $3$).
\end{example}

\begin{lemma} \label{le:HomoSltnSpanVecs}
%<*le:HomoSltnSpanVecs>
For any linear system of homogeneous equations there exist
vectors $\vec{\beta}_1$, \ldots, $\vec{\beta}_k$ such that the 
solution set of the system is
\begin{equation*}
  \set{c_1\vec{\beta}_1+\cdots+c_k\vec{\beta}_k \suchthat c_1,\ldots,c_k\in\Re}
\end{equation*}
where the $c_i$ are the
free variables in an echelon form version of the system.
%</le:HomoSltnSpanVecs>
\end{lemma}

The proof verifies that from a system of homogeneous equations
in echelon form,
we can use back substitution to express 
all of the leading variables in terms of the free variables.
To explain why this establishes the lemma, before the 
proof we walk through an example.
We will also take the opportunity to 
illustrate the proof's notation.

%<*ex:HomoSystemGivesSpanningSet0>
We will use  this system of homogeneous equations in echelon form.
\begin{equation*}
  \begin{linsys}{5}
     x_1  &+  &x_2   &+  &2x_3  &+  &x_4  &+  &x_5  &=  &0  \\
          &   &x_2   &+  &x_3   &+  &x_4  &-  &x_5  &=  &0  \\
          &   &      &   &      &   &2x_4 &+  &x_5  &=  &0  
  \end{linsys}
\end{equation*}
Start with the bottom row, equation~$m=3$.
Solve for its leading variable, $x_{\ell_3}=x_4$, using its
leading term $a_{m,\ell_m}\cdot x_m=a_{3,4}x_4=2\cdot x_4$
(`$\ell$' stands for ``leading'' so $\ell_m$ is the column number of 
row~$m$'s leading variable), giving $x_4=-(1/2)x_5$. 
That is the leading variable in terms of the
free variable~$x_5$.
%</ex:HomoSystemGivesSpanningSet0>

Now back substitution enters its iteration phase.
%<*ex:HomoSystemGivesSpanningSet1>
Let $t$ denote how many rows the current one is above the bottom.
For this step $t=1$ and the current row is~$m-t=3-1=2$.
We come to this row having so far expressed the leading variable 
in terms of free variables for all lower-down rows, which here is just row~$3$.
Into row~$2$, substitute for the lower-down leading variable
$x_{\ell_3}=x_4$, giving 
$x_2+x_3+(-(1/2)x_5)-x_5=0$.
Solve for this row's leading variable $x_{\ell_{m-t}}=x_{\ell_2}=x_2$, 
yielding $x_2=-x_3+(3/2)x_5$.

Next in the iteration is the row that is $t=2$ up from the bottom, row $m-t=1$.
Again we come to the row knowing  
for lower rows 
how to express the leading variables $x_{\ell_3}=x_4$ and $x_{\ell_2}=x_2$ 
in terms of free variables.
Substitute, giving 
$x_1+(-x_3+(3/2)x_5)+2_3+(-(1/2)x_5)+x_5=0$.
Solve for the final leading variable
as $x_{\ell_1}=x_1=-x_3-2x_5$.
%</ex:HomoSystemGivesSpanningSet2>

%<*ex:HomoSystemGivesSpanningSet3>
We are done. 
Put the solution into vector notation
\begin{equation*}
  \set{
  \colvec{-1 \\ -1 \\ 1 \\ 0 \\ 0}x_3
   +\colvec{-2 \\ 3/2 \\ 0 \\ -1/2 \\ 1}x_5 
  \suchthat x_3,x_5\in\Re}
\end{equation*}
and recognize this as the combination of vectors 
associated with the free variables that is called for in
the lemma.
%</ex:HomoSystemGivesSpanningSet3>

One more point before the proof.
Back substitution moves up the system, using the work
from lower rows to do the next row.
This suggests the proof strategy of mathematical 
induction.\index{mathematical induction}\index{proof techniques!induction}\index{induction}%

Induction is an important and non-obvious technique that we will
use a number of times in this book.
It involves verifying a statement, using a base step and
an inductive step.
In the base step, we verify that the statement is true in a first 
instance, here that for an echelon form linear system we can write 
the bottom equation's leading variable in terms of free variables.
In the inductive step we verify an implication, that 
if the statement
is true for all prior cases then it follows 
that the statement holds also for the next case.
Here we will show that if for the
lower-down rows we can express the leading variables 
in terms of the free variables, then for the next row up (if there is one)
we can express the leading variable in those terms also.

Verifying these two proves that the statement holds for all of the rows. 
By the base step it is true for the bottom equation. 
By the inductive step the fact that it is true for the bottom
equation shows that 
it is true for the next one up. 
Then, invoking the inductive step again, 
the fact that it is true for the bottom two rows implies that it is 
true for the third equation up,
etc.\appendrefs{mathematical induction}%

\begin{proof}
%<*pf:HomoSltnSpanVecs0>
Here is the statement that we will prove using induction:~in an 
echelon form system of homogeneous equations we can express 
all leading variables in terms of the free variables.

% Ignore any  \( 0=0 \) equations. 
% (if the
% system consists only of \( 0=0 \)~equations then the lemma is trivially true
% because there are no leading variables).
% If there are any variables that appear only with with a coefficient of
% $0$ then they are free, and will never lead a row, so we will
% ignore that possibility.  
%</pf:HomoSltnSpanVecs0>

%<*pf:HomoSltnSpanVecs1>
For the base step, consider the bottom-most non-$0=0$ equation.
(The system has only homogeneous equations so it is consistent, 
and thus while it may have $0=0$ rows,  
it has no rows like $0=7$.)
\begin{equation*}
  a_{m,\ell_m}x_{\ell_m}+a_{m,\ell_{m}+1}x_{\ell_{m}+1}+\cdots+a_{m,n}x_n=0
  \qquad
  a_{m,\ell_m}\neq 0
   % \tag{$*$}
\end{equation*}
%</pf:HomoSltnSpanVecs1>
%<*pf:HomoSltnSpanVecs2>
Since this is the bottom row 
any variables % $x_{\ell_{m}+1}$, \ldots{} 
after the
leading one must be free. 
Move these to the right hand side and divide by $a_{m,\ell_m}$
\begin{equation*}
  x_{\ell_m}
  =(-a_{m,\ell_{m}+1}/a_{m,\ell_m})x_{\ell_{m}+1}+\cdots+(-a_{m,n}/a_{m,\ell_m})x_n
\end{equation*}
to express the leading variable in terms of free variables.
%</pf:HomoSltnSpanVecs2>
(A fine point:~if the bottom equation
has no variables to the right 
of~$x_{\ell_m}$ then \( x_{\ell_m}=0 \).
This satisfies the statement that we are verifying because
it has \( x_{\ell_m} \) written 
as a sum of a number of the free variables, namely as the sum of zero many, 
and the convention is that 
a trivial sum totals to~$0$.)

%<*pf:HomoSltnSpanVecs3>
For the inductive step, assume 
that for the \( m \)-th equation,
and the \text{\( (m-1) \)-th}
equation, etc., up to and including the \mbox{\( m-(t-1) \)-th}~equation,
we can
express the leading variable of that equation in terms of free variables. 
% That is,
% assume that the statement that we can
% express the leading variable of the row in terms of the free variables
% holds for the bottom-most $t$~rows, with $1\leq t<m$.
We will check that 
therefore for the next equation up, 
equation \( m-t \), we can also express the leading variable in terms
of those that are free.

Into that equation, for each leading variable from a lower-down row,
\( x_{\ell_m} \), \ldots{} \( x_{\ell_{m-(t-1)}} \), 
substitute its expression in terms of free variables.
The result has a leading term  
$a_{m-t,\ell_{m-t}}x_{\ell_{m-t}}$ with \( a_{m-t,\ell_{m-t}}\neq 0 \),
and the rest of the left hand side
is a combination of free variables.
Move the free variables to
the right side and divide by
\( a_{m-t,\ell_{m-t}} \). 
That leaves
\( x_{\ell_{m-t}} \) expressed in terms of free variables.

We have done both the base step and the inductive step, 
so the statement is proved.
%</pf:HomoSltnSpanVecs3>
\end{proof}

\begin{definition}
The set of vector combinations
$\set{c_1\vec{\beta}_1+\cdots+c_k\vec{\beta}_k \suchthat c_1,\ldots,c_k\in\Re}$
is \definend{generated by}\index{generated by}\index{generated} 
or~\definend{spanned by}\index{spanned by}\index{span} 
the set 
\( \set{\smash{\vec{\beta}_1},\ldots,\smash{\vec{\beta}_k}} \).
\end{definition}

% The proof mentions a tricky point.
% We follow the convention that the sum of an empty set of vectors is the 
% zero vector.
% In particular, we needs this where 
% a homogeneous system has a unique solution because it
% takes the \( c \)'s to be the free variables
% and if there is a unique solution then there are no free variables.

The next result finishes the proof of  \nearbytheorem{th:GenEqPartPlusHomo}.

\begin{lemma}  \label{th:GenEqPartHomo}
%<*th:GenEqPartHomo>
For any linear system and for any particular solution $\vec{p}$ of that 
system,
the solution set equals
$
  \set{\vec{p}+\vec{h} \suchthat \text{ \( \vec{h} \) satisfies the
                                associated homogeneous system}     }$.
%</th:GenEqPartHomo>
\end{lemma}

% So fixing any particular solution gives the above 
% description of the solution set.

\begin{proof}
Fix a particular solution~$\vec{p}$.
Let $S$ be this set.
\begin{equation*}
  \set{\vec{p}+\vec{h} \suchthat \text{ \( \vec{h} \) satisfies the
                                associated homogeneous system} }
\end{equation*}
We will establish mutual set inclusion, that any solution to the system is 
a member of~$S$, and that any member of $S$ is a solution to the 
system.\appendrefs{set equality}

%<*pf:GenEqPartHomo0>
For set inclusion the first way,
assume that \( \vec{v} \) solves the system.
Consider the difference \( \vec{v}-\vec{p} \).
Plug it into any equation.
\begin{multline*}
  a_{i,1}(v_1-p_1)+\cdots+a_{i,n}(v_n-p_n) \\
  =(a_{i,1}v_1+\cdots+a_{i,n}v_n)       
  -(a_{i,1}p_1+\cdots+a_{i,n}p_n)  
  =d_i-d_i                 
  =0
\end{multline*}
Thus $\vec{v}-\vec{p}=\vec{h}$, where~$\vec{h}$ solves the 
associated system of homogeneous equations.
Rearranging as $\vec{v}=\vec{p}+\vec{h}$ has $\vec{v}$
decomposed into the form for~$S$, as the sum of the fixed
particular solution and a vector that solves the associated system.
%</pf:GenEqPartHomo0>

%<*pf:GenEqPartHomo1>
For set inclusion the other way, take a member of~$S$, $\vec{p}+\vec{h}$,
where \( \vec{p} \) is a particular solution and \( \vec{h} \) solves the
associated system of homogeneous equations. 
Plugging $\vec{p}+\vec{h}$ into any equation 
\begin{multline*}
  a_{i,1}(p_1+h_1)+\cdots+a_{i,n}(p_n+h_n)  \\
  =(a_{i,1}p_1+\cdots+a_{i,n}p_n)      
   +(a_{i,1}h_1+\cdots+a_{i,n}h_n)  
  =d_i+0                                
  =d_i
\end{multline*}
shows that it is a solution of the given system. 
%</pf:GenEqPartHomo1>
\end{proof}

Solution sets of systems of linear equations fall into three types, 
many solutions, unique solution, and no solution.
We next illlustrate how each type fits the 
`$\text{General}=\text{Particular}+\text{Homogeneous}$' pattern.

\begin{example}
Gauss's Method on this system
\begin{equation*}
  \begin{linsys}{2}
    x  &+  &4y  &=  &-2 \\
   3x  &+  &12y &=  &-6 
  \end{linsys}
  \grstep{-3\rho_1+\rho_2}
  \begin{linsys}{2}
    x  &+  &4y  &=  &-2 \\
       &   &0   &=  &0 
  \end{linsys}
  \tag{$*$}
\end{equation*}
shows that we can write the leading variable in terms of the free variable
as $x=-2-4y$. 
The general solution set is this.
\begin{equation*}
  \set{\colvec{x \\ y} 
        =\colvec{-2 \\ 0}+\colvec{-4 \\ 1}y
       \suchthat y\in\R}
\end{equation*}
Reduction of the system of homogeneous equations
proceeds in the same way.
\begin{equation*}
  \begin{linsys}{2}
    x  &+  &4y  &=  &0 \\
   3x  &+  &12y &=  &0 
  \end{linsys}
  \grstep{-3\rho_1+\rho_2}
  \begin{linsys}{2}
    x  &+  &4y  &=  &0 \\
       &   &0   &=  &0 
  \end{linsys}
\end{equation*}
Its general solution is just ($*$)'s solution without the particular
solution.
\begin{equation*}
  \set{% 
        \colvec{-4 \\ 1}y
       \suchthat y\in\R}
\end{equation*}
\end{example}

\begin{example} \label{ex:IllusGenEqPartHomo}
This system has a unique solution.
\begin{equation*}
  \begin{linsys}{3}
    x  &+  &2y  &-  &z  &=  &1  \\
    2x &+  &4y  &   &   &=  &2  \\
       &   &y   &-  &3z &=  &0
  \end{linsys}
  \grstep{-2\rho_1+\rho_2}
  \begin{linsys}{3}
    x  &+  &2y  &-  &z  &=  &1  \\
       &   &    &   &2z &=  &0  \\
       &   &y   &-  &3z &=  &0
  \end{linsys}                           
  \grstep{\rho_2\leftrightarrow\rho_3} 
  \begin{linsys}{3}
      x  &+  &2y  &-  &z  &=  &1  \\      
         &   &y   &-  &3z &=  &0  \\
         &   &    &   &2z &=  &0
   \end{linsys}
\end{equation*}
Its general solution is a singleton set, comprised just of the
particular solution.
\begin{equation*}
  \set{\colvec{1 \\ 0 \\ 0} }
  \tag{$*$}
\end{equation*}
The associated system of homogeneous equations 
reduces via the same row operations 
\begin{equation*}
  \begin{linsys}{3}
    x  &+  &2y  &-  &z  &=  &0  \\
    2x &+  &4y  &   &   &=  &0  \\
       &   &y   &-  &3z &=  &0
  \end{linsys}
  \grstep{-2\rho_1+\rho_2}
  \repeatedgrstep{\rho_2\swap\rho_3} 
  \begin{linsys}{3}
      x  &+  &2y  &-  &z  &=  &0  \\      
         &   &y   &-  &3z &=  &0  \\
         &   &    &   &2z &=  &0
   \end{linsys}
\end{equation*}
Its solution set is also a singleton
(which we can think of as the collection of unrestricted combinations
of the empty set of vectors). 
\begin{equation*}
  \set{\colvec{0 \\ 0 \\ 0} }
\end{equation*}
Adding the particular solution to all the members of this set
gives back the singleton set ($*$).
\end{example}

\begin{example}
When a system has no solution, that is, 
when the general solution set is empty,
then the system also fits the pattern.
\begin{equation*}
  \begin{linsys}{4}
    x  &   &  &+  &z  &+ &w  &=  &-1  \\
   2x  &-  &y &   &   &+ &w  &=  &3   \\
    x  &+  &y &+  &3z &+ &2w &=  &1   
  \end{linsys}
  \grstep[-\rho_1+\rho_3]{-2\rho_1+\rho_2}
  \begin{linsys}{4}
    x  &   &  &+  &z  &+ &w  &=  &-1  \\
       &   &-y&-  &2z &- &w  &=  &5   \\
       &   &y &+  &2z &+ &w  &=  &2   
   \end{linsys}
\end{equation*}
The final two equations
conflict so it has no particular solution. 
But the associated system of homogeneous equations does have a solution, 
as do all such systems, because the zero vector solves any homogeneous equation.
\begin{equation*}
  \begin{linsys}{4}
    x  &   &  &+  &z  &+ &w  &=  &0   \\
   2x  &-  &y &   &   &+ &w  &=  &0   \\
    x  &+  &y &+  &3z &+ &2w &=  &0   
  \end{linsys}
  \grstep[-\rho_1+\rho_3]{-2\rho_1+\rho_2}
  \grstep{\rho_2+\rho_3}
  \begin{linsys}{4}
    x  &   &  &+  &z  &+ &w  &=  &0   \\
       &   &-y&-  &2z &- &w  &=  &0   \\
       &   &  &   &   &  &0  &=  &0         
  \end{linsys}
\end{equation*}
In fact, its solution set is infinite. 
\begin{equation*}
  \set{\colvec{-1 \\ -2 \\ 1 \\ 0}z+\colvec{-1 \\ -1 \\ 0 \\ 1}w
         \suchthat z,w\in\Re}
\end{equation*}
Nonetheless, because the original system has no particular solution, its
general solution set is empty\Dash there are no vectors of the form
$\vec{p}+\vec{h}$ because there are no $\vec{p}\:$'s.
\end{example}

\begin{corollary} \label{co:ThreeKindsSolutionSets}
%<*co:ThreeKindsSolutionSets>
Solution sets of linear systems are either empty, have one element, or
have infinitely many elements.
%</co:ThreeKindsSolutionSets>
\end{corollary}

\begin{proof}
%<*pf:ThreeKindsSolutionSets0>
We've seen examples of all three happening so we need only prove
that there are no other possibilities.

First observe that a system of homogeneous equations with
at least one non-\( \zero \) solution $\vec{v}$ has infinitely many
solutions.
This is because any scalar multiple of~$\vec{v}$ also solves that
system, and there are infinitely many vectors in the set of scalar 
multiples of $\vec{v}$: if $s,t\in\Re$ are unequal then $s\vec{v}\neq t\vec{v}$,
since $s\vec{v}-t\vec{v}=(s-t)\vec{v}$ is
non-$\zero$ as  any non-$0$ component of $\vec{v}$, when 
rescaled by the non-$0$ factor $s-t$, will give a non-$0$ value.
%</pf:ThreeKindsSolutionSets0>

%<*pf:ThreeKindsSolutionSets1>
Now apply \nearbylemma{th:GenEqPartHomo} to conclude that a solution set
\begin{equation*}
  \set{\vec{p}+\vec{h}\suchthat
    \text{\( \vec{h} \) solves the associated homogeneous system}}
\end{equation*}
is either empty (if there is no particular solution \( \vec{p} \)),
or has one element (if there is a \( \vec{p} \) and the homogeneous system
has the unique solution \( \zero \)), or is infinite (if there is a
\( \vec{p} \) and the homogeneous system has a non-$\zero$ solution,
and thus by the prior paragraph has infinitely many solutions).
%</pf:ThreeKindsSolutionSets1>
\end{proof}

This table summarizes the factors affecting the size of a
general solution.

\smallskip
%<*table:KindsSolutionSets>
\begin{center} \small 
\begin{tabular}{r@{}c}
  &\hspace*{2.5em}\begin{tabular}{c} 
      \textit{number of solutions of the} \\[-.5ex]
      \textit{homogeneous system}
    \end{tabular}                            \\[1.7ex] 
  \begin{tabular}{r@{\hspace*{.5em}}} 
     \ \\[.6ex]
     \textit{particular} \\[-.55ex]
     \textit{solution}   \\[-.5ex]
     \textit{exists?}   
  \end{tabular}
  &\begin{tabular}{r|c@{\hspace*{1em}}c} % \cline{2-3}   
     \multicolumn{1}{c}{\ }
         &\textit{one}    &\textit{infinitely many}                    \\ 
     \cline{2-3}
     \textit{yes}    
        &\rule{0ex}{16pt}\begin{tabular}{@{}c@{}} unique \\[-.5ex] solution \end{tabular}
        &\begin{tabular}{@{}c@{}} infinitely many \\[-.5ex] solutions \end{tabular}
         \\[2ex] % \hline  
     \textit{no}    
        &\begin{tabular}{@{}c@{}} no \\[-.5ex] solutions \end{tabular}
        &\begin{tabular}{@{}c@{}} no \\[-.5ex] solutions \end{tabular} 
         \\ % \hline
   \end{tabular}
\end{tabular}
\end{center}
%</table:KindsSolutionSets>
\smallskip

The dimension on the top of the table is the simpler one.
When we perform Gauss's Method on a linear system, ignoring the
constants on the right side and so paying attention only
to the coefficients on the left-hand side,
we either end with every variable leading some row or else 
we find some variable that does not lead a row, that is,
we find some variable that is free. 
(We formalize ``ignoring the constants on the right'' by
considering the associated homogeneous system.)

A notable special case is that of
systems having the same number of equations as unknowns. 
Such a system will have a solution, and that solution will be unique, 
if and only if it
reduces to an echelon form system where every variable leads its row.
This will happen if and only if
the associated system of  homogeneous equations has a unique solution.

\begin{definition} \label{df:Nonsingular}
%<*df:Nonsingular>
A matrix is \definend{square}\index{square matrix}\index{matrix!square}
if it has the same number of rows and columns.
A square matrix is \definend{nonsingular}\index{nonsingular!matrix}
\index{matrix!nonsingular}
if it is the matrix of coefficients of a
homogeneous system with a unique solution.
It is
\definend{singular}\index{singular!matrix}\index{matrix!singular}
otherwise, that is,
if it is the matrix of coefficients of a homogeneous system with 
infinitely many solutions.
%</df:Nonsingular>
\end{definition}

\begin{example}
The first of these matrices is nonsingular while the second is singular
\begin{equation*}
  \begin{mat}[r]
    1  &2  \\
    3  &4
  \end{mat}
  \qquad
  \begin{mat}[r]
    1  &2  \\
    3  &6
  \end{mat}
\end{equation*}
because the first of these homogeneous systems has a unique solution 
while the second has infinitely many solutions.
\begin{equation*}
  \begin{linsys}[b]{2}
    x &+  &2y  &=  &0  \\
   3x &+  &4y  &=  &0  
  \end{linsys}
  \qquad
  \begin{linsys}[b]{2}
    x &+  &2y  &=  &0  \\
   3x &+  &6y  &=  &0
  \end{linsys}
\end{equation*}  
A system
with the same number of equations as variables
behaves in one of two ways.
If the matrix of coefficients is nonsingular then the system 
has a unique solution no matter what constants are on the right 
side:~for instance, Gauss's Method shows that this system
\begin{equation*}
  \begin{linsys}{2}
    x  &+  &2y  &=  &a \\
    3x &+  &4y  &=  &b
  \end{linsys}
\end{equation*}
has the unique solution $x=b-2a$ and  $y=(3a-b)/2$.
% sage: a, b, x, y = var('a, b, x, y')
% sage: sys = [x+2*y==a, 3*x+4*y==b]
% sage: solve(sys, x, y)
% [[x == -2*a + b, y == 3/2*a - 1/2*b]]
On the other hand, if the matrix of coefficients is
singular then the system never has a unique solution:~it 
either has no solutions or else infinitely many, as with these.
\begin{equation*}
  \begin{linsys}[b]{2}
    x  &+  &2y  &=   &1   \\
   3x  &+  &6y  &=   &2   
  \end{linsys}
  \qquad
  \begin{linsys}[b]{2}
    x  &+  &2y  &=   &1   \\
   3x  &+  &6y  &=   &3
  \end{linsys}
\end{equation*} 
\end{example}

\begin{remark}
One meaning of `singular'
is ``departing from general expectation.''
People often naively expect that a system 
with the same number of variables as equations has a unique solution.
% So we can think of the word as connoting 
% ``troublesome,'' or at least ``not ideal.''
(That `singular' applies to the systems that never have exactly one solution 
is ironic, but that is the standard term.)  
\end{remark}

% \begin{example}
% The systems from \nearbyexample{ex:FirstExHomoSys},
% \nearbyexample{ex:HomoZeroOnlySol},
% and \nearbyexample{ex:IllusGenEqPartHomo}
% each have an associated homogeneous system with a unique solution.
% Thus these matrices are nonsingular.
% \begin{equation*}
%   \begin{mat}[r]
%     3  &4  \\
%     2  &-1
%   \end{mat}
%   \qquad
%   \begin{mat}[r]
%     3  &2   &1  \\
%     6  &-4  &0  \\
%     0  &1   &1
%   \end{mat}
%   \qquad
%   \begin{mat}[r]
%     1  &2  &-1 \\
%     2  &4  &0  \\
%     0  &1  &-3
%   \end{mat}
% \end{equation*}
% The Chemistry problem from \nearbyexample{ex:SolnChemProb} 
% is a homogeneous system with more than one solution so its matrix
% is singular. 
% \begin{equation*}
%   \begin{mat}[r]
%     7  &0  &-7 &0  \\
%     8  &1  &-5 &-2 \\
%     0  &1  &-3 &0  \\
%     0  &3  &-6 &-1
%   \end{mat}
% \end{equation*}
% \end{example}

The table above has two dimensions.
We have considered the one on top:~we can tell
into which column a given linear system goes
solely by considering the system's left-hand side; the 
constants on the right-hand side play no role in this.

The table's other dimension, 
determining whether a particular solution exists, is tougher.
Consider these two systems with the same left side but different right sides.
\begin{equation*}
  \begin{linsys}[b]{2}
    3x &+ &2y &= &5  \\
    3x &+ &2y &= &5
  \end{linsys}
  \qquad
  \begin{linsys}[b]{2}
    3x &+ &2y &= &5  \\
    3x &+ &2y &= &4
  \end{linsys}
\end{equation*}
The first has a solution while the second does not, so
here the constants on the right side decide if the system has a solution.
We could conjecture that the left side of a linear system determines
the number of solutions while the right side determines whether solutions
exist but that guess is not correct.
Compare these two,
with the same right sides but different left sides.
\begin{equation*}
  \begin{linsys}[b]{2}
    3x &+ &2y &= &5  \\
    4x &+ &2y &= &4
  \end{linsys}
  \qquad
  \begin{linsys}[b]{2}
    3x &+ &2y &= &5  \\
    3x &+ &2y &= &4
  \end{linsys}
\end{equation*}
The first has a solution but the second does not.
Thus the constants on the right side
don't alone determine whether a solution exists.
Rather, that depends on some interaction between the left and
right.

For some intuition about that interaction,
consider this system with one of the coefficients left unspecified, as  
the variable~$c$.
\begin{equation*}
  \begin{linsys}{3}
    x  &+  &2y  &+  &3z  &=  &1  \\
    x  &+  &y   &+  &z   &=  &1  \\
   cx  &+  &3y  &+  &4z  &=  &0
  \end{linsys}
\end{equation*}
If \( c=2 \) then this system has no solution because the left-hand side 
has the third row as the sum of the first two, while the right-hand does not.
If \( c\neq 2 \) then this system has a unique solution (try it with \( c=1 \)).
For a system to have a solution, if one row of the matrix of coefficients on
the left is a linear combination of other rows
then on the right the constant from that row must be the same
combination of constants from the same rows.

More intuition about the interaction comes from studying linear
combinations.
That will be our focus in the second chapter, after we finish the study
of Gauss's Method itself in the rest of this chapter.

\begin{exercises}
  \item Solve this system.
    Then solve the associated homogeneous system.
    \begin{equation*}
      \begin{linsys}{4}
        x  &+ &y  &- &2z &= &0 \\
        x  &- &y  &  &   &= &-3 \\
        3x &- &y  &- &2z &= &-6  \\
           &  &2y &- &2z &= &3  
      \end{linsys}
   \end{equation*}
   \begin{answer}
     This reduction solves the system.
     \begin{align*}
       \begin{amat}{3}
         1 &1   &-2 &0 \\
         1 &-1  &0  &3  \\
         3 &-1  &-2 &-6 \\
         0 &2   &-2 &3
       \end{amat}
       &\grstep[-3\rho_1+\rho_3]{-\rho_1+\rho_2}
       \begin{amat}{3}
         1 &1   &-2 &0 \\
         0 &-2  &2  &-3  \\
         0 &-4  &4  &-6 \\
         0 &2   &-2 &3
       \end{amat}                                 \\
       &\grstep[\rho_2+\rho_4]{-2\rho_2+\rho_3}
       \begin{amat}{3}
         1 &1   &-2 &0 \\
         0 &-2  &2  &-3  \\
         0 &0   &0  &0 \\
         0 &0   &0  &0
       \end{amat}
     \end{align*}
     The solution set is this.
     \begin{equation*}
       \set{\colvec{-3/2 \\ 3/2 \\ 0}
            +\colvec{1 \\ 1 \\ 1}z
            \suchthat z\in\Re}
     \end{equation*}
     Similarly we can reduce the associated homogeneous system
     \begin{equation*}
       \begin{amat}{3}
         1 &1   &-2 &0 \\
         1 &-1  &0  &0  \\
         3 &-1  &-2 &0 \\
         0 &2   &-2 &0
       \end{amat}
       \grstep[-3\rho_1+\rho_3]{-\rho_1+\rho_2}
       % \begin{amat}{3}
       %   1 &1   &-2 &0 \\
       %   0 &-2  &2  &0  \\
       %   0 &-4  &4  &0 \\
       %   0 &2   &-2 &0
       % \end{amat}                                         \\
       \grstep[\rho_2+\rho_4]{-2\rho_2+\rho_3}
       \begin{amat}{3}
         1 &1   &-2 &0 \\
         0 &-2  &2  &0  \\
         0 &0   &0  &0 \\
         0 &0   &0  &0
       \end{amat}
     \end{equation*}
     to get its solution set.
     \begin{equation*}
       \set{
            \colvec{1 \\ 1 \\ 1}z
            \suchthat z\in\Re}
     \end{equation*}
   \end{answer}
  \recommended \item 
    Solve each system.
    Express the solution set using vectors.
    Identify a particular solution and the solution set of the
    homogeneous system.
    % (These systems also appear in 
    % Exercise~2.\nearbyexercise{exer:SolveInMatrixNotation}.)
    \begin{exparts*}
      \partsitem \( \begin{linsys}[t]{2}
                  3x  &+  &6y  &=  &18  \\
                   x  &+  &2y  &=  &6   
                   \end{linsys}  \)
      \partsitem \( \begin{linsys}[t]{2}
                   x  &+  &y   &=  &1  \\
                   x  &-  &y   &=  &-1   
                    \end{linsys}  \)
      \partsitem \( \begin{linsys}[t]{3}
                   x_1  &   &     &+  &x_3   &=  &4  \\
                   x_1  &-  &x_2  &+  &2x_3  &=  &5  \\
                  4x_1  &-  &x_2  &+  &5x_3  &=  &17  
                    \end{linsys}  \)
      \partsitem \( \begin{linsys}[t]{3}
                   2a   &+  &b    &-  &c     &=  &2  \\
                   2a   &   &     &+  &c     &=  &3  \\
                    a   &-  &b    &   &      &=  &0   
                    \end{linsys}  \)
      \partsitem \( \begin{linsys}[t]{4}
                     x  &+  &2y   &-   &z   &    &    &=  &3  \\
                    2x  &+  &y    &    &    &+   &w   &=  &4  \\
                     x  &-  &y    &+   &z   &+   &w   &=  &1  
                    \end{linsys}  \)
      \partsitem \( \begin{linsys}[t]{4}
                     x  &   &     &+   &z   &+   &w   &=  &4  \\
                    2x  &+  &y    &    &    &-   &w   &=  &2  \\
                    3x  &+  &y    &+   &z   &    &    &=  &7  
                    \end{linsys}  \)
    \end{exparts*}
    \begin{answer}
      \begin{exparts}
        \partsitem  Here is the Gauss-Jordan reduction.
          \begin{equation*}
            \begin{amat}{2}
              3  &6  &18  \\ 
              1  &2  &6  \\ 
            \end{amat}
            \grstep{(-1/3)\rho_{1}+\rho_{2}}
            \begin{amat}{2}
              3  &6  &18  \\ 
              0  &0  &0  \\ 
            \end{amat}
            \grstep{(1/3)\rho_{1}}
            \begin{amat}{2}
              1  &2  &6  \\ 
              0  &0  &0  \\ 
            \end{amat}
          \end{equation*}
          This is the solution set.
          \begin{equation*}
            S=\set{\colvec{6 \\ 0}+\colvec{-2 \\ 1}y
              \suchthat y\in\Re}
          \end{equation*}
          Here is a particular solution, and the solution set 
          to the associated homogeneous system.
          \begin{equation*}
            \colvec{6 \\ 0}
              \quad\text{and}\quad
            \set{\colvec{-2 \\ 1}y
              \suchthat y\in\Re}
          \end{equation*}
          (There is more than one
          correct answer for a particular solution.
          For instance, in the description given above of $S$, setting 
          the parameter~$y$ to be~$2$ gives a vector with first component~$2$ 
          and second component also~$2$.)
        \partsitem
          The reduction
          \begin{multline*}
            \begin{amat}{2}
              1  &1  &1  \\ 
              1  &-1  &-1  \\ 
            \end{amat}
            \grstep{-\rho_{1}+\rho_{2}}
            \begin{amat}{2}
              1  &1  &1  \\ 
              0  &-2  &-2  \\ 
            \end{amat}               \\
            \grstep{(-1/2)\rho_{2}}
            \begin{amat}{2}
              1  &1  &1  \\ 
              0  &1  &1  \\ 
            \end{amat}
            \grstep{-\rho_{2}+\rho_{1}}
            \begin{amat}{2}
              1  &0  &0  \\ 
              0  &1  &1  \\ 
            \end{amat}
          \end{multline*}
          gives that this is the solution set.
          \begin{equation*}
            \set{\colvec{0 \\ 1} }
          \end{equation*}
          Here is a particular solution, 
          and the solution set for the associated
          homogeneous system.
          \begin{equation*}
            \colvec{0 \\ 1}
              \qquad
            \set{\colvec{0 \\ 0} }
          \end{equation*}
        \partsitem
          This reduction
          \begin{align*}
            \begin{amat}{3}
              1  &0  &1  &4  \\ 
              1  &-1  &2  &5  \\ 
              4  &-1  &5  &17  \\ 
            \end{amat}
            &\grstep[-4\rho_{1}+\rho_{3}]{-\rho_{1}+\rho_{2}}
            \begin{amat}{3}
              1  &0  &1  &4  \\ 
              0  &-1  &1  &1  \\ 
              0  &-1  &1  &1  \\ 
            \end{amat}                 \\
            &\grstep{-\rho_{2}+\rho_{3}}
            \begin{amat}{3}
              1  &0  &1  &4  \\ 
              0  &-1  &1  &1  \\ 
              0  &0  &0  &0  \\ 
            \end{amat}
          \end{align*}
          shows that the solution set is infinite.
          \begin{equation*}
            \set{\colvec{4 \\ -1 \\ 0}+\colvec{-1 \\ 1 \\ 1}x_3
              \suchthat x_3\in\Re}
          \end{equation*}
          Here is a particular solution and the solution set for the associated
          homogeneous system.
          \begin{equation*}
            \colvec{4 \\ -1 \\ 0}
              \qquad
            \set{\colvec{-1 \\ 1 \\ 1}x_3
              \suchthat x_3\in\Re}
          \end{equation*}
        \partsitem
          Gauss's Method is routine.
          \begin{equation*}
            % \begin{amat}{3}
            %   2  &1  &-1  &2  \\ 
            %   2  &0  &1  &3  \\ 
            %   1  &-1  &0  &0  \\ 
            % \end{amat}
            \grstep[(-1/2)\rho_{1}+\rho_{3}]{-\rho_{1}+\rho_{2}}  
            % \begin{amat}{3}
            %   2  &1  &-1  &2  \\ 
            %   0  &-1  &2  &1  \\ 
            %   0  &-3/2  &1/2  &-1  \\ 
            % \end{amat}                            \\
            \repeatedgrstep{(-3/2)\rho_{2}+\rho_{3}}
            % \begin{amat}{3}
            %   2  &1  &-1  &2  \\ 
            %   0  &-1  &2  &1  \\ 
            %   0  &0  &-5/2  &-5/2  \\ 
            % \end{amat}
            \repeatedgrstep[-\rho_{2} \\ (-2/5)\rho_{3}]{(1/2)\rho_{1}}
            \begin{amat}{3}
              1  &1/2  &-1/2  &1  \\ 
              0  &1  &-2  &-1  \\ 
              0  &0  &1  &1  \\ 
            \end{amat}
          \end{equation*}
          The solution set is a singleton.
          \begin{equation*}
            \set{\colvec{1 \\ 1 \\ 1}}
          \end{equation*}
          This is a particular solution and the solution set for the associated
          system of homogeneous equations.
          \begin{equation*}
            \colvec{1 \\ 1 \\ 1}
              \qquad
            \set{\colvec{0 \\ 0 \\ 0}}
          \end{equation*}
        \partsitem
          Gauss's Method
          \begin{equation*}
            % \begin{amat}{4}
            %   1  &2  &-1  &0  &3  \\ 
            %   2  &1  &0  &1  &4  \\ 
            %   1  &-1  &1  &1  &1  \\ 
            % \end{amat}
            \grstep[-\rho_{1}+\rho_{3}]{-2\rho_{1}+\rho_{2}}
            % \begin{amat}{4}
            %   1  &2  &-1  &0  &3  \\ 
            %   0  &-3  &2  &1  &-2  \\ 
            %   0  &-3  &2  &1  &-2  \\ 
            % \end{amat}
            \repeatedgrstep{-\rho_{2}+\rho_{3}}
            \begin{amat}{4}
              1  &2  &-1  &0  &3  \\ 
              0  &-3  &2  &1  &-2  \\ 
              0  &0  &0  &0  &0  \\ 
            \end{amat}
          \end{equation*}
          gives that the solution set is infinite.
          \begin{equation*}
            \set{\colvec{5/3 \\ 2/3 \\ 0 \\ 0}
                 +\colvec{-1/3 \\ 2/3 \\ 1 \\ 0}z
                 +\colvec{-2/3 \\ 1/3 \\ 0 \\ 1}w
                 \suchthat z,w\in\Re}
          \end{equation*}
          A particular solution and the solution set for the associated
          homogeneous system are here.
          \begin{equation*}
            \colvec{5/3 \\ 2/3 \\ 0 \\ 0}
              \qquad
            \set{\colvec{-1/3 \\ 2/3 \\ 1 \\ 0}z
                 +\colvec{-2/3 \\ 1/3 \\ 0 \\ 1}w
                 \suchthat z,w\in\Re}
          \end{equation*}
        \partsitem 
          This reduction shows that 
          the system's solution set is empty.
          \begin{equation*}
            % \begin{amat}{4}
            %   1  &0  &1  &1  &4  \\ 
            %   2  &1  &0  &-1  &2  \\ 
            %   3  &1  &1  &0  &7  \\ 
            % \end{amat}
            \grstep[-3\rho_{1}+\rho_{3}]{-2\rho_{1}+\rho_{2}}
            % \begin{amat}{4}
            %   1  &0  &1  &1  &4  \\ 
            %   0  &1  &-2  &-3  &-6  \\ 
            %   0  &1  &-2  &-3  &-5  \\ 
            % \end{amat}
            \repeatedgrstep{-\rho_{2}+\rho_{3}}
            \begin{amat}{4}
              1  &0  &1  &1  &4  \\ 
              0  &1  &-2  &-3  &-6  \\ 
              0  &0  &0  &0  &1  \\ 
            \end{amat}
          \end{equation*}
          Thus there is no particular solution.
          In contrast, a system of homogeneous equations always has a solution.
          Here is the solution set for the system associated with this exercise's
          starting system.
          \begin{equation*}
            \set{\colvec{-1 \\ 2 \\ 1 \\ 0}z
                 +\colvec{-1 \\ 3 \\ 0 \\ 1}w
                 \suchthat z,w\in\Re}
          \end{equation*}
      \end{exparts}  
    \end{answer}
  \item 
    Solve each system, giving
    the solution set in vector notation.
    Identify a particular solution and the solution of the
    homogeneous system.
    \begin{exparts*}
      \partsitem \( \begin{linsys}[t]{3}
                  2x  &+  &y  &-  &z  &=  &1  \\
                  4x  &-  &y  &   &   &=  &3  
                  \end{linsys}  \)
      \partsitem \( \begin{linsys}[t]{4}
                   x  &   &   &-  &z  &   &   &=  &1  \\
                      &   &y  &+  &2z &-  &w  &=  &3  \\
                   x  &+  &2y &+  &3z &-  &w  &=  &7  
                   \end{linsys}  \)
      \partsitem \( \begin{linsys}[t]{4}
                   x  &-  &y  &+  &z  &   &   &=  &0  \\
                      &   &y  &   &   &+  &w  &=  &0  \\
                  3x  &-  &2y &+  &3z &+  &w  &=  &0  \\
                      &   &-y &   &   &-  &w  &=  &0  
                  \end{linsys}  \)
      \partsitem \( \begin{linsys}[t]{5}
                   a  &+  &2b &+  &3c &+  &d  &-  &e  &=  &1  \\
                  3a  &-  &b  &+  &c  &+  &d  &+  &e  &=  &3  
                  \end{linsys}  \)
    \end{exparts*}
    \begin{answer}
    (These exercises also appeared in the prior section, so 
    it may save effort to reuse your work from there.)
    \begin{exparts}
      \partsitem
        Here is the reduction.
        \begin{equation*}
          \begin{amat}{3}
            2  &1  &-1  &1  \\ 
            4  &-1  &0  &3  \\ 
          \end{amat}
          \grstep{-2\rho_{1}+\rho_{2}}
          \begin{amat}{3}
            2  &1  &-1  &1  \\ 
            0  &-3  &2  &1  \\ 
          \end{amat}
        \end{equation*}
        Back substitution gives this solution set.
        \begin{equation*}
          \set{\colvec{2/3 \\ -1/3 \\ 0}
               +\colvec{1/6 \\ 2/3 \\ 1}z
              \suchthat z\in\Re}
        \end{equation*}
        Here are a particular solution and the solution set for the associated
        system of homogeneous equations.
        \begin{equation*}
          \colvec{2/3 \\ -1/3 \\ 0}
            \qquad
        \set{\colvec{1/6 \\ 2/3 \\ 1}z
            \suchthat z\in\Re}
        \end{equation*}
      \partsitem
        The reduction goes as usual.
        \begin{equation*}
          % \begin{amat}{4}
          %   1  &0  &-1  &0  &1  \\ 
          %   0  &1  &2  &-1  &3  \\ 
          %   1  &2  &3  &-1  &7  \\ 
          % \end{amat}
          \grstep{-\rho_{1}+\rho_{3}}
          % \begin{amat}{4}
          %   1  &0  &-1  &0  &1  \\ 
          %   0  &1  &2  &-1  &3  \\ 
          %   0  &2  &4  &-1  &6  \\ 
          % \end{amat}
          \repeatedgrstep{-2\rho_{2}+\rho_{3}}
          \begin{amat}{4}
            1  &0  &-1  &0  &1  \\ 
            0  &1  &2  &-1  &3  \\ 
            0  &0  &0  &1  &0  \\ 
          \end{amat}
        \end{equation*}
        The solution set is infinite.
        \begin{equation*}
          \set{\colvec{1 \\ 3 \\ 0 \\ 0}
               +\colvec{1 \\-2 \\ 1 \\ 0}z
              \suchthat z\in\Re}
        \end{equation*}
        Here are 
        a particular solution and the solution set for the associated
        system of homogeneous equations.
        \begin{equation*}
          \colvec{1 \\ 3 \\ 0 \\ 0}
            \qquad
          \set{\colvec{1 \\ -2 \\ 1 \\ 0}z
              \suchthat z\in\Re}
        \end{equation*}
      \partsitem
        Gauss's Method
        \begin{equation*}
          % \begin{amat}{4}
          %   1  &-1  &1  &0  &0  \\ 
          %   0  &1  &0  &1  &0  \\ 
          %   3  &-2  &3  &1  &0  \\ 
          %   0  &-1  &0  &-1  &0  \\ 
          % \end{amat}
          \grstep{-3\rho_{1}+\rho_{3}}
          % \begin{amat}{4}
          %   1  &-1  &1  &0  &0  \\ 
          %   0  &1  &0  &1  &0  \\ 
          %   0  &1  &0  &1  &0  \\ 
          %   0  &-1  &0  &-1  &0  \\ 
          % \end{amat}
          \repeatedgrstep[\rho_{2}+\rho_{4}]{-\rho_{2}+\rho_{3}}
          \begin{amat}{4}
            1  &-1  &1  &0  &0  \\ 
            0  &1  &0  &1  &0  \\ 
            0  &0  &0  &0  &0  \\ 
            0  &0  &0  &0  &0  \\ 
          \end{amat}
        \end{equation*}
        gives this solution set.
        \begin{equation*}
          \set{\colvec{0 \\ 0 \\ 0 \\ 0}
               +\colvec{-1 \\ 0 \\ 1 \\ 0}z
               +\colvec{-1 \\ -1 \\ 0 \\ 1}w
              \suchthat z,w\in\Re}
        \end{equation*}
        We can decompose it into
        a particular solution and the solution set for the associated
        system of homogeneous equations.
        \begin{equation*}
          \colvec{0 \\ 0 \\ 0 \\ 0}
            \qquad
          \set{\colvec{-1 \\ 0 \\ 1 \\ 0}z
               +\colvec{-1 \\ -1 \\ 0 \\ 1}w
              \suchthat z,w\in\Re}
        \end{equation*}
      \partsitem
        The linear reduction is easy.
        \begin{equation*}
          % \begin{amat}{5}
          %   1  &2  &3  &1  &-1  &1  \\ 
          %   3  &-1  &1  &1  &1  &3  \\ 
          % \end{amat}
          \grstep{-3\rho_{1}+\rho_{2}}
          \begin{amat}{5}
            1  &2  &3  &1  &-1  &1  \\ 
            0  &-7  &-8  &-2  &4  &0  \\ 
          \end{amat}
        \end{equation*}
        The general solution set
        \begin{equation*}
          \set{\colvec{1 \\ 0 \\ 0 \\ 0 \\ 0}
               +\colvec{-5/7 \\ -8/7 \\ 1 \\ 0 \\ 0}c
               +\colvec{-3/7 \\ -2/7 \\ 0 \\ 1 \\ 0}d
               +\colvec{-1/7 \\ 4/7 \\ 0 \\ 0 \\ 1}e
              \suchthat c,d,e\in\Re}
        \end{equation*}
        breaks into 
        this particular solution and
        solution set for the associated system.
        \begin{equation*}
          \colvec{1 \\ 0 \\ 0 \\ 0 \\ 0}
            \qquad
          \set{\colvec{-5/7 \\ -8/7 \\ 1 \\ 0 \\ 0}c
               +\colvec{-3/7 \\ -2/7 \\ 0 \\ 1 \\ 0}d
               +\colvec{-1/7 \\ 4/7 \\ 0 \\ 0 \\ 1}e
              \suchthat c,d,e\in\Re}
        \end{equation*}
    \end{exparts}  
   \end{answer}
  \recommended \item 
    For the system
    \begin{equation*}
      \begin{linsys}{4}
       2x  &-  &y  &   &    &-  &w  &=  &3  \\
           &   &y  &+  &z   &+  &2w &=  &2  \\
        x  &-  &2y &-  &z   &   &   &=  &-1
      \end{linsys}
    \end{equation*}
    which of these can be used as the particular solution part of some
    general solution?
    \begin{exparts*}
      \partsitem   \( \colvec{0 \\ -3 \\ 5 \\ 0} \)
      \partsitem   \( \colvec{2 \\ 1 \\ 1 \\ 0} \)
      \partsitem   \( \colvec{-1 \\ -4 \\ 8 \\ -1} \)
    \end{exparts*}
    \begin{answer}
      See if each satisfies all three equations.
      \begin{exparts}
        \partsitem No, it does not the third equation.
        \partsitem Yes, it satisfies all three.
        \partsitem Yes.
      \end{exparts}  
    \end{answer}
  \recommended \item  
    \nearbylemma{th:GenEqPartHomo} says that we can use any particular solution 
    for $\vec{p}$.
    Find, if possible, a general solution to this system
    \begin{equation*}
      \begin{linsys}{4}
        x  &-  &y  &   &    &+  &w  &=  &4  \\
       2x  &+  &3y &-  &z   &   &   &=  &0  \\
           &   &y  &+  &z   &+  &w  &=  &4  
      \end{linsys}
    \end{equation*}
    that uses the given vector as its particular solution.
    \begin{exparts*}
      \partsitem   \( \colvec{0 \\ 0 \\ 0 \\ 4} \)
      \partsitem   \( \colvec{-5 \\ 1 \\ -7 \\ 10} \)
      \partsitem   \( \colvec{2 \\ -1 \\ 1 \\ 1} \)
    \end{exparts*}
    \begin{answer}
      Gauss's Method on the associated homogeneous system 
      \begin{align*}
        \begin{amat}[r]{4}
           1  &-1  &0  &1  &0  \\
           2  &3   &-1 &0  &0  \\
           0  &1   &1  &1  &0
        \end{amat}
        &\grstep{-2\rho_1+\rho_2}
        \begin{amat}[r]{4}
           1  &-1  &0  &1  &0  \\
           0  &5   &-1 &-2 &0  \\
           0  &1   &1  &1  &0
        \end{amat}                                \\
        &\grstep{-(1/5)\rho_2+\rho_3}
        \begin{amat}[r]{4}
           1  &-1  &0  &1  &0  \\
           0  &5   &-1 &-2 &0  \\
           0  &0   &6/5&7/5&0
        \end{amat}
      \end{align*}
      gives this is the solution to the homogeneous problem.
      \begin{equation*}
        \set{\colvec{-5/6 \\ 1/6 \\ -7/6 \\ 1}w\suchthat w\in\Re}
      \end{equation*}
      \begin{exparts}
        \partsitem This vector is indeed a particular solution,
          since plugging it in shows that it satisfies all three equations. 
          This is the desired general solution.
          \begin{equation*}
            \set{\colvec{0 \\ 0 \\ 0 \\ 4}+
                 \colvec{-5/6 \\ 1/6 \\ -7/6 \\ 1}w\suchthat w\in\Re}
          \end{equation*}
        \partsitem This vector is also a particular solution, so the
          general solution is this.
          \begin{equation*}
            \set{\colvec{-5 \\ 1 \\ -7 \\ 10}+
                 \colvec{-5/6 \\ 1/6 \\ -7/6 \\ 1}w\suchthat w\in\Re}
          \end{equation*}
        \partsitem That vector is not a solution of the system since
          it does not satisfy the third equation.
          No such general solution exists.
      \end{exparts} 
    \end{answer}
  \item 
     One is nonsingular while the other is singular.
     Which is which?
     \begin{exparts*}
       \partsitem $\begin{mat}[r]
           1  &3   \\
           4  &-12    
         \end{mat}$
       \partsitem $\begin{mat}[r]
           1  &3  \\
           4  &12  
         \end{mat}$
     \end{exparts*}
     \begin{answer}
       Both matrices are square.
       For each, the answer is `nonsingular' if after Gauss's Method 
       the echelon form result only non-$0$ numbers on the diagonal.
       \begin{exparts}
         \partsitem Gauss's Method leads to an echelon form
         \begin{equation*}
           \begin{mat}
             1  &3  \\ 
             4  &-12  \\ 
           \end{mat}
           \grstep{-4\rho_{1}+\rho_{2}}
           \begin{mat}
             1  &3  \\ 
             0  &-24  \\ 
           \end{mat}
         \end{equation*}
         whose diagonal entries are all nonzero.
         So the starting matrix is nonsingular.
         \partsitem
         The zeros on the diagonal at the end
         \begin{equation*}
           \begin{mat}
             1  &3  \\ 
             4  &12  \\ 
           \end{mat}
           \grstep{-4\rho_{1}+\rho_{2}}
           \begin{mat}
             1  &3  \\ 
             0  &0  \\ 
           \end{mat}
         \end{equation*}
         show that the starting matrix is nonsingular.
       \end{exparts}
     \end{answer}
  \recommended \item 
    Singular or nonsingular?
    \begin{exparts*}
      \partsitem \(
        \begin{mat}[r]
          1  &2  \\
          1  &3
        \end{mat}   \)
      \partsitem \(
        \begin{mat}[r]
          1  &2  \\
         -3  &-6
        \end{mat}   \)
      \partsitem \(
        \begin{mat}[r]
          1  &2  &1  \\
          1  &3  &1
        \end{mat}   \)
      \partsitem \(
        \begin{mat}[r]
          1  &2  &1  \\
          1  &1  &3  \\
          3  &4  &7
        \end{mat}   \)
      \partsitem \(
        \begin{mat}[r]
          2  &2  &1  \\
          1  &0  &5  \\
         -1  &1  &4
        \end{mat}   \)
    \end{exparts*}
    \begin{answer}
      \begin{exparts}
      \partsitem Nonsingular, because
        \begin{equation*}
          \grstep{-\rho_1+\rho_2}
          \begin{mat}[r]
            1  &2  \\
            0  &1
          \end{mat}
        \end{equation*}
        in echelon form  each row containing a leading entry,
        a nonzero entry on the diagonal.
      \partsitem Singular, because
        \begin{equation*}
          \grstep{3\rho_1+\rho_2}
          \begin{mat}[r]
            1  &2  \\
            0  &0
          \end{mat}
        \end{equation*}
        in echelon form the second row is without a leading entry.
      \partsitem Neither.
        A matrix must be square for either word to apply.
      \partsitem Singular.
        The echelon form has a row with no leading entry.
      \begin{equation*}
        % \begin{mat}
        %   1  &2  &1  \\ 
        %   1  &1  &3  \\ 
        %   3  &4  &7  \\ 
        % \end{mat}
        \grstep[-3\rho_{1}+\rho_{3}]{-\rho_{1}+\rho_{2}}
        % \begin{mat}
        %   1  &2  &1  \\ 
        %   0  &-1  &2  \\ 
        %   0  &-2  &4  \\ 
        % \end{mat}
        \repeatedgrstep{-2\rho_{2}+\rho_{3}}
        \begin{mat}
          1  &2  &1  \\ 
          0  &-1  &2  \\ 
          0  &0  &0  \\ 
        \end{mat}        
      \end{equation*}
      \partsitem Nonsingular.
      \begin{equation*}
        % \begin{mat}
        %   2  &2  &1  \\ 
        %   1  &0  &5  \\ 
        %   -1  &1  &4  \\ 
        % \end{mat}
        \grstep[(1/2)\rho_{1}+\rho_{3}]{(-1/2)\rho_{1}+\rho_{2}}
        % \begin{mat}
        %   2  &2  &1  \\ 
        %   0  &-1  &9/2  \\ 
        %   0  &2  &9/2  \\ 
        % \end{mat}
        \repeatedgrstep{2\rho_{2}+\rho_{3}}
        \begin{mat}
          2  &2  &1  \\ 
          0  &-1  &9/2  \\ 
          0  &0  &27/2  \\ 
        \end{mat}
      \end{equation*}
      Every entry on the diagonal of the echelon form is nonzero.
     \end{exparts}  
    \end{answer}
  \recommended \item 
    Is the vector in the set generated by the
    given set?
      \begin{exparts}
        \partsitem \( \colvec{2 \\ 3}, \)\
          \( \set{\colvec{1 \\ 4},
                \colvec{1 \\ 5}} \)
        \partsitem \( \colvec{-1 \\ 0 \\ 1}, \)\
          \( \set{\colvec{2 \\ 1 \\ 0},
                \colvec{1 \\ 0 \\ 1}} \)
        \partsitem \( \colvec{1 \\ 3 \\ 0}, \)\
          \( \set{\colvec{1 \\ 0 \\ 4},
                \colvec{2 \\ 1 \\ 5},
                \colvec{3 \\ 3 \\ 0},
                \colvec{4 \\ 2 \\ 1}} \)
        \partsitem \( \colvec{1 \\ 0 \\ 1 \\ 1}, \)\
          \( \set{\colvec{2 \\ 1 \\ 0 \\ 1},
                \colvec{3 \\ 0 \\ 0 \\ 2}} \)
      \end{exparts}
      \begin{answer}
        In each case we must decide if the vector is a linear combination
        of the vectors in the set.
        \begin{exparts}
          \partsitem Yes.
            Solve
            \begin{equation*}
              c_1\colvec{1 \\ 4}+c_2\colvec{1 \\ 5}=\colvec{2 \\ 3}
            \end{equation*}
            with
            \begin{equation*}
              \begin{amat}[r]{2}
                1  &1  &2  \\
                4  &5  &3
              \end{amat}
              \grstep{-4\rho_1+\rho_2}
              \begin{amat}[r]{2}
                1  &1  &2  \\
                0  &1  &-5
              \end{amat}
            \end{equation*}
            to conclude that there are $c_1$ and $c_2$ giving the combination. 
          \partsitem No.
            The reduction
            \begin{align*}
              \begin{amat}[r]{2}
                2  &1  &-1 \\
                1  &0  &0  \\
                0  &1  &1
              \end{amat}
              &\grstep{-(1/2)\rho_1+\rho_2}
              \begin{amat}[r]{2}
                2  &1     &-1 \\
                0  &-1/2  &1/2  \\
                0  &1     &1
              \end{amat}                        \\
              &\grstep{2\rho_2+\rho_3}
              \begin{amat}[r]{2}
                2  &1     &-1 \\
                0  &-1/2  &1/2  \\
                0  &0     &2
              \end{amat}
            \end{align*}
            shows that
            \begin{equation*}
              c_1\colvec{2 \\ 1 \\ 0}+c_2\colvec{1 \\ 0 \\ 1}
                =\colvec{-1 \\ 0 \\ 1}
            \end{equation*}
            has no solution.
          \partsitem Yes.
            We want to know if there are $c_i$'s giving this combination.
            \begin{equation*}
              \colvec{1 \\ 3 \\ 0}=
              c_1\colvec{1 \\ 0 \\ 4}+
              c_2\colvec{2 \\ 1 \\ 5}+
              c_3\colvec{3 \\ 3 \\ 0}+
              c_4\colvec{4 \\ 2 \\ 1}
            \end{equation*}
            The reduction
            \begin{equation*}
              % \begin{amat}[r]{4}
              %   1  &2  &3  &4  &1  \\
              %   0  &1  &3  &2  &3  \\
              %   4  &5  &0  &1  &0
              % \end{amat}
              \grstep{-4\rho_1+\rho_3}
              % \begin{amat}[r]{4}
              %   1  &2  &3  &4  &1  \\
              %   0  &1  &3  &2  &3  \\
              %   0  &-3 &-12&-15&-4
              % \end{amat}                   \\
              \repeatedgrstep{3\rho_2+\rho_3}
              \begin{amat}[r]{4}
                1  &2  &3  &4  &1  \\
                0  &1  &3  &2  &3  \\
                0  &0  &-3 &-9 &5
              \end{amat}
            \end{equation*}
            shows that there are infinitely many 
            four-tuples $(c_1, c_2, c_3, c_4)$.
            \begin{equation*}
              \set{\colvec{c_1 \\ c_2 \\ c_3 \\ c_4}=
                   \colvec{-10 \\ 8 \\ -5/3 \\ 0}+
                   \colvec{-9 \\ 7 \\ -3 \\ 1}c_4
                    \suchthat c_4\in\Re}
            \end{equation*}
          \partsitem No.
            Even without setting up the reduction, we can see from 
            the third components that no combination will suffice.
        \end{exparts} 
      \end{answer}
  \item 
     Prove that any linear system with a nonsingular matrix of 
     coefficients has a solution, and that the solution is unique.
    \begin{answer}
        Because the matrix of coefficients is nonsingular, Gauss's Method
        produces an echelon form system where each variable leads an equation.
        Back substitution gives a unique solution.

      (Another way to see that the solution is unique is to note that
      with a nonsingular matrix of coefficients the associated
      homogeneous system has a unique solution, by definition.
      Since the general solution is the sum of a particular solution with
      each homogeneous solution, the general solution has 
      at most one element.)
     \end{answer}
  \item 
    In the
    proof of
    \nearbylemma{le:HomoSltnSpanVecs},
    what happens if there are no non-\( 0=0 \) equations?
    \begin{answer}
      The set of equations make no restrictions on which 
      $n$-tuples are in the solution set.
      So the solution set is all of \( \Re^n \)\!. 
      We can express it in the required form as below
      (no variable leads an equation, so every variable is free).
      \begin{equation*}
        \set{x_1\colvec{1 \\ 0 \\ \vdotswithin{1} \\ 0}
             +x_2\colvec{0 \\ 1 \\ \vdotswithin{0} \\ 0}
             +\cdots
             +x_n\colvec{0 \\ 0 \\ \vdotswithin{0} \\ 1}
             \suchthat x_1,\ldots,x_n\in\Re}
      \end{equation*}  
     \end{answer}
  \recommended \item 
    Prove that if \( \vec{s} \) and \( \vec{t} \)
    satisfy a system of homogeneous equations then so do these vectors.
    \begin{exparts*}
      \partsitem \( \vec{s}+\vec{t} \)
      \partsitem \( 3\vec{s} \)
      \partsitem \( k\vec{s}+m\vec{t} \) for \( k,m\in\Re \)
    \end{exparts*}
    What's wrong with this argument: ``These three show that if a 
    system of homogeneous
    equations has one solution then it has many solutions\Dash any multiple of 
    a solution is another solution, and any sum of solutions is a solution
    also\Dash so there are no
    homogeneous systems with exactly one solution.''?
    \begin{answer}
      Let \( \vec{s},\vec{t} \) be members of \( \Re^n \) and write them as here.
      \begin{equation*}
        \vec{s}=\colvec{s_1 \\ \vdotswithin{s_1} \\ s_n}
          \qquad
        \vec{t}=\colvec{t_1 \\ \vdotswithin{t_1} \\ t_n}
      \end{equation*}
      Also let \( a_{i,1}x_1+\cdots+a_{i,n}x_n=0 \) be the \( i \)-th equation
      in the system.
      \begin{exparts}
        \partsitem The check is easy.
          The components of $\vec{s}+\vec{t}$ are $s_1+t_1$, $s_2+t_2$, etc.
          Substitute them into the equation, distribute, and regroup.
          \begin{multline*}
            a_{i,1}(s_1+t_1)+\cdots+a_{i,n}(s_n+t_n)      \\
            =
            (a_{i,1}s_1+\cdots+a_{i,n}s_n)
            +(a_{i,1}t_1+\cdots+a_{i,n}t_n)           
            =
            0+0
          \end{multline*}
        \partsitem This is similar to the prior one.
          \begin{equation*}
            a_{i,1}(3s_1)+\cdots+a_{i,n}(3s_n)
            =3(a_{i,1}s_1+\cdots+a_{i,n}s_n)
            =3\cdot 0=0
          \end{equation*}
        \partsitem This one is not much harder.
          \begin{multline*}
            a_{i,1}(ks_1+mt_1)+\cdots+a_{i,n}(ks_n+mt_n)  \\
            \begin{split}
              &=
              k(a_{i,1}s_1+\cdots+a_{i,n}s_n)
              +m(a_{i,1}t_1+\cdots+a_{i,n}t_n)      \\         
              &=
              k\cdot 0+m\cdot 0
              =0
            \end{split}
          \end{multline*}
      \end{exparts}  
     What is wrong with that argument is that any linear combination 
     involving only the 
     zero vector yields the zero vector.
   \end{answer}
  \item
    Prove that if a system with only rational coefficients
    and constants
    has a solution then it has at least one all-rational solution.
    Must it have infinitely many?
    \begin{answer}
      First the proof.
      If the system has all rational coefficients then
      Gauss's Method with back substitution will use only rationals (e.g.,
      \( -(m/n)\rho_i+\rho_j \)).
      Thus we can express the solution set using only rational numbers as
      the components of each vector.
      Therefore, the particular solution is all rational.

      As to the question about infinitely many, 
      there are infinitely many rational vector solutions if and only if the
      associated homogeneous system has infinitely many 
      real vector solutions.
      That's because setting any parameters to be rationals will produce an
      all-rational solution.  
   \end{answer}
\end{exercises}

% \endinput




\subsection{Comparing Set Descriptions} \label{subsec:ComparingSetDescriptions}
\emph{This material is optional, not required for later work.}

A set can have many different descriptions.
Here are two different ways to describe one set, 
the collection of vectors with components in ratio $1\mathbin{:}2\mathbin{:}3$.
\begin{equation*}
  \set{\colvec{1 \\ 2 \\ 3}a\suchthat a\in\Re}
  \qquad
  \set{\colvec{2 \\ 4 \\ 6}b\suchthat b\in\Re}
\end{equation*}
For instance, this set contains the vector below
on the left (take $a=5$ and $b=5/2$), 
but not the one on the right
\begin{equation*}
  \colvec{5 \\ 10 \\ 15}
  \qquad
  \colvec{4 \\ 8 \\ 11}
\end{equation*}
(the first component gives $a=4$ but that clashes with the third component).
So if we compute one description and the answer key has another, we want to 
be able to verify that the two are the same set of vectors.

Sets are the same if and only if they have the same elements.
Usually we demonstrate that $S=T$ by showing 
mutual inclusion, \index{sets!mutual inclusion}\index{mutual inclusion}
that any member of $S$ is also a member of~$T$\!, and 
that any member of $T$ is also a member of~$S$.\appendrefs{set equality}

\begin{example}
To show that these two are equal
\begin{equation*}
  S=
  \set{\colvec{1 \\ -1 \\ 0}c+\colvec{1 \\ 1 \\ 0}d\suchthat c,d\in\Re}
  \qquad
  T=
  \set{\colvec{4 \\ 1 \\ 0}m+\colvec{-1 \\ -3 \\ 0}n\suchthat m,n\in\Re}
\end{equation*}
we we will first show that $S\subseteq T$ and then that $T\subseteq S$. 

We start the argument's $S\subseteq T$ half with a model of how to proceed.
Fix a vector $\vec{v}\in S$ by picking \( c=1 \) and \( d=1 \), making 
$\vec{v}$ have components $2$, $0$, and~$0$.
To show that $\vec{v}$ is expressible in the form for $T$\!, 
we look for \( m \) and~$n$ such that this holds.
\begin{equation*}
  \colvec{4 \\ 1 \\ 0}m
  +\colvec{-1 \\ -3 \\ 0}n
  =\colvec{2 \\ 0 \\ 0}
  \qquad
  \begin{linsys}{2}
    4m  &-  &n  &=  &2  \\
    1m  &-  &3n &=  &0  \\
        &   &0  &=  &0 
  \end{linsys}  
\end{equation*}
Gauss's Method shows that $m=6/11$ and $n=2/11$ work, so $\vec{v}\in T$\!. 

In general,
to show that any vector from \( S \) is a member of \( T \)\!, 
fix a pair \( c, d \) and find \( m \) and~\( n \) satisfying these.
\begin{equation*}
  \colvec{4 \\ 1 \\ 0}m
  +\colvec{-1 \\ -3 \\ 0}n
  =\colvec{c+d \\ -c+d \\ 0}
  \qquad
  \begin{linsys}{2}
    4m  &-  &n  &=  &c+d\hfill  \\
     m  &-  &3n &=  &-c+d\hfill  \\
        &   &0  &=  &0\hfill  
  \end{linsys}
\end{equation*}
Gauss's Method
\begin{equation*}
  \begin{linsys}{2}
    4m  &-  &n  &=  &c+d\hfill  \\
     m  &-  &3n &=  &-c+d\hfill  
  \end{linsys}
  \grstep{-(1/4)\rho_1+\rho_2}
  \begin{linsys}{2}
    4m  &-  &n        &=  &c+d\hfill            \\
        &   &-(11/4)n &=  &-(5/4)c+(3/4)d\hfill  
   \end{linsys}
\end{equation*}
gives that \( n=(5/11)c-(3/11)d \) and \( m=(4/11)c+(2/11)d \).
Therefore, any member of $S$ is a member of~$T$.
\begin{equation*}
   \colvec{c+d \\ -c+d \\ 0}
   =\colvec{4 \\ 1 \\ 0}\cdot\biggl(\frac{4}{11}c+\frac{2}{11}d\biggr)+
   \colvec{-1 \\ -3 \\ 0}\cdot\biggl(\frac{5}{11}c-\frac{3}{11}d\biggr)
\end{equation*}

For the other inclusion, \( T\subseteq S \), we
want to show that for any member of~$T$\!, 
that is, if we fix $m$ and~$n$, 
then there are appropriate
$c$ and~$d$.
\begin{equation*}
   \begin{linsys}{2}
     c  &+ &d  &= &4m-n\hfill \\
    -c  &+ &d  &= &m-3n\hfill 
   \end{linsys}
   \grstep{\rho_1+\rho_2}
   \begin{linsys}{2}
     c  &+ &d  &= &4m-n\hfill \\
        &  &2d &= &5m-4n\hfill 
    \end{linsys}
\end{equation*}
We get that \( d=(5/2)m-2n \) and \( c=(3/2)m+n \).
Thus any vector described as in \( T \)
can also be expressed in the right form for~\( S \).
\begin{equation*}
  \colvec{4 \\ 1 \\ 0}m+\colvec{-1 \\ -3 \\ 0}n
% \end{equation*}
% \begin{equation*}
  =
  \colvec{1 \\ -1 \\ 0}\cdot\biggl(\frac{3}{2}m+n\biggr)
    +\colvec{1 \\ 1 \\ 0}\cdot\biggl(\frac{5}{2}m-2n\biggr)
\end{equation*}
\end{example}

\begin{example}
We will show this by mutual inclusion.
\begin{equation*}
  S=\set{\colvec{6 \\ 0}+\colvec{-4 \\ 2}y\suchthat y\in\R}
  =\set{\colvec{2 \\ 2}+\colvec{2 \\ -1}k\suchthat k\in\R}=T
\end{equation*}

For $S\subseteq T$, we will verify that any $\vec{s}\in S$ is also a member
of~$T$.
Considering $\vec{s}$ is the same as fixing a~$y$.
We then look for a matching~$k$.
\begin{equation*}
  \begin{linsys}{2}
     2  &+  &2k  &=  &6-4y\hfill  \\
     2  &-  &k   &=  &2y\hfill  
  \end{linsys}
  \quad\Longrightarrow\quad
  \begin{linsys}{1}
     2k  &=  &4-4y\hfill  \\
     -k  &=  &-2+2y\hfill  
  \end{linsys}
\end{equation*}
From~$y$ we can compute $k$ as $k=2-2y$.

For $T\subseteq S$, we will verify that any $\vec{t}\in T$ is a member of~$S$.
For each~$k$, find a~$y$.
\begin{equation*}
  \begin{linsys}{2}
    6  &-  &4y  &= &2+2k\hfill\\
       &   &2y  &= &2-k\hfill  \\
  \end{linsys}
  \quad\Longrightarrow\quad
  \begin{linsys}{1}
    4y  &= &-4+2k\hfill\\
    2y  &= &2-k\hfill  \\
  \end{linsys}
\end{equation*}
Given~$k$, we can take $y=1-(k/2)$.
\end{example}





\begin{exercises}
  % \item 
  %   Decide if the vector is a member of the set.
  %   \begin{exparts}
  %     \partsitem $\colvec{2 \\ 3}$, 
  %        $\set{\colvec{1 \\ 2}k\suchthat k\in\Re}$
  %     \partsitem $\colvec{-3 \\ 3}$, 
  %        $\set{\colvec{1 \\ -1}k\suchthat k\in\Re}$
  %     \partsitem $\colvec{-3 \\ 3 \\ 4}$, 
  %            $\set{\colvec{1 \\ -1 \\ 2}k\suchthat k\in\Re}$
  %     \partsitem $\colvec{-3 \\ 3 \\ 4}$, 
  %            $\set{\colvec{1 \\ -1 \\ 2}k+\colvec{0 \\ 0 \\ 2}m
  %               \suchthat k,m\in\Re}$
  %     \partsitem $\colvec{1 \\ 4 \\ 14}$, 
  %            $\set{\colvec{2 \\ 2 \\ 5}k+\colvec{-1 \\ 0 \\ 2}m
  %               \suchthat k,m\in\Re}$
  %     \partsitem $\colvec{1 \\ 4 \\ 6}$, 
  %            $\set{\colvec{2 \\ 2 \\ 5}k+\colvec{-1 \\ 0 \\ 2}m
  %               \suchthat k,m\in\Re}$
  %   \end{exparts}
  %   \begin{answer}
  %     \begin{exparts}
  %       \partsitem No.
  %       \partsitem Yes.
  %       \partsitem No.
  %       \partsitem Yes.
  %       \partsitem Yes; use Gauss's Method to get $k=4$ and $m=-3$.
  %       \partsitem No; use Gauss's Method to conclude that there is no solution.
  %     \end{exparts}
  %   \end{answer}
  \item 
     Produce two descriptions of this set that are different from the given one. 
     \begin{equation*}
       \set{\colvec{2 \\ -5}k\suchthat k\in\Re}
     \end{equation*}
     \begin{answer}
       One easy thing to do is to double and triple the vector.
       \begin{equation*}
         \set{\colvec{4 \\ -10}k\suchthat k\in\Re}
         \qquad
         \set{\colvec{6 \\ -15}k\suchthat k\in\Re}
       \end{equation*}
      \end{answer}
  \recommended\item Show that the two sets of vectors are equal by
     first showing that $S$ is a subset of~$T$\!,
     that any member of~$S$ is a member of~$T$\!.
     Finish by showing that $T\subseteq S$.
     \begin{exparts}
       \partsitem $S=\set{\colvec{3 \\ 0}x+\colvec{0 \\ 2}y\suchthat x,y\in\R}$,
         $T=\set{\colvec{1 \\ 0}a+\colvec{0 \\ 1}b\suchthat a, b\in\R}$
       \partsitem $S=\set{\colvec{5 \\ 1 \\ 0}x+\colvec{2 \\ -1 \\ 0}y\suchthat x,y\in\R}$,
         $T=\set{\colvec{1 \\ 0 \\ 0}a+\colvec{0 \\ 1 \\ 0}b\suchthat a, b\in\R}$
     \end{exparts}
     \begin{answer}
       \begin{exparts}
         \partsitem
            First we verify that $S\subseteq T$\!, 
            that any member of~$S$ is a member of~$T$\!.
            We will show that for all $x,y$ there are suitable $a$ and~$b$
            making this true.
            \begin{equation*}
              \colvec{1 \\ 0}a+\colvec{0 \\ 1}b
               =
              \colvec{3 \\ 0}x+\colvec{0 \\ 2}y
              \tag{$*$}
            \end{equation*}
            In brief, fixing $x$ and~$y$, we must produce
            real numbers~$a,b\in\R$ satisfying these.
            \begin{equation*}
              \begin{linsys}{2}
                a  &    &   &=  &3x  \\
                   &\   &b  &=  &2y
              \end{linsys}
            \end{equation*}
            Obviously we can take $a$ to be~$3x$ and~$b$ to be $2y$.

            The second half is to show that $T\subseteq S$, 
            that any member of~$T$ is also a member of~$S$.
            Here we fix $a$ and~$b$ and look for 
            real numbers $x$ and~$y$ satisfying this system.
            \begin{equation*}
              \begin{linsys}{2}
                3x &    &   &=  &a  \\
                   &\   &2y &=  &b
              \end{linsys}
            \end{equation*}
            Again we don't need Gauss's Method.
            We can take $x$ to be~$a/3$ and $y$ to be~$b/2$.
         \partsitem
            The first half is to verify that $S\subseteq T$\!, 
            that any member of~$S$ is a member of~$T$.
            We will show that for any $x$ and~$y$ there are suitable
            real numbers $a$ and~$b$.
            So fix $x$ and~$y$. 
            We need
            real numbers~$a,b\in\R$ satisfying this.
            \begin{equation*}
              \begin{linsys}{2}
                a  &    &   &=  &5x+2y  \\
                   &\   &b  &=  &x-y
              \end{linsys}
            \end{equation*}
            This half of the proof is easy:~take $a$ to be~$5x+2y$ 
            and $b$ to be~$x-y$.

            The second half is to show that  $T\subseteq S$, 
            that any member of~$T$ is also a member of~$S$.
            We fix $a$ and~$b$, and look for 
            real numbers $x$ and~$y$ satisfying these relationships.
            \begin{equation*}
              \begin{linsys}{2}
                5x &+   &2y &=  &a  \\
                 x &-   &y  &=  &b
              \end{linsys}
            \end{equation*}
            The computation is straightforward.
            \begin{align*}
              \begin{amat}{2}
                5  &2   &a  \\ 
                1  &-1  &b  \\ 
              \end{amat}
              &\grstep{(-1/5)\rho_{1}+\rho_{2}}
                \begin{amat}{2}
                  5  &2     &a  \\ 
                  0  &-7/5  &(-1/5)a+b  \\ 
                \end{amat}                       %  \\
              % &\grstep[(-5/7)\rho_{2}]{(1/5)\rho_{1}}
              %   \begin{amat}{2}
              %     1  &2/5  &(1/5)a  \\ 
              %     0  &1    &(1/7)a-(5/7)b  \\ 
              %   \end{amat}                       \\
              % &\grstep{(-2/5)\rho_{2}+\rho_{1}}
              %   \begin{amat}{2}
              %     1  &0  &(1/7)a+(2/7)b  \\ 
              %     0  &1  &(1/7)a-(5/7)b   \\ 
              %   \end{amat}
            \end{align*}
            Back substitution gives that 
            we can take $x$ to be $(1/7)a+(2/7)b$ and $y$ to be $(1/7)a-(5/7)b$.
       \end{exparts}
     \end{answer}
  \recommended \item 
    Show that these sets are equal
    \begin{equation*}
      S=\set{\colvec{1 \\ 4 \\ 1 \\ 1}
           +\colvec{-1 \\ 0 \\ 1 \\ 0}m\suchthat m\in\Re  }
      \qquad
      T=\set{\colvec{0 \\ 4 \\ 2 \\ 1}
           +\colvec{-1 \\ 0 \\ 1 \\ 0}k\suchthat k\in\Re  },
    \end{equation*}
    and that both describe the solution set of this system.
    \begin{equation*}  
       \begin{linsys}{4}
         x  &-  &y  &+  &z  &+  &w  &=  &-1  \\
            &   &y  &   &   &-  &w  &=  &3   \\
         x  &   &   &+  &z  &+  &2w &=  &4
       \end{linsys}
    \end{equation*}
    \begin{answer}
      We first show that $S=T$ by mutual inclusion. 

      For $S\subseteq T$ we want that any vector in the form for~$S$
      can also be expressed in the form for~$T$\!.
      That is, we fix~$m$ and look for~$k$ to satisfy these.
      \begin{equation*}
        \begin{linsys}{2}
            &  &-k  &=  &1-m\hfill \\
         4  &  &    &=  &4\hfill \\
         2  &+ &k   &=  &1+m\hfill \\
         1  &  &    &=  &1\hfill \\
        \end{linsys}
        \qquad\Longrightarrow\qquad
        \begin{linsys}{2}
          \   &   &-k  &=  &1-m\hfill \\
              &   &0   &=  &0\hfill \\
              &\  &k   &=  &-1+m\hfill \\
              &   &0   &=  &0\hfill \\
        \end{linsys}
      \end{equation*}
      It suffices to take $k$ to be $-1+m$.
      So $S\subseteq T$.

      As to $T\subseteq S$, fix~$k$ and look for a suitable~$m$.
      \begin{equation*}
        \begin{linsys}{2}
          1  &-  &m  &= &-k\hfill  \\
          4  &   &   &= &4\hfill  \\
          1  &+  &m  &= &2+k\hfill  \\
          1  &   &   &= &1\hfill  \\
        \end{linsys}
        \qquad\Longrightarrow\qquad
        \begin{linsys}{2}
           \   &\    &-m &= &-1-k\hfill  \\
               &   &0  &= &0\hfill  \\
               &   &m  &= 1+k\hfill  \\
               &   &0  &= &0\hfill  \\
        \end{linsys}
      \end{equation*}
      Take $m$ to be $1+k$.
      Thus $T\subseteq S$.       

      That system reduces like this,
      \begin{equation*}
        \begin{amat}{4}
          1  &-1  &1  &1  &-1  \\ 
          0  &1  &0  &-1  &3  \\ 
          1  &0  &1  &2  &4  \\ 
        \end{amat}
        \grstep{-\rho_{1}+\rho_{3}}
        % \begin{amat}{4}
        %   1  &-1  &1  &1  &-1  \\ 
        %   0  &1  &0  &-1  &3  \\ 
        %   0  &1  &0  &1  &5  \\ 
        % \end{amat}
        \repeatedgrstep{-\rho_{2}+\rho_{3}}
        \begin{amat}{4}
          1  &-1  &1  &1  &-1  \\ 
          0  &1  &0  &-1  &3  \\ 
          0  &0  &0  &2  &2  \\ 
        \end{amat}
      \end{equation*}
      giving this solution set.
      \begin{equation*}
        U=\set{\colvec{2 \\ 4 \\ 0 \\ 1}+ \colvec{-1 \\ 0 \\ 1 \\ 0}z
             \suchthat z\in\R}
      \end{equation*}
      If we show either $S=U$ or $T=U$ then we are done.
      We will do the first, just an arbitrary choice.

      For $S\subseteq U$ we fix an element of~$S$ and look for a way
      to express it in the form used in the description of~$U$.
      That is, we fix~$m$ and look for~$z$ to satisfy these.
      \begin{equation*}
        \begin{linsys}{2}
         2  &  &-z  &=  &1-m\hfill \\
         4  &  &    &=  &4\hfill \\
         0  &+ &z   &=  &1+m\hfill \\
         1  &  &    &=  &1\hfill \\
        \end{linsys}
        \qquad\Longrightarrow\qquad
        \begin{linsys}{2}
          \   &   &-z  &=  &-1-m\hfill \\
              &   &0   &=  &0\hfill \\
              &\  &z   &=  &1+m\hfill \\
              &   &0   &=  &0\hfill \\
        \end{linsys}
      \end{equation*}
      It suffices to take $z$ to be $1+m$.
      So $S\subseteq U$.

      For $U\subseteq S$, fix~$z$ and look for a suitable~$m$.
      \begin{equation*}
        \begin{linsys}{2}
          1  &-  &m  &= &2-z\hfill  \\
          4  &   &   &= &4\hfill  \\
          1  &+  &m  &= &z\hfill  \\
          1  &   &   &= &1\hfill  \\
        \end{linsys}
        \qquad\Longrightarrow\qquad
        \begin{linsys}{2}
           \   &\    &-m &= &1-z\hfill  \\
               &   &0  &= &0\hfill  \\
               &   &m  &= -1+z\hfill  \\
               &   &0  &= &0\hfill  \\
        \end{linsys}
      \end{equation*}
      Take $m$ to be $-1+z$.
      Thus $U\subseteq S$.  

      \authorsNote{Of course, all three sets $S$, $T$, and~$U$ have the same
        vector associated with their parameter.
        And in all three the vector not associated with the parameter 
        is a particular solution of the system.
        In the special case of solution sets for linear systems there 
        are faster ways to check set equality then mutual inclusion,
        which works in the general case of any two sets.
        We will see more in the next chapter.}     
    \end{answer}
  \item
    These sets are not equal.
    \begin{equation*}
      S=\set{\colvec{x+y \\ x \\ y}\suchthat x,y\in\Re}
      \qquad
      T=\set{\colvec{m+p \\ n \\ p}\suchthat m,n,p\in\Re}
    \end{equation*}
    Show that $S$ is a proper subset, $S\subseteq T$ but $T\nsubseteq S$.
    \begin{answer}
      First we argue that $S\subseteq T$.
      Fix a vector from \( S \). 
      To express it in the form for \( T \) we 
      solve for appropriate $m$, $n$, and~$p$.
      \begin{equation*}
        \begin{linsys}{3}
          m  & &   &+  &p  &=  &x+y\hfill  \\
             & &n  &   &   &=  &x\hfill   \\
             & &   &   &p  &=  &y\hfill    
        \end{linsys}
      \end{equation*}
      By eye we see that we can take $p$ to be~$y$, take $n$ to be~$x$,
      and take $m$ to be~$x$.
      % So we can express any $\vec{v}\in S$
      % \begin{equation*}
      %   \vec{v}=
      %   \colvec{1 \\ 1 \\ 0}x+
      %   \colvec{1 \\ 0 \\ 1}y
      % \end{equation*}
      % as a member of \( T \) with
      % \( p=y \) and \( m=n=x \).
      % \begin{equation*}
      %   \vec{v}=
      %   \colvec{1 \\ 0 \\ 0}x+
      %   \colvec{0 \\ 1 \\ 0}2x+
      %   \colvec{1 \\ 0 \\ 1}y.
      % \end{equation*}

      For the $T \nsubseteq S$ argument, it suffices to note that 
      \begin{equation*}
        \vec{w}=\colvec{1 \\ 2 \\ 3}\in T 
      \end{equation*}
      (where $m=-2$, $n=2$, and~$p=3$).
      But $\vec{w}\notin S$ because members of~$S$ have that the second and
      third components add to the first, and $2+3\neq 1$.

      \authorsNote{
        We can wonder where the approach in the first part fails.
        Attempting to show that $T$ is a subset of~$S$, we want to show
        that any member of~$T$ can be expressed in the form of a member of~$S$.
        Fix $m$, $n$, and $p$ for the member of~$T$ and look for some 
        $x$ and~$y$ marking it as a member also of~$S$.
        \begin{align*}
          \begin{linsys}{2}
            x  &+  &y  &=  &m+p\hfill  \\
            x  &   &   &=  &n\hfill    \\
               &   &y  &=  &p\hfill    
          \end{linsys}
          &\grstep{-\rho_1+\rho_2}
            \begin{linsys}{2}
              x  &+  &y  &=  &m+p\hfill  \\
                 &   &-y  &=  &-m+n-p\hfill \\
                 &   &y  &=  &p\hfill    
               \end{linsys}                                  \\
          &\grstep{\rho_2+\rho_3}
            \begin{linsys}{2}
              x  &+  &y  &=  &m+p\hfill  \\
                 &   &-y &=  &-m+n-p\hfill \\
                 &   &0  &=  &-m+n\hfill 
            \end{linsys}
        \end{align*}
        Thus the only vectors
        \begin{equation*}
          \colvec{m+p \\ n \\ p}\in T
        \end{equation*}
        representable in the form
        \begin{equation*}
          \colvec{x+y \\ x \\ y}\in S
        \end{equation*}
        are those where \( 0=-m+n \).
        For instance, taking $m=n=1$ and $p=0$ gives that 
        \begin{equation*}
          \colvec{1 \\ 1 \\ 0}
        \end{equation*}
        is a member of \( T \) but not~\( S \).
      }
    \end{answer}
  \recommended \item 
    Decide if the sets are equal.
    \begin{exparts}
      \partsitem \( \set{\colvec{1 \\ 2}
                        +\colvec{0 \\ 3}t
                     \suchthat t\in\Re} \)
            and
            \( \set{\colvec{1 \\ 8}
                        +\colvec{0 \\ -1}s
                     \suchthat s\in\Re} \)
      \partsitem \( \set{\colvec{1 \\ 3 \\ 1}t
                        +\colvec{2 \\ 1 \\ 5}s
                     \suchthat t,s\in\Re} \)
            and
            \( \set{\colvec{4 \\ 7 \\ 7}m
                        +\colvec{-4 \\ -2 \\ -10}n
                     \suchthat m,n\in\Re} \)
      \partsitem \( \set{\colvec{1 \\ 2}t
                     \suchthat t\in\Re} \)
            and
            \( \set{\colvec{2 \\ 4}m
                        +\colvec{4 \\ 8}n
                     \suchthat m,n\in\Re} \)
      \partsitem \( \set{\colvec{1 \\ 0 \\ 2}s
                        +\colvec{-1 \\ 1 \\ 0}t
                     \suchthat s,t\in\Re} \)
            and
            \( \set{\colvec{-1 \\ 1 \\ 1}m
                        +\colvec{0 \\ 1 \\ 3}n
                     \suchthat m,n\in\Re} \)
      \partsitem \( \set{\colvec{1 \\ 3 \\ 1}t
                        +\colvec{2 \\ 4 \\ 6}s
                     \suchthat t,s\in\Re} \)
            and
            \( \set{\colvec{3 \\ 7 \\ 7}t
                        +\colvec{1 \\ 3 \\ 1}s
                     \suchthat t,s\in\Re} \)
    \end{exparts}
    \begin{answer}
      For each item, we call the first set \( S_1 \) and the
      other \( S_2 \).
     \begin{exparts}
      \partsitem They are equal.

        To see that \( S_1\subseteq S_2 \), we must show that any
        element of the
        first set is in the second, that is, for any vector of the form
        \begin{equation*}
          \vec{v}=\colvec{1 \\ 2}
                  +\colvec{0 \\ 3}t
        \end{equation*}
        there is an appropriate \( s \) such that
        \begin{equation*}
          \vec{v}=\colvec{1 \\ 8}
                  +\colvec{0 \\ -1}s.
        \end{equation*}
        Restated, we fix \( t \) and look for \( s \) so that this holds.
        \begin{equation*}
          \begin{linsys}{2}
            1  &+  &0s  &=  &1+0t\hfill  \\
            8  &-  &1s  &=  &2+3t\hfill  
          \end{linsys}
          \qquad\Longrightarrow\qquad
          \begin{linsys}{1}
            0  &= &0 \hfill \\
            s  &= &6-3t\hfill
          \end{linsys}
        \end{equation*}
        That is,
        \begin{equation*}
          \colvec{1 \\ 2}
          +\colvec{0 \\ 3}t
          =\colvec{1 \\ 8}
          +\colvec{0 \\ -1}(6-3t)
        \end{equation*}
        and so any vector in the form for \( S_1 \) can be put
        in the form
        needed for inclusion in \( S_2 \).

        For \( S_2\subseteq S_1 \), we look for \( t \) so that
        these equations hold.
        \begin{equation*}
          \begin{linsys}{2}
            1  &+  &0t  &=  &1+0s\hfill  \\
            2  &+  &3t  &=  &8-1s\hfill  
          \end{linsys}
          \qquad\Longrightarrow\qquad
          \begin{linsys}{1}
            0 &= &0\hfill   \\
            t &= &2-(1/3)s
          \end{linsys}
        \end{equation*}
        Thus 
        \begin{equation*}
          \colvec{1 \\ 8}
          +\colvec{0 \\ -1}s
          =\colvec{1 \\ 2}
          +\colvec{0 \\ 3}(2-(1/3)s).
        \end{equation*}
        and $S_2\subseteq S_1$,
        and so the two are equal.
      \partsitem These two are equal.

        To show that \( S_1\subseteq S_2 \), we check that for any \( t,s \)
        we can find an appropriate \( m,n \) so that these hold.
        \begin{equation*}
          \begin{linsys}{2}
           4m  &-  &4n   &=  &1t+2s\hfill  \\
           7m  &-  &2n   &=  &3t+1s\hfill  \\
           7m  &-  &10n  &=  &1t+5s\hfill  
          \end{linsys}
        \end{equation*}
        Use Gauss's Method
        \begin{align*}
          \begin{amat}{2}
            4  &-4  &1t+2s  \\
            7  &-2  &3t+1s  \\
            7  &-10 &1t+5s
          \end{amat}
          &\grstep[(-7/4)\rho_1+\rho_3]{(-7/4)\rho_1+\rho_2}
          \begin{amat}{2}
            4  &-4  &1t+2s           \\
            0  &5   &(5/4)t-(10/4)s  \\
            0  &-3  &-(3/4)t+(6/4)s
          \end{amat}                              \\
          &\grstep{(3/5)\rho_2+\rho_3}
          \begin{amat}{2}
            4  &-4  &1t+2s           \\
            0  &5   &(5/4)t-(10/4)s  \\
            0  &0   &0
          \end{amat}
        \end{align*}
        to conclude that
        \begin{equation*}
          \colvec{1 \\ 3 \\ 1}t
          +\colvec{2 \\ 1 \\ 5}s
          =\colvec{4 \\ 7 \\ 7}((1/2)t)
          +\colvec{-4 \\ -2 \\ -10}((1/4)t-(1/2)s)
        \end{equation*}
        and so \( S_1\subseteq S_2 \).

        For \( S_2\subseteq S_1 \), solve
        \begin{equation*}
          \begin{linsys}{2}
           1t  &+  &2s   &=  &4m-4n\hfill  \\
           3t  &+  &1s   &=  &7m-2n\hfill  \\
           1t  &+  &5s   &=  &7m-10n\hfill  
          \end{linsys}
        \end{equation*}
        with Gaussian reduction
        \begin{align*}
          \begin{amat}{2}
            1  &2   &4m-4n  \\
            3  &1   &7m-2n  \\
            1  &5   &7m-10n
          \end{amat}
          &\grstep[-\rho_1+\rho_3]{-3\rho_1+\rho_2}
          \begin{amat}{2}
            1  &2   &4m-4n  \\
            0  &-5  &-5m+10n\\
            0  &3   &3m-6n
          \end{amat}                                    \\
          &\grstep{(3/5)\rho_2+\rho_3}
          \begin{amat}{2}
            1  &2   &4m-4n  \\
            0  &-5  &-5m+10n\\
            0  &0   &0
          \end{amat}
        \end{align*}
        to get
        \begin{equation*}
          \colvec{4 \\ 7 \\ 7}m
          +\colvec{-4 \\ -2 \\ -10}n
          =\colvec{1 \\ 3 \\ 1}(2m)
          +\colvec{2 \\ 1 \\ 5}(m-2n)
        \end{equation*}
        and so we can express any member of \( S_2 \) in the form needed for
        \( S_1 \).
      \partsitem These sets are equal.

        To prove that \( S_1\subseteq S_2 \), we must be able to solve
        \begin{equation*}
          \begin{linsys}{2}
           2m  &+  &4n  &=  &1t\hfill  \\
           4m  &+  &8n  &=  &2t\hfill  
          \end{linsys}
        \end{equation*}
        for \( m \) and \( n \) in terms of \( t \).
        Apply Gaussian  reduction
        \begin{equation*}
          \begin{amat}{2}
            2  &4   &1t  \\
            4  &8   &2t
          \end{amat}
          \grstep{-2\rho_1+\rho_2}
          \begin{amat}{2}
            2  &4   &1t  \\
            0  &0   &0
          \end{amat}
        \end{equation*}
        to conclude that
        any pair \( m,n \) where \( 2m+4n=t \) will do.
        For instance,
        \begin{equation*}
          \colvec{1 \\ 2}t
          =\colvec{2 \\ 4}((1/2)t)
          +\colvec{4 \\ 8}(0)
        \end{equation*}
        or
        \begin{equation*}
          \colvec{1 \\ 2}t
          =\colvec{2 \\ 4}((-3/2)t)
          +\colvec{4 \\ 8}(t).
        \end{equation*}
        Thus \( S_1\subseteq S_2 \).

        For \( S_2\subseteq S_1 \), we solve
        \begin{equation*}
          \begin{linsys}{2}
           1t  &=  &2m+4n\hfill  \\
           2t  &=  &4m+8n\hfill  
          \end{linsys}
        \end{equation*}
        with Gauss's Method
        \begin{equation*}
          \begin{amat}{1}
            1  &2m+4n  \\
            2  &4m+8n
          \end{amat}
          \grstep{-2\rho_1+\rho_2}
          \begin{amat}{1}
            1  &2m+4n  \\
            0  &0
          \end{amat}
        \end{equation*}
        to deduce that any vector in \( S_2 \) is also in \( S_1 \).
        \begin{equation*}
          \colvec{2 \\ 4}m
          +\colvec{4 \\ 8}n
          =\colvec{1 \\ 2}(2m+4n)\in S_1.
        \end{equation*}
      \partsitem Neither set is a subset of the other.

        For \( S_1\subseteq S_2 \) to hold we must be able to solve
        \begin{equation*}
          \begin{linsys}{2}
          -1m  &+  &0n   &=  &1s-1t\hfill  \\
           1m  &+  &1n   &=  &0s+1t\hfill  \\
           1m  &+  &3n   &=  &2s+0t\hfill  
          \end{linsys}
        \end{equation*}
        for \( m \) and \( n \) in terms of \( t \) and \( s \).
        Gauss's Method
        \begin{align*}
          \begin{amat}{2}
           -1  &0   &1s-1t  \\
            1  &1   &0s+1t  \\
            1  &3   &2s+0t
          \end{amat}
          &\grstep[\rho_1+\rho_3]{\rho_1+\rho_2}
          \begin{amat}{2}
           -1  &0   &1s-1t  \\
            0  &1   &1s+0t  \\
            0  &3   &3s-1t
          \end{amat}                        \\
          &\grstep{-3\rho_2+\rho_3}
          \begin{amat}{2}
           -1  &0   &1s-1t  \\
            0  &1   &1s+0t  \\
            0  &3   &0s-1t
          \end{amat}
        \end{align*}
        shows that $m=t-s$, and that $n=s$, and that $3n=-t$ so $n=-t/3$.
        We have that the only way there are matching $m$ and~$n$ is for those
        members of $S_1$ where $s=-t/3$.
        An example of a vector in both sets is where $t=3$ and $s=-1$.
        \begin{equation*}
          \colvec{1 \\ 0 \\ 2}\cdot (-1)
            +\colvec{-1 \\ 1 \\ 0}\cdot 3
            =\colvec{-4 \\ 3 \\ -2}
            =\colvec{-1 \\ 1 \\ 1}\cdot 4
            +\colvec{0 \\ 1 \\ 3}\cdot (-1)
        \end{equation*}
        However, if $s\neq -t/3$ then the vector is not expressible in the form 
        for~$S_2$.
        For example, we can take $t=3$ and~$s=0$.
        \begin{equation*}
          \colvec{1 \\ 0 \\ 2}\cdot 0
            +\colvec{-1 \\ 1 \\ 0}\cdot 3
            =\colvec{-3 \\ 3 \\ 0}
        \end{equation*}
        To double check that this is not a member of~$S_2$ we can use 
        a linear system.
        \begin{equation*}
          \colvec{-3 \\ 3 \\ 0}
            =\colvec{-1 \\ 1 \\ 1}\cdot m
            +\colvec{0 \\ 1 \\ 3}\cdot n
          \qquad
          \begin{linsys}{2}
            -m  &  &  &=  &-3  \\
             m  &+ &n &=  &3  \\
             m  &+ &3n&=  &0  \\
          \end{linsys}
        \end{equation*}
        Gauss's Method gives that there is no solution.
        \begin{equation*}
          \begin{amat}{2}
            -1  &0  &-3  \\ 
            1  &1  &3  \\ 
            1  &3  &0  \\ 
          \end{amat}
          \grstep[\rho_{1}+\rho_{3}]{\rho_{1}+\rho_{2}}
          % \begin{amat}{2}
          %   -1  &0  &-3  \\ 
          %   0  &1  &0  \\ 
          %   0  &3  &-3  \\ 
          % \end{amat}
          \repeatedgrstep{-3\rho_{2}+\rho_{3}}
          \begin{amat}{2}
            -1  &0  &-3  \\ 
            0  &1  &0  \\ 
            0  &0  &-3  \\ 
          \end{amat}
        \end{equation*}

        Having shown that \( S_1 \) is not a subset of \( S_2 \), we know that
        \( S_1\neq S_2 \). 
        So, strictly speaking, we need not go further.
        But for folks who tried to do the other half 
        we shall also show that \( S_2 \) is not a subset of~\( S_1 \).

        For \( S_2\subseteq S_1 \), we fix $m$ and~$n$ and look
        to solve for \( s \) and~\( t \).
        \begin{equation*}
          \begin{linsys}{2}
           s  &-  &t   &=  &-m\hfill  \\
              &   &t   &=  &m+n\hfill  \\
          2s  &   &    &=  &m+3n\hfill  
          \end{linsys}
        \end{equation*}
        Apply row reduction.
        \begin{align*}
          \begin{amat}{2}
            1  &-1  &-m  \\
            0  &1   &m+n  \\
            2  &0   &m+3n
          \end{amat}
          &\grstep{-2\rho_1+\rho_3}
          \begin{amat}{2}
            1  &-1  &-m  \\
            0  &1   &m+n  \\
            0  &2   &3m+3n
          \end{amat}                                    \\
          &\grstep{-2\rho_2+\rho_3}
          \begin{amat}{2}
            1  &-1  &-m  \\
            0  &1   &m+n  \\
            0  &0   &m+n
          \end{amat}
        \end{align*}
        We get that $0=m+n$, and that $t=m+n$, and that $s-t=-m$.
        The first gives $m=-n$.
        So the only vectors from \( S_2 \) that are also in
        \( S_1 \) are of this form.
        \begin{equation*}
          \colvec{-1 \\ 1 \\ 1}\cdot m
          +\colvec{0 \\ 1 \\ 3}\cdot (-m)
          =\colvec{-1 \\ 0 \\ -2}\cdot m
          =\colvec{1 \\ 0 \\ 2}\cdot (-m)
           +\colvec{-1 \\ 1 \\ 0}\cdot 0
        \end{equation*}
        For instance, taking $m\neq -n$ gives a vector from $S_2$ that is not an element
        of~$S_1$
        \begin{equation*}
          \colvec{-1 \\ 1 \\ 1}\cdot 1
            +\colvec{0 \\ 1 \\ 3}\cdot 0
          =\colvec{-1 \\ 1 \\ 1}
        \end{equation*}
        We could double check that with a linear system.
        \begin{equation*}
          \colvec{-1 \\ 1 \\ 1}
            =\colvec{1 \\ 0 \\ 2}s
            +\colvec{-1 \\ 1 \\ 0}t
          \qquad\qquad
          \begin{linsys}{2}
            s  &- &t  &=  &-1  \\
               &  &t  &=  &1   \\
           2s  &  &   &=  &1
          \end{linsys}
        \end{equation*}
        There is no need for Gauss's Method.
        The third equation gives $s=1/2$ and the second gives $t=1$.
        Those two are inconsistent with the first equation.
      \partsitem These sets are equal.

        First, to prevent confusion we change the letters used for the parameters.
        \begin{equation*}
          S_2=\set{\colvec{3 \\ 7 \\ 7}m
                   +\colvec{1 \\ 3 \\ 1}n
                   \suchthat m,n\in\Re}
        \end{equation*}

        To show that \( S_1\subseteq S_2 \) we will select an element of~$S_1$ and
        show that it is an element of~$S_2$.
        That is, fix $t$ and~$s$ and
        solve for $m$ and~$n$.
        \begin{equation*}
          \begin{linsys}{2}
           3m  &+  &n   &=  &t+2s\hfill  \\
           7m  &+  &3n  &=  &3t+4s\hfill  \\
           7m  &+  &n   &=  &t+6s\hfill  
          \end{linsys}
        \end{equation*}
        Gauss's Method
        \begin{equation*}
          % \begin{amat}{2}
          %   3  &1   &t+2s  \\
          %   7  &3   &3t+4s  \\
          %   7  &1   &t+6s
          % \end{amat}
          \grstep[(-7/3)\rho_1+\rho_3]{(-7/3)\rho_1+\rho_2}
          % \begin{amat}{2}
          %   3  &1    &t+2s  \\
          %   0  &2/3  &(2/3)t-(2/3)s  \\
          %   0  &-4/3 &(-4/3)t+(4/3)s
          % \end{amat}                                    \\
          \repeatedgrstep{2\rho_2+\rho_3}
          \begin{amat}{2}
            3  &1    &t+2s  \\
            0  &2/3  &(2/3)t-(2/3)s  \\
            0  &0    &0
          \end{amat}
        \end{equation*}
        has a unique solution.
        With $(2/3)n=(2/3)t-(2/3)s$ we get $n=t-s$.
        Substituting~$n$ into $3m+n=t+2s$ gives
        $3m+(t-s)=t+2s$ and so $m=s$.
        In brief, we have this.
        \begin{equation*}
          \colvec{1 \\ 3 \\ 1}t
          +\colvec{2 \\ 4 \\ 6}s
          =\colvec{3 \\ 7 \\ 7}\cdot s
          +\colvec{1 \\ 3 \\ 1}\cdot (t-s)
        \end{equation*}
        Therefore \( S_1\subseteq S_2 \).

        For the converse, \( S_2\subseteq S_1 \), fix an element of $S_2$
        and show it is expressible in the form for~$S_1$.
        That is, fix $m$ and~$n$
        and solve for $t$ and~$s$.
        \begin{equation*}
          \begin{linsys}{2}
           t   &+  &2s   &=  &3m+n\hfill  \\
           3t  &+  &4s   &=  &7m+3n\hfill  \\
           t   &+  &6s   &=  &7m+n\hfill  
          \end{linsys}
        \end{equation*}
        Gaussian reduction
        \begin{align*}
          \begin{amat}{2}
            1  &2   &3m+n  \\
            3  &4   &7m+3n  \\
            1  &6   &7m+n
          \end{amat}
          &\grstep[-\rho_1+\rho_3]{-3\rho_1+\rho_2}
          \begin{amat}{2}
            1  &2   &3m+n  \\
            0  &-2  &-2m    \\
            0  &4   &4m
          \end{amat}                                \\
          &\grstep{2\rho_2+\rho_3}
          \begin{amat}{2}
            1  &2   &3m+n  \\
            0  &-2  &-2m    \\
            0  &0   &0
          \end{amat}
        \end{align*}
        gives that we can take $s$ to be~$m$,
        and substitution into the first equation gives that we can take
        $t=m+n$.
        \begin{equation*}
          \colvec{3 \\ 7 \\ 7}m
          +\colvec{1 \\ 3 \\ 1}n
          =\colvec{1 \\ 3 \\ 1}\cdot (m+n)
          +\colvec{2 \\ 4 \\ 6}\cdot m
        \end{equation*}
        Hence any vector in~\( S_2 \) is also in~\( S_1 \).
        Therefore the two sets are equal.
    \end{exparts}  
   \end{answer}

\end{exercises}
